\newif\ifanon\anonfalse 

\documentclass[11pt]{article}
\usepackage[margin=1in]{geometry}
\usepackage[utf8]{inputenc}
\usepackage[T1]{fontenc} % Added to make curly braces look nice in \texttt mode

\usepackage{mathpazo}
\linespread{1.09} % Mathpazo looks nicer with line spread. 


% ============================================================================== 
% SECTION: Typography and Color Enhancements. 
% This section includes all packages and commands for enhancing text and color.
%   ==============================================================================
\usepackage{xcolor}
\usepackage{color}
\definecolor{niceRed}{RGB}{190,38,38}
\definecolor{niceYellow}{HTML}{f5b400}
\definecolor{blueGrotto}{HTML}{059DC0}
\definecolor{royalBlue}{HTML}{057DCD}
\definecolor{navyBlue}{HTML}{0B579C}
\definecolor{limeGreen}{HTML}{81B622}
\definecolor{nicePurple}{HTML}{9c27b0}
\definecolor{lightRoyalBlue}{HTML}{def2ff}  
\definecolor{gold}{HTML}{ffa300}

\newcommand{\white}[1]{\textcolor{white}{#1}}
\newcommand*{\red}[1]{\textcolor{red}{#1}}
\newcommand*{\blue}[1]{\textcolor{blue}{#1}}
\newcommand*{\gray}[1]{\textcolor{lightgray}{#1}}

% ============ Comment Management  ============
\usepackage{color-edits} 
\addauthor[Jane]{jl}{blue}
\addauthor[Anay]{am}{gold}
\addauthor[Manolis]{mz}{teal}

% ============ Reference and Bib Setup  ============
\usepackage{hyperref}
\hypersetup{
  colorlinks = true, 
  urlcolor = {blueGrotto},
  linkcolor = {royalBlue},
  citecolor = {navyBlue}
}
% 
\usepackage[
  backref=true,
  backend=biber,
  natbib=true,
  style=alphabetic,
  sorting=alphabeticlabel,
  sortcites=true,
  minbibnames=3,
  maxbibnames=999,
  mincitenames=1,
  maxcitenames=4,
  minalphanames=1,
  maxalphanames=6,
  doi=false, url=false, eprint=false, isbn=false, 
]{biblatex}

\renewcommand*{\bibfont}{\small}

\DeclareSortingTemplate{alphabeticlabel}{
  \sort[final]{%
    \field{labelalpha} % Sort by the alphabetic label (\eg{}, DKL, DLL)
  }
  \sort{%
    \field{year}   % If labels are the same, sort by year
  }
  \sort{%
    \field{title}  % If both label and year are the same, sort by title
  }
}

\AtBeginRefsection{\GenRefcontextData{sorting=ynt}}
\AtEveryCite{\localrefcontext[sorting=ynt]}
\addbibresource{refs.bib}


\newcommand{\plaincite}[1]{%
  \begingroup
    \renewcommand*{\bibopenbracket}{}%
    \renewcommand*{\bibclosebracket}{}%
    \cite{#1}%
  \endgroup
}
% ==============================================================================
% SECTION: Packages
% This section includes all other packages relevant
% ==============================================================================

% ============ Packages for tables  ============
\usepackage{booktabs} 
\usepackage{multicol} 
\usepackage{multirow} 
\usepackage{longtable}



% ============ Packages for figures  ============
\usepackage{subcaption}
\usepackage{graphicx}
\usepackage{float} 
\usepackage{tikz}
\usepackage{pgfplots}
\pgfplotsset{compat=1.17}
\usepgfplotslibrary{fillbetween}
\usetikzlibrary{arrows.meta}
\usetikzlibrary{calc}


\tikzset{
  myNodeFlex/.style={
    draw,
    rectangle,
    rounded corners,
    text centered,
    minimum height=1.5em,
  }
}

\tikzset{
  myNode/.style={
    draw,
    rectangle,
    rounded corners,
    text centered,
    minimum height=1.5em,
    minimum width=3cm,
    text width=5cm,    
  }
}

\tikzset{
  myNodeNarrow/.style={
    draw,
    rectangle,
    rounded corners,
    text centered,
    minimum height=1.5em,
    minimum width=1cm,
  }
}

\tikzset{
  myNodeWide/.style={
    draw,
    rectangle,
    rounded corners,
    text centered,
    minimum height=1.5em,
    minimum width=6cm,
  }
}


% ============ Packages for math  ============
\usepackage{physics} % useful math marcos
\usepackage{amsmath}
\usepackage{amssymb}
\usepackage{amsthm}
\usepackage{bbm}
\usepackage{blkarray}
\usepackage{mathtools}
\usepackage{nicefrac}
\usepackage{dsfont}
\usepackage{pifont} % for tick and cross symbols
\usepackage{thm-restate} % to re-state theorems before their proofs
\usepackage[capitalise,noabbrev,nameinlink,sort]{cleveref} % Added cleveref -- easy referencing

\usepackage{setspace}

% ============ Packages for lists and boxes ============
\usepackage{typed-checklist}
\usepackage{enumerate} 
\usepackage{enumitem} 
\usepackage{tcolorbox}
\usepackage[framemethod=TikZ]{mdframed}
\mdfsetup{%
backgroundcolor=gray!10, , %royalBlue!20, 
roundcorner=4pt,
skipabove=15pt,
skipbelow=5pt,
linewidth=1pt}
\usepackage{changepage} % to indent theorems, definitions etc which appear inside remarks

% ============ Packages for Algorithms ============
\usepackage{algorithm}
\usepackage{algpseudocode}[1]
\renewcommand{\algorithmicrequire}{\textbf{Input:}}
\renewcommand{\algorithmicensure}{\textbf{Oracles:}}


% ============ Packages for extended font selection ============
\usepackage{mathrsfs} % for mathscript
\usepackage{soul} % for underline that wraps
\setuldepth{Berlin}

 
% ================================================================================
% SECTION: Theorem Formatting and Numbering
% This section includes all commands for formatting and numbering thm-environments
% ================================================================================

% ============ Theorem Style Environments ============
%%%% ------------- Chained Numbering -------------
\theoremstyle{plain} 
\newtheorem{theorem}{Theorem}[section]
\newtheorem{corollary}[theorem]{Corollary}
\newtheorem{proposition}[theorem]{Proposition}
\newtheorem{lemma}[theorem]{Lemma}
\newtheorem{fact}[theorem]{Fact}
\newtheorem{claim}[theorem]{Claim}
\newtheorem{conjecture}[theorem]{Conjecture}

%%%% ------------- Environments with Separate Numbering Scheme -------------
\newtheorem{assumption}{Assumption}
\newtheorem{infassumption}{Informal Assumption}
\newtheorem{inftheorem}[theorem]{Informal Theorem}
\newtheorem{infcorollary}[theorem]{Informal Corollary}
\newtheorem{minftheorem}[theorem]{Main Informal Theorem}
\newtheorem{maintheorem}{Main Theorem}
\newtheorem{definition}{Definition}
\newtheorem{infdefinition}{Informal Definition}
\newtheorem*{definition*}{Definition}
\newtheorem{problem}{Problem}
\newtheorem{openproblem}{Open Problem}

%%%% ------------- Definition Style Environment -------------       
\theoremstyle{definition} 
\newtheorem{example}[theorem]{Example}
\newtheorem{remark}[theorem]{Remark}

%%%% ------------- Remark Style Environment -------------       
\theoremstyle{remark}
\newtheorem{observation}[theorem]{Observation}

% ============ Indenting Post Environments ============
\AfterEndEnvironment{definition}{\noindent\ignorespaces}
\AfterEndEnvironment{infdefinition}{\noindent\ignorespaces}
\AfterEndEnvironment{assumption}{\noindent\ignorespaces}
\AfterEndEnvironment{problem}{\noindent\ignorespaces}
\AfterEndEnvironment{openproblem}{\noindent\ignorespaces}
\AfterEndEnvironment{lemma}{\noindent\ignorespaces}
\AfterEndEnvironment{theorem}{\noindent\ignorespaces}
\AfterEndEnvironment{proposition}{\noindent\ignorespaces}
\AfterEndEnvironment{fact}{\noindent\ignorespaces}
\AfterEndEnvironment{question}{\noindent\ignorespaces}
\AfterEndEnvironment{corollary}{\noindent\ignorespaces}
\AfterEndEnvironment{model}{\noindent\ignorespaces}
\AfterEndEnvironment{remark}{\noindent\ignorespaces}
\AfterEndEnvironment{proof}{\noindent\ignorespaces}
\AfterEndEnvironment{fact}{\noindent\ignorespaces}
\AfterEndEnvironment{minftheorem}{\noindent\ignorespaces}
\AfterEndEnvironment{inftheorem}{\noindent\ignorespaces}
\AfterEndEnvironment{maintheorem}{\noindent\ignorespaces}
\AfterEndEnvironment{restatable}{\noindent\ignorespaces}
\AfterEndEnvironment{infassumption}
{\noindent\ignorespaces}
\AfterEndEnvironment{infcorollary}{\noindent\ignorespaces}

% ============ Cleverref Macros ============
\crefname{section}{Section}{Sections}
\crefname{theorem}{Theorem}{Theorems}
\crefname{lemma}{Lemma}{Lemmas}
\crefname{problem}{Problem}{Problems}
\crefname{program}{Program}{Progams}
\crefname{definition}{Definition}{Definitions}
\crefname{conjecture}{Conjecture}{Conjectures}
\crefname{corollary}{Corollary}{Corollaries}
\crefname{construction}{Construction}{Constructions}
\crefname{conjecture}{Conjecture}{Conjectures}
\crefname{claim}{Claim}{Claims}
\crefname{observation}{Observation}{Observations}
\crefname{proposition}{Proposition}{Propositions}
\crefname{fact}{Fact}{Facts}
\crefname{question}{Question}{Questions}
\crefname{problem}{Problem}{Problems}
\crefname{remark}{Remark}{Remarks}
\crefname{example}{Example}{Examples}
\crefname{equation}{Equation}{Equations}
\crefname{appendix}{Section}{Sections}
\crefname{algorithm}{Algorithm}{Algorithms}
\crefname{model}{Model}{Models}
\crefname{figure}{Figure}{Figures}
\crefname{infassumption}{Informal Assumption}{Informal Assumptions}
\crefname{inftheorem}{Informal Theorem}{Informal Theorems}
\crefname{infdefinition}{Informal Definition}{Informal Definitions}
\crefname{minftheorem}{Main Informal Theorem}{Main Informal Theorems}
\crefname{maintheorem}{Main Theorem}{Main Theorems}
\crefname{assumption}{Assumption}{Assumptions}
\crefname{step}{Step}{Steps}
\crefname{result}{Result}{Results}
\crefname{event}{Event}{Events}
\crefname{none}{}{}


\newlist{asmpenum}{enumerate}{1} % should only occur inside assumption env.
\setlist[asmpenum]{label={\arabic*.},ref=\theassumption.{\arabic*}}
\crefname{asmpenumi}{Assumption}{Assumptions}


% ================================================================================
% SECTION: Helpers (Alignment and tagging)
% This section includes all commands aligning and tagging (for proofs)
% ================================================================================

% ============ Indented Environment for Subproofs  ============
\usepackage{etoolbox}

\newenvironment{subenvironment}{%
    % \setlength{\parskip}{0pt}% Remove space between paragraphs
    % \vspace{-5mm}% Remove space before the content
    \begin{adjustwidth}{2em}{2em}% {#1}{#2}%
    % The % signs remove unwanted spaces
}{
    \end{adjustwidth}%
    % \vspace{-5mm}% Remove space after the content
}

\BeforeBeginEnvironment{subenvironment}{\white{.}\\ \vspace{-10mm}}
\AfterEndEnvironment{subenvironment}{\vspace{4mm}}

% ============ Helpers for Naming Equations ============
\newcommand{\yesnum}{\addtocounter{equation}{1}\tag{\theequation}} 
\makeatletter
\newcommand{\tagnum}[2]{%
    \refstepcounter{equation}%
    \tag{#1) \ (\theequation}%
    \protected@write \@auxout {}{%
        \string \newlabel {#2}{{\theequation}{\thepage}{}{equation.\theequation}{}}%
    }%
}
\makeatother

\makeatletter
% detach tags and references
\renewcommand{\eqref}[1]{\textup{\eqrefform@{\ref{#1}}}}
\let\eqrefform@\tagform@

\newcommand{\changetag}[1]{%
  \renewcommand\tagform@[1]{\maketag@@@{(\ignorespaces#1\unskip\@@italiccorr)}}%
}
\makeatother

% ============ Helpers for Adding Explanations in Proofs ============
\newcommand{\Stackrel}[2]{\stackrel{\mathmakebox[\widthof{\ensuremath{#2}}]{#1}}{#2}}

% ============ Helpers for Fromatting Text in Equations ============
\newcommand{\quadtext}[1]{\quad\text{#1}\quad}
\newcommand{\qquadtext}[1]{\qquad\text{#1}\qquad}
 
\newcommand{\quadand}{\quadtext{and}}
\newcommand{\qquadand}{\qquadtext{and}}

\newcommand{\quadwhere}{\quadtext{where}}
\newcommand{\qquadwhere}{\qquadtext{where}}

% ================================================================================
% SECTION: Math Macros
% This section includes all macros for math
% ================================================================================

% ============  Math Macros: Auto \left and \right ============
\newcommand{\mat}[1]{\begin{bmatrix}#1\end{bmatrix}}
\def\abs#1{\left| #1 \right|}
\def\sabs#1{| #1 |}
\newcommand{\given}{\;\middle|\;}
\newcommand{\sinparen}[1]{(#1)}
\newcommand{\binparen}[1]{\big(#1\big)}
\newcommand{\bbinparen}[1]{\bigg(#1\bigg)}
\newcommand{\sinbrace}[1]{\{#1\}}
\newcommand{\binbrace}[1]{\big\{#1\big\}}
\newcommand{\bbinbrace}[1]{\bigg\{#1\bigg\}}
\newcommand{\sinsquare}[1]{[#1]}
\newcommand{\inbrace}[1]{\left\{#1\right\}}
\newcommand{\sinangle}[1]{\langle#1\rangle}
\newcommand{\inparen}[1]{\left(#1\right)}
\newcommand{\insquare}[1]{\left[#1\right]}
\newcommand{\inangle}[1]{\left\langle#1\right\rangle}
\newcommand{\floor}[1]{\left\lfloor#1\right\rfloor}
\newcommand{\sfloor}[1]{\lfloor#1\rfloor}
\newcommand{\ceil}[1]{\left\lceil#1\right\rceil}
\newcommand{\sceil}[1]{\lceil#1\rceil}
\newcommand{\snorm}[1]{\ensuremath{\| #1 \|}}
\let\norm\relax
\newcommand{\norm}[1]{\ensuremath{\left\lVert #1 \right\rVert}}

% ============ Table formatting commands  ============
\newcommand{\midsepremove}{\aboverulesep = 0mm \belowrulesep = 0mm}
\newcommand{\midsepadd}{\aboverulesep = 1mm \belowrulesep = 1mm}
\newcommand{\midsepdefault}{\aboverulesep = 0.605mm \belowrulesep = 0.984mm}


% ============ Common Sets ============
\newcommand{\N}{\mathbb{N}}
\newcommand{\R}{\mathbb{R}}
\newcommand{\Q}{\mathbb{Q}}
\newcommand{\Z}{\mathbb{Z}}

% ============ Events ============
\newcommand{\evE}{\ensuremath{\mathscr{E}}}
\newcommand{\evF}{\ensuremath{\mathscr{F}}}
\newcommand{\evG}{\ensuremath{\mathscr{G}}}
\newcommand{\evH}{\ensuremath{\mathscr{H}}}

% ============ Math Operators  ============
\newcommand{\E}{\operatornamewithlimits{\mathbb{E}}} 
\newcommand{\Ex}{\E}
\newcommand{\Exp}{\Ex}
\newcommand{\cov}{\ensuremath{\operatornamewithlimits{\rm Cov}}}
% \newcommand{\var}{\ensuremath{\operatornamewithlimits{\rm Var}}}
% \newcommand{\tr}{\ensuremath{\operatornamewithlimits{\rm Tr}}}
\newcommand\ind{\mathds{1}}
\newcommand{\sgn}{\operatornamewithlimits{sgn}}
\newcommand{\argmin}{\operatornamewithlimits{arg\,min}}
\newcommand{\argmax}{\operatornamewithlimits{arg\,max}}

% ============ Functions, metrics ============
\newcommand{\tv}[2]{\operatorname{d}_{\mathsf{TV}}\sinparen{#1,#2}}
\newcommand{\kl}[2]{\operatornamewithlimits{\mathsf{KL}}\sinparen{#1\|#2}} 
\newcommand{\chidiv}[2]{\chi^2\inparen{#1\|#2}}
\newcommand{\srenyi}[3]{\cR_{#1}\sinparen{#2\|#3}}
\newcommand{\renyi}[3]{\cR_{#1}\inparen{#2\|#3}}

% ============ Math Shorthand ============
\newcommand{\zo}{\ensuremath{\inbrace{0, 1}}}
\newcommand{\oo}{\ensuremath{\inbrace{-1, 1}}}
\newcommand{\IntsTo}[1]{\setb{1, 2, \dotsc, #1}}
\newcommand{\IntsTon}{\IntsTo{n}}
\newcommand{\sfrac}[2]{{#1/#2}} 
\newcommand{\nfrac}[2]{\nicefrac{#1}{#2}}
\newcommand{\OPT}{\mathrm{OPT}} 
\newcommand{\opt}{\ensuremath{\mathrm{OPT}}}
\newcommand{\Err}{\mathrm{Err}}
\newcommand{\err}{\textrm{\rm Err}}
\newcommand{\poly}{\mathrm{poly}}
\newcommand{\polylog}{\mathrm{polylog}}
\newcommand{\vc}{\textrm{\rm VC}}
\newcommand{\vol}{\textrm{\rm vol}}
\newcommand{\supp}{\operatorname{supp}}
\newcommand{\cdf}{\ensuremath{{\rm cdf}}}
\newcommand{\iid}{i.i.d.}
\newcommand{\unif}{{\rm Unif}}
\newcommand{\nll}{\mathrm{Null}}

\newcommand{\cupdot}{\charfusion[\mathbin]{\cup}{\cdot}}
\newcommand{\bigcupdot}{\charfusion[\mathop]{\bigcup}{\cdot}}
\newcommand*\circled[1]{\tikz[baseline=(char.base)]{
            \node[shape=circle,draw,inner sep=1pt] (char) {#1};}}

% ============ Special Symbols  ============
\newcommand{\tick}{\ding{51}}
\let\cross\relax
\newcommand{\cross}{\ding{55}}

% ============ Nicer Math Symbols  ============
\newcommand{\eps}{\varepsilon}
\renewcommand{\epsilon}{\varepsilon}
% http://tex.stackexchange.com/a/217624/100722
\makeatletter
\newcommand*{\tran}{{\mathpalette\@tran{}}}
\newcommand*{\@tran}[2]{\raisebox{\depth}{$\m@th#1\intercal$}}
\makeatother

\mathchardef\NABLA"272
\newcommand*{\Nabla}{\boldsymbol\NABLA}
\let\nabla\Nabla

% ============ Hat-Symbols ============
\renewcommand{\hat}{\widehat}
\newcommand{\wh}[1]{\widehat{#1}}
\newcommand{\hf}{\wh{f}}

\newcommand{\hD}{\wh{D}}
\newcommand{\hF}{\wh{F}}
\newcommand{\hR}{\wh{R}}
\newcommand{\hS}{\wh{S}}
\newcommand{\hW}{\wh{W}}

% ============ Star-Symbols ============
\newcommand{\sS}{S^\star}

% ============ Overline-Symbols ============
\renewcommand{\bar}{\overline}
\newcommand{\oE}{\overline{E}}
\newcommand{\oS}{\overline{S}}

% ============ Tilde-Symbols ============
\renewcommand{\tilde}{\widetilde}
\newcommand{\wt}[1]{\widetilde{#1}}
\newcommand{\tS}{\wt{S}}
\newcommand{\tT}{\wt{T}}

% ============ Cal-Symbols ============
\newcommand{\customcal}[1]{\euscr{#1}}
\newcommand{\cA}{\customcal{A}}
\newcommand{\cB}{\customcal{B}}
\newcommand{\cC}{\customcal{C}}
\newcommand{\cD}{\customcal{D}}
\newcommand{\cE}{\customcal{E}}
\newcommand{\cF}{\customcal{F}} 
\newcommand{\cG}{\customcal{G}} 
\newcommand{\cH}{\customcal{H}} 
\newcommand{\cI}{\customcal{I}} 
\newcommand{\cJ}{\customcal{J}} 
\newcommand{\cK}{\customcal{K}}
\newcommand{\cL}{\customcal{L}}
\newcommand{\cM}{\customcal{M}}
\newcommand{\cN}{\customcal{N}}
\newcommand{\cO}{\customcal{O}}
\newcommand{\cP}{\customcal{P}} 
\newcommand{\cQ}{\customcal{Q}} 
\newcommand{\cR}{\customcal{R}} 
\newcommand{\cS}{\customcal{S}}
\newcommand{\cT}{\customcal{T}} 
% \newcommand{\cU}{\customcal{U}} 
\newcommand{\cV}{\customcal{V}} 
\newcommand{\cW}{\customcal{W}}
\newcommand{\cX}{\customcal{X}}
\newcommand{\cY}{\customcal{Y}}
\newcommand{\cZ}{\customcal{Z}}

% ================================================================================
% SECTION: Helpers (Style)
% This section includes all macros for Styling and Font Selection
% ================================================================================

% ============ Fancy fonts ============
\DeclareMathAlphabet{\mathdutchcal}{U}{dutchcal}{m}{n}
\SetMathAlphabet{\mathdutchcal}{bold}{U}{dutchcal}{b}{n}
\DeclareMathAlphabet{\mathdutchbcal}{U}{dutchcal}{b}{n}

% \usepackage{newtxtext}
% \usepackage{newtxmath}

\DeclareMathAlphabet\urwscr{U}{urwchancal}{b}{n}%
\DeclareMathAlphabet\rsfscr{U}{rsfso}{m}{n}
\DeclareMathAlphabet\euscr{U}{eus}{m}{n}
\DeclareFontEncoding{LS2}{}{}
\DeclareFontSubstitution{LS2}{stix}{m}{n}
\DeclareMathAlphabet\stixcal{LS2}{stixcal}{m} {n}

% ============ Miscellaneous Style ============
\newcommand{\toa}[1]{\textit{#1}} % term of art 
\renewcommand{\paragraph}[1]{\medskip \noindent\textbf{#1}~}
\newcommand{\itparagraph}[1]{\medskip \noindent\textit{#1}~}

\newcommand{\ie}{\textit{i.e.}}
\newcommand{\eg}{\textit{e.g.}} 


% ================================================================================
% SECTION: Project-Specific Macros
% This section includes all project-specific macros
% ================================================================================
\newcommand{\hypo}[1]{\mathdutchcal{#1}}
\newcommand{\hyA}{\hypo{A}}
\newcommand{\hyB}{\hypo{B}}
\newcommand{\hyC}{\hypo{C}}
\newcommand{\hyD}{\hypo{D}}
\newcommand{\hyE}{\hypo{E}}
\newcommand{\hyF}{\hypo{F}}
\newcommand{\hyG}{\hypo{G}}
\newcommand{\hyH}{\hypo{H}}
\newcommand{\hyI}{\hypo{I}}
\newcommand{\hyJ}{\hypo{J}}
\newcommand{\hyK}{\hypo{K}}
\newcommand{\hyL}{\hypo{L}}
\newcommand{\hyM}{\hypo{M}}
\newcommand{\hyN}{\hypo{N}}
\newcommand{\hyO}{\hypo{O}}
\newcommand{\hyP}{\hypo{P}}
\newcommand{\hyQ}{\hypo{Q}}
\newcommand{\hyR}{\hypo{R}}
\newcommand{\hyS}{\hypo{S}}
\newcommand{\hyT}{\hypo{T}}
\newcommand{\hyU}{\hypo{U}}
\newcommand{\hyV}{\hypo{V}}
\newcommand{\hyW}{\hypo{W}}
\newcommand{\hyX}{\hypo{X}}
\newcommand{\hyY}{\hypo{Y}}
\newcommand{\hyZ}{\hypo{Z}} 

\newcommand{\pERM}{Pessimistic-ERM}
\newcommand{\instancePERM}{\ensuremath{I_{\rm PERM}}}

\newcommand{\lreg}{\textsf{L1-Regression}}
\newcommand{\constrainedreg}{$\mathsf{Constrained\text{-}Regression}$} % Polynomial\text{-}

\renewcommand{\d}{{\rm d}}

% New macros 
\newcommand{\optset}{{H^\star}}
\newcommand{\opthypo}{S^\star}
\newcommand{\cDtrue}{\cD^\star}
\newcommand{\cDgiven}{\cD}
\newcommand{\cPtrue}{\cP^\star}

\newcommand{\unlabeledtrue}{{\cDtrue}}
\newcommand{\unlabeled}{{{\cDgiven}}}
\newcommand{\positivetrue}{{\cPtrue}}
\newcommand{\positive}{{\wh{\cP}}}

\renewcommand{\hat}{\widehat}
\renewcommand{\tilde}{\widetilde}

\newcommand{\pstar}{p^\star}
\newcommand{\hstar}{h^\star}
\newcommand{\muStar}{\mu^\star}
\newcommand{\SigmaStar}{\Sigma^\star}

\usepackage{array}
\newcolumntype{L}[1]{>{\raggedright\let\newline\\\arraybackslash\hspace{0pt}}m{#1}}
\newcolumntype{C}[1]{>{\centering\let\newline\\\arraybackslash\hspace{0pt}}m{#1}}
\newcolumntype{R}[1]{>{\raggedleft\let\newline\\\arraybackslash\hspace{0pt}}m{#1}}

\newcommand{\np}{\textsf{NP}}
\newcommand{\nL}{L}
 


\title{Learning with Positive and Imperfect Unlabeled Data}  

\ifanon
\author{Anonymous Author(s)}
\else 
\author{
    \centering
    \hspace{4mm}
    \begin{tabular}{C{4.75cm}C{4.75cm}C{4.75cm}}
    {\bf Jane H. Lee} & {\bf Anay Mehrotra} & {\bf Manolis Zampetakis}\\[2mm]
    Yale University & Yale University & Yale University\\[-4mm]
    {\small\phantom{....................}} \mbox{\small\href{mailto:jane.h.lee@yale.edu}{\texttt{jane.h.lee@yale.edu}}} & {\small \phantom{............}} \mbox{\small\href{mailto:anaymehrotra1@gmail.com}{\texttt{anaymehrotra1@gmail.com}}} & \mbox{\small\href{mailto:manolis.zampetakis@yale.edu}{\texttt{manolis.zampetakis@yale.edu}}}
    \\
    \end{tabular} 
}
\fi 

\date{}

\begin{document}

\maketitle
\thispagestyle{empty}

\begin{abstract}
    We study the problem of learning binary classifiers from positive and unlabeled data, in the case that the distribution of unlabeled data is shifted. We call this problem \textit{learning from positive and imperfect unlabeled data (PIU Learning)}. In the absence of covariate shifts, \ie{}, with \emph{perfect} unlabeled data, Denis (1998) reduced this problem to learning under Massart noise. This reduction, however, fails even under small covariate shifts.
    
    Our main results on PIU learning are: the characterization of the sample complexity of PIU learning and a computationally and sample-efficient algorithm achieving misclassification error $\eps$. 
    We illustrate the importance of PIU Learning by proving that our main result for PIU learning implies new algorithms for the following problems:
    \begin{enumerate}[leftmargin=5mm]
        \item \textbf{Learning from smooth distributions}, where we give algorithms that learn interesting concept classes from \textit{only} positive samples under \textit{smooth} feature distributions, which circumvents existing impossibility results and contributes to recent advances in learning algorithms under smoothness {(Haghtalab, Roughgarden, and Shetty, J.~ACM'24) (Chandrasekaran, Klivans, Kontonis, Meka, and Stavropoulos, COLT'24)}.


        \item \textbf{Learning with a list of unlabeled distributions,} where we design new algorithms that apply to a broad class of concept classes under the assumption that we are given a list of unlabeled distributions, one of which -- unknown to the learner -- is $O(1)$-close to the true distribution of features.
        
        
        \item \textbf{Estimation in the presence of unknown truncation}, where we give the first polynomial sample and time algorithm for estimating the parameters of an exponential family distribution from samples truncated to an unknown set $S^\star$ that is approximable by polynomials in $\nL_1$-norm.
            This improves the algorithm by Lee, Mehrotra, and Zampetakis (FOCS'24) 
            that requires approximation in $\nL_2$-norm -- a significantly stronger condition.
        \item \textbf{Detecting truncation}, where we present the first algorithm for detecting whether given samples have been truncated (or not) for a broad class of distributions, including non-product distributions. This improves upon the algorithm by De, Li, Nadimpalli, and Servedio (STOC'24), which was limited to product distributions.
    \end{enumerate}

\end{abstract}
 
\clearpage

\thispagestyle{empty}
{   
    \tableofcontents
}
\thispagestyle{empty}

\clearpage
\pagenumbering{arabic}
    
\section{Introduction} \label{sec:intro}
    \vspace{-2mm}
    
     In the classical PAC Learning framework \cite{valiant1984theory}, the goal is to learn a binary concept from labeled samples by minimizing misclassification error. Agnostic PAC model makes no assumptions on the data distribution and allows arbitrary noise, but is often computationally hard \cite{guruswami2009hardness, feldman2006new, daniely2016complexity}. The realizable setting, by contrast, assumes perfectly labeled data and admits more efficient algorithms. Beyond these models there exists a wide range of practically and theoretically motivated models that relax the realizable assumption while retaining tractability, motivated by label noise, distribution shift, and missing labels \cite{massartNoise2006, yishay_2009_domain_adapt, hanneke_kpotufe_2019_transfer_exp}. %

     A popular variant of PAC Learning, that arises in many real-world settings, is when only positive examples are readily available, while negative labels are scarce or costly to obtain. This arises, for example, in medical diagnosis, where confirmed cases of a condition are recorded but healthy controls are unlabeled \cite{elkan2008pu}, or in web search and recommendation systems, where user engagement indicates positive feedback but non-engagement is ambiguous \cite{liu2003building}. The problem of \textit{learning from positive samples} was first theoretically explored by \citet{natarajan1987learning}. While they did not characterize the learnability of binary concept classes from positive data, they proved information-theoretic impossibility results, showing that even simple hypothesis classes (such as halfspaces) cannot be learned from positive samples alone. 
 
     
     The practical challenge of solving classification without access to negative samples led to the incorporation of \textit{unlabeled data} to increase the amount of usable information. 
     This gives rise to a more tractable learning problem -- \textit{learning from positive and unlabeled samples (PU learning)}. \citet{denis1990positive} introduced this model and showed that PU learning can be reduced to agnostic PAC learning, thereby establishing the classical VC-\mbox{dimension-based upper bound on its sample complexity.}\footnote{\citet{denis1990positive} reduced PU learning to learning under Massart noise and due to \citet{massartNoise2006}, this yields a sample complexity of $\wt{O}(\nfrac{1}{\eps})$ for PU learning, compared to the $\wt{O}(\nfrac{1}{\eps^2})$ bound for agnostic learning \cite{kearns1992toward}.}
 
     In this paper, we revisit the problem of PU learning by observing that a natural generalization of PU learning leads to a broader learning framework that underlies several fundamental problems in statistics and learning theory -- problems for which the original PU model is too restrictive. Notably, \citet{denis1990positive} assumed that the distribution of positive data is identical to the distribution of unlabeled data, except that it is truncated to include only positive samples. Our key insight is that allowing for some discrepancy between these distributions leads to a more flexible model, which we call \textit{learning from positive and imperfect unlabeled data (PIU learning)}. This generalization captures a variety of important theoretical settings:
     \vspace{-1.5mm}
     \begin{enumerate}[leftmargin=15pt,itemsep=-3pt]
     \item[$\blacktriangleright$] \textbf{Learning with Smooth Positive Samples.} Inspired by recent work on smoothness-based learning \cite{haghtalab2020smooth, haghtalab2021smooth, pmlr-v247-chandrasekaran24a, blanchard2024agnostic}, we show that bypassing the negative result of \citet{natarajan1987learning} for learning from positive-only data under smoothness reduces to PIU learning.
     
     \item[$\blacktriangleright$] \textbf{Truncated Statistics.} PIU learning arises naturally in the estimation of high-dimensional distributions under unknown truncation or censoring \cite{daskalakis2018efficient, Kontonis2019EfficientTS}.


     \item[$\blacktriangleright$] \textbf{Detecting Truncation.} PIU learning also provides a framework for detecting whether a dataset has been subject to censoring -- a problem explored in recent work \cite{de2023detectingConvex, de2024detecting}.

     \end{enumerate}
     \vspace{-2mm}

    \noindent Motivated by these and further applications discussed later, we seek to answer the following.

    \vspace{1mm}
    
     \begin{mdframed}[skipbelow=0pt,skipabove=4pt]
       \textbf{Question:} \textit{Can we learn with positive and unlabeled data when unlabeled data is imperfect?}
     \end{mdframed}
     



    \paragraph{Roadmap.} 
        In the next section (\cref{sec:intro:PIU}), we give a more precise formulation of PIU learning and then provide an informal statement of our main result in \cref{sec:informal}.
        Followed, in \cref{sec:intro:application}, by informal statements of applications that we highlighted above.
        Next, we overview the key ideas in our proof in \cref{sec:technicalOverview}. 
        Finally, we discuss some real-world applications of PIU Learning, open problems, and additional related work in \cref{sec:intro:realWorldApplications,sec:intro:takeawaysOpenProblems,sec:intro:relatedWorks} respectively. 
        
     
   
\subsection{Learning with Positive and Imperfect Unlabeled Samples} \label{sec:intro:PIU}

In the PU learning model, one is given a concept class $\hyH$ and access to two sample-streams:
\begin{enumerate}[leftmargin=19pt,itemsep=0pt]
  \item[(a)] \textit{Unlabeled samples} drawn from a distribution $\cDtrue$ over the feature space, and 
  \item[(b)] \textit{Positive samples} drawn from a distribution $\cPtrue$, which corresponds to the conditional distribution of $\cDtrue$ restricted to the set $\optset \coloneqq \{x : h^{\star}(x) = 1\}$ for some target concept $h^{\star} \in \hyH$. 
\end{enumerate}  
That is, $\cPtrue$ is the truncation of $\cDtrue$ to the region where the label is positive. This precise structural relationship between $\cPtrue$ and $\cDtrue$ -- positives being drawn from the true distribution conditioned on label 1 -- is fundamental to existing PU learning guarantees \cite{denis1990positive}. When this assumption fails, the standard techniques break down, and the learnability of the problem is not well understood.

\paragraph{PIU Learning Model.} In this work, we initiate the study of \textit{Positive and Imperfect Unlabeled (PIU) Learning}. As in the standard PU learning setting, positive samples are drawn from a distribution $\cPtrue$ corresponding to the conditional distribution of a target function $h^\star \in \hyH$. However, instead of accessing unlabeled samples from the true distribution $\cDtrue$, we are given samples from a different distribution $\cDgiven$, which may not align exactly with $\cDtrue$. We assume that $\cDtrue$ is \textit{smooth} with respect to $\cDgiven$ -- a relaxed relationship that allows for bounded discrepancies between the two. This generalization captures settings where the unlabeled data may be corrupted, biased, or subject to distribution shift, making the PIU model strictly more general than classical PU learning. The precise relationship between $\cDtrue$ and $\cDgiven$ is given by the following assumption.
\begin{restatable}[Generalized Smoothness]{assumption}{smoothness}\label{asmp:smoothness}
    There are known constants $\sigma \in (0,1]$ and $q\geq 1$, such that, for any measurable set $S$, $\cDtrue(S)\leq {(\nfrac{1}{\sigma})}\cdot \cDgiven(S)^{\sfrac{1}{q}}$.
\end{restatable}
For $q=1$, this assumption reduces to the notion of smoothness considered by the line of recent works on smoothened analysis and is natural to expect for real-world  distributions \cite{manthey2020smoothed,haghtalab2022efficientSmooth,haghtalab2020smooth,haghtalab2024smoothJACM,block2024performanceempiricalriskminimization,blanchard2024agnostic}. When $q=2$, generalized smoothness is implied by $\chidiv{\cDtrue}{\cD} = O(1)$ and, for larger $q$, it is implied by bounds on the suitable Rényi divergence (see \cref{sec:distCloseness}).\footnote{The $\chi^2$-divergence of $\cP$ with respect to $\cQ$ is $\chidiv{\cP}{\cQ}=\int_x \inparen{\nfrac{\d \cP(x)}{\d \cQ(x)}}^2 \d \cQ(x) - 1$. We define Rényi divergences in \cref{sec:preliminaries}, although understanding their definition is not crucial for our results.} Furthermore, this assumption has also appeared in recent works on transfer learning, where $q$ relates to the \textit{transfer exponent} defined in \cite{hanneke_kpotufe_2019_transfer_exp,pmlr-v195-hanneke23c,hanneke2024unifiedtheorytransferlearning}.
\smallskip






As we will present in the next section, generalized smoothness suffices to design sample efficient algorithms. For our computationally efficient algorithms though, we need to additionally assume that the positive sample set $\optset$ has non-trivial mass.
\begin{restatable}[Non-Trivial Fraction of Positives]{assumption}{posFraction}\label{asmp:posFraction}
    For some known $\alpha>0$, $\Pr_{x\sim \cDtrue}\inparen{h^\star(x)=1}\geq \alpha$.
\end{restatable} 

\noindent This is required by prior works on PU learning that design computationally efficient algorithms (\eg{}, \cite{lee2024unknown}), and in virtually all the existing work on truncated statistics \cite{daskalakis2018efficient, Kontonis2019EfficientTS, lee2023learning, lee2024unknown} which is one of the main application of our model.%


\subsection{Our Main Results for PIU Learning}\label{sec:informal}
In this section, we present our main results for learning from positive and imperfect unlabeled examples. Then, in \cref{sec:intro:application}, we present the applications of these results to other learning theory problems, as we highlighted before. %
\subsubsection{Sample Efficient PIU Learning}
Our first result provides the sample complexity of PIU learning.
\begin{restatable}[Sample Complexity]{theorem}{sampleComplexity}
    \label{infthm:sampleComplexity}\label{thm:sampleComplexity}
    Suppose \cref{asmp:smoothness} holds. %
    Fix any $\eps,\delta\in (0,\nfrac{1}{2})$.
    There is an algorithm that, given $\eps,\delta$ and $n=\wt{O}\binparen{{{(\eps \sigma)}}^{-2q}\cdot \inparen{\vc{}(\hyH)+\log{\nfrac{1}{\delta}}}}$ independent samples from $\cPtrue$ and $\cDgiven$, outputs a hypothesis $\wh{h}$ such that, with probability $1-\delta$,
    \[ 
        \Pr_{\cDtrue}\binparen{\wh{h}(x)\neq h^\star(x)}\leq \eps\,.
    \]
\end{restatable}
An implication of this theorem is that PIU learning is also characterized by the VC-dimension. Interestingly, however, the algorithm in \cref{thm:sampleComplexity} is not empirical risk minimization (ERM). In fact, ERM provably fails in the PIU learning setting. Instead, \cref{thm:sampleComplexity} proposes an iterative procedure that solves a constrained optimization problem, which we refer to as \textit{\pERM{}}. This yields an improper learning algorithm, and we further show that proper learning is in fact impossible for PIU learning (see \cref{rem:necessityOfImproper}). Both the failure of ERM and the impossibility of proper learning underscore the fundamental gap between classical PU learning and the generalized framework introduced in this paper. 
\smallskip

\noindent \textbf{Improved Sample Complexity for $q = 1$.} Under an additional assumption that complements generalized smoothness (\cref{asmp:smoothness}) by providing a lower bound on $\cDtrue(S)$ for measurable sets, we further improve the sample complexity in the case $q = 1$ to $\wt{O}\sinparen{\nfrac{\vc{}(\hyH)}{\eps}}$ (\cref{thm:sampleComplexity:bothSided}). This matches the optimal rate for realizable learning and is therefore optimal for PIU learning with $q = 1$, despite the latter being a harder problem. 

\subsubsection{Computationally Efficient PIU Learning} \label{sec:intro:computationallyEfficientResult}

We next investigate the computational complexity of PIU learning under the additional assumption of a non-trivial mass of positive samples (\cref{asmp:posFraction}).

\begin{inftheorem}[Computational Complexity; see \cref{thm:main}]\label{mainthm:algo}\label{infthm:main}
    Suppose \cref{asmp:smoothness,asmp:posFraction} hold. 
    Fix $\eps,\delta\in (0,\nfrac{1}{2})$, and let $k$ be such that degree-$k$ polynomials $\zeta$-approximate $\hyH$ in \mbox{$\nL_1$-norm} with respect to $\cDgiven$ for $\zeta\leq \inparen{{\eps\alpha}{\sigma}}^{\poly(q)}$.
    Then, there is an algorithm that, given $n=\poly(d^k,\nfrac{1}{\zeta},\log{\nfrac{1}{\delta}})$ independent samples from $\cPtrue$ and $\cDgiven$, outputs a hypothesis $\wh{h}$ such such that, with probability $1-\delta$, 
    \[
        \Pr_{\cDtrue}\binparen{\wh{h}(x)\neq h^\star(x)}\leq \eps\,.
    \]
    The algorithm runs in time $\poly(n)$.
\end{inftheorem}
        Therefore, efficient PIU learning is possible when the optimal hypothesis $\optset$ has a low complexity in the sense that low-degree polynomials can well approximate it under distributions $\cDgiven$ (\cref{def:poly_approx}). 
        %
        The degree of polynomials $k$ required typically scales as $k = \poly\inparen{\nfrac{1}{\eps}}$. 
        This exponential dependence on $\nfrac{1}{\eps}$ is necessary for any SQ algorithm (such as the one in \cref{mainthm:algo}) due to a recent lower bound by \citet*{diakonikolas2024statistical} (see \cref{rem:SQlowerbound} for details).

        Notably, \cref{mainthm:algo} only requires an approximation with respect to the accessible unlabeled distribution $\cDgiven$ and not the inaccessible distribution $\cDtrue$.
        This enables the use of recent testable learning techniques when certain sandwiching polynomials exist \cite{rubinfeld2023testing,gollakota2023momentMatching,gollakota2023universal} to ensure that the required polynomial approximation guarantee holds by, roughly, verifying that the moments of samples from $\cDgiven$ match those of a known reference distribution (see \cref{rem:testing} for details). 
        Finally, we note that the requirement of $\nL_1$-approximability is natural and is required by nearly all existing computationally efficient agnostic learning algorithms.
        \cref{tab:computational}
        compares our method with prior PU Learning works.
        
        

        
        \begin{table}[htb!]
            \centering
            \vspace{4mm}
            \midsepremove{}
            \begin{tabular}{C{1.75cm}|C{1.9cm}|C{2.6cm}|C{2.65cm}|C{2cm}|C{3cm}|}
                & Runtime & Works without Density Lower Bound? & Polynomial Approximation & Constant $q$ & Additional Assumptions \\
                \cmidrule(r){2-6} 
             \cite{denis1990positive} & $d^{\,L_1\inparen{\poly(\alpha\eps)}}$ & \cross{} & $L_1$-norm & $q=1$ & --- \\
             \cite{lee2024unknown} & $d^{\,{L_2}\sinparen{ \eps^{\poly(1/\alpha)}}}$ & \cross{} & $L_2$-norm & $q\in [1,2]$ & Gaussian $\cD,\cDtrue$\textsuperscript{($\dagger$)} \\
             This Work & $d^{\,L_1\inparen{\poly(\alpha\eps)}}$ & \tick{} & $L_1$-norm & $q=O(1)$ & ---\\
            \end{tabular} 
            \midsepdefault{}
            \caption{\small \textit{Computational Complexity of Learning in the PIU Learning Model:}\quad 
                The table presents run times of existing algorithms under the PIU learning model (\cref{prob:PUlearning}). %
                All works require \cref{asmp:smoothness,asmp:posFraction}.\textsuperscript{~($\star$)}
                For any $\zeta>0$, let $L_1(\zeta)$ and $L_2(\zeta)$ denote the minimum degrees required to $\zeta$-approximate the underlying hypothesis class $\hyH$ in $\nL_1$-norm and $\nL_2$-norm with respect to $\cD$ respectively.
                Since $\nL_1$-approximation is less stringent, $L_1(\zeta)\leq L_2(\zeta)$, and for some hypothesis classes, $L_1(\zeta)$ is much smaller.
                Our method improves upon \cite{lee2024unknown} in several respects: (1) we can handle $q=O(1)$, do not require a density lower bound (\cref{asmp:densityLB}), and extend beyond Gaussians, (2) our running time scales with $L_1(\cdot)$ instead of $L_2(\cdot)$, and (3) we avoid doubly-exponential dependence on $\alpha$.
            \newline
            \vrule height 0.2pt width 15cm\newline
            {\footnotesize \white{.}\quad  ($\star$)~~To obtain a computationally efficient learner, we combine \lreg{} with the reduction from \cite{denis1990positive} and}\\[-1mm]
            \mbox{\footnotesize \white{.}\quad ~~~~~~ require \cref{asmp:posFraction} to guarantee polynomial approximability with respect to the distribution in the reduction.}
            \\
            \mbox{\footnotesize \white{.}\quad  ($\dagger$)~~\cite{lee2024unknown} need $\cD$ and $\cDtrue$ to have a $\chi^2$-bridge (see Condition \eqref{unknowntrunc:chisquarebridge}) which is only known to exist when both are}
            \\[-1mm]
            \mbox{\footnotesize \white{.}\quad ~~~~~~ Gaussians, or when both are product exponentials.}
            }
            \vspace{-4mm}
            \label{tab:computational}
        \end{table}
          

\subsection{Applications to Other Learning Theory Settings} \label{sec:intro:application}

In this section, we present applications of our main results across various learning theory settings. In many of these cases, we give the first algorithms for the respective problems or obtain significant improvements over the best-known results.  

    
\subsubsection{Learning with Smooth Positive Samples}\label{sec:intro:application:smooth}

\citet{natarajan1987learning} characterized the classes that can be \textit{properly} PAC learned from only positive samples showing that even simple classes such as two-dimensional halfspaces are not properly learnable from positive samples alone.

\begin{mdframed}
  \textbf{Question:} \textit{Can we improperly learn interesting hypothesis classes using only positive examples?}
\end{mdframed}

\noindent Our first result in this section shows that the answer to this question is negative. In fact, most interesting hypothesis classes (including two-dimensional halfspaces) are not PAC learnable from positive examples alone, either properly or improperly.

\begin{restatable}[Impossiblity of Learning from Positive Samples]{theorem}{impossiblityPositiveSampleLearning}
    \label{thm:impossibilityPositiveSamples}
    Let $\hyH$ be any hypothesis class satisfying $\lim_{k\to\infty} \vc{}(\hyH_{\cap k})=\infty$ where $\hyH_{\cap k}$ is the hypothesis class formed by taking intersections of $k$ hypothesis from $\hyH$, \ie{}, 
    $\hyH_{\cap k}\coloneqq \inbrace{H_1\cap H_2\cap \dots\cap H_k\colon H_1,H_2,\dots,H_k\in \hyH}.$
    
    For any (randomized) algorithm, there is an $h^\star\in \hyH$ and distributions $\cD,\cDtrue$ satisfying \cref{asmp:smoothness} and \cref{asmp:posFraction} with parameters $(\sigma,q,\alpha)=(\nfrac{1}{2},1,\nfrac{1}{2})$, such that given sample access to $\cD$ and $\cPtrue$ (the truncation of $\cDtrue$ to the set of positive samples), after any finite time, with probability $\nfrac{1}{2}$, 
    the algorithm outputs a hypothesis $\wh{h}$ such that, with probability $1-\delta$,
    \[ 
        \Pr_{\cDtrue}\binparen{\wh{h}(x)\neq h^\star(x)}\geq \Omega(1)\,.
    \]
\end{restatable}
This is a strong impossibility result: it applies to most common classes, although two exceptions are the family of axis-aligned rectangles and monotone conjunctions, which can be \textit{properly} learned from just positive samples. Following the paradigm of \cite{haghtalab2020smooth, haghtalab2021smooth} we ask whether this impossibility result can be bypass if we make smoothness assumptions for the input distribution.

\begin{mdframed}
  \textbf{Question:} \textit{Can we learn interesting classes using only positive samples under smoothness assumptions for the input distribution?}
\end{mdframed}

\noindent An immediate corollary of \cref{thm:sampleComplexity} and \cref{infthm:main} is  an affirmative answer to this. %
\begin{infcorollary}[see \cref{cor:smooth_classification}]\label{infcor:smooth_classification}\label{infcor:smooth}
    Consider a PAC learning instance with class $\hyH$, feature distribution $\cDtrue$ that is $(\sigma,q)$-smooth (for $\sigma=\Omega(1)$; see \cref{asmp:smoothness}) with respect to a known reference $\cD$ that one can efficiently sample from, and target $h^\star\in \hyH$.  
    \begin{itemize}
        \item[$\triangleright$] \textit{(Sample Complexity)} ~~
            $\wt{O}\sinparen{{\eps^{-2q}}\cdot \inparen{\vc{}(\hyH) + \log{\nfrac{1}{\delta}}}}$ independent positive samples are sufficient to $(\eps,\delta)$-PAC learn $\hyH$ with respect to $\cDtrue$.
        \item[$\triangleright$] \textit{(Computational Complexity)}~~
            Suppose $\cDtrue(\optset)\geq \Omega(1)$.
            Let $k\geq 1$ be such that $\hyH$ is $\eps^{O(q)}$-approximable by degree-$k$ polynomials in the $\nL_1$-norm with respect to $\cD$.
            Then, $n=\poly(d^k,\nfrac{1}{\eps},\log{\nfrac{1}{\delta}})$ positive samples and $\poly(n)$ time is sufficient to $(\eps,\delta)$-PAC learn $\hyH$ with respect to $\cDtrue$.
    \end{itemize}
\end{infcorollary}
Hence, learning from positive samples is also characterized by the VC dimension provided the distribution of covariates is smooth with respect to a known reference distribution $\cDgiven$. Moreover, the efficient algorithm works whenever $\hyH$ can be approximated in $\nL_1$-norm by low-degree polynomials with respect to the reference distribution $\cDgiven$; which is the standard requirement for efficient agnostic learning algorithms (\eg{}, \cite{kalai2008agnostically,pmlr-v247-chandrasekaran24a}).

    

    These results add to the growing line of work on smoothened (online) learning \cite{haghtalab2022efficientSmooth,haghtalab2020smooth,haghtalab2024smoothJACM,block2024performanceempiricalriskminimization} that shows that online learning becomes significantly easier when the adversary is restricted to pick responses from smooth distributions (their notions of smoothness corresponds to the strongest notion of smoothness in \cref{asmp:smoothness}; namely, $q=1$).
    


\subsubsection{Learning with Positive Samples and Lists of Unlabeled Samples}\label{sec:intro:application:listDecoding}

Following the application in the previous section, one may ask if learning is possible from just positive samples without access to an underlying ``reference'' distribution $\cD$. In this section, we show that learning is possible with some minimal information. 
    
In particular, we show that, in addition to the positive samples, it is sufficient to have a list of distributions that contains one \textit{unknown} distribution that ``witnesses'' the smoothness of $\cDtrue$.
\begin{definition}[List Decoding PIU Model]
    \label{def:list}
    Fix a distribution $\cDtrue$ over domain $\cX$ and numbers $\ell,q\geq 1$ and $\sigma\in (0,1]$.
    In the List Decoding model specified by parameters $(\cDtrue,\ell,\sigma,q)$, one is provided a list of $\ell$ distributions $\cDgiven_1, \cDgiven_2, \dots, \cDgiven_\ell$ over $\cX$ such that for some \textit{unknown} index $1\leq i^\star\leq \ell$,
    $\cD_{i^\star}$ is $(\sigma,q)$-smooth with respect to $\cD_{i^\star}$ in the sense of \cref{asmp:smoothness}. 
    The distributions $\unlabeled_j$ (for $j\neq i^\star$) can be arbitrary.
\end{definition}
This model generalizes PIU learning, which assumes sample access to $\cD_{i^\star}$. It is inspired by real-world scenarios where the unlabeled samples are scraped from unreliable sources that may be adversarially corrupted, which, in some cases, can be used to obtain a list of unlabeled samples satisfying \cref{def:list} (see \cref{rem:list}).

Our main result for this list-decoding model shows that it is also characterized by VC dimension and its computational complexity can also be bounded by polynomial approximability.
\begin{inftheorem}[see \cref{thm:listDecoding}]\label{infthm:list}
    Consider a PAC learning instance with      
        class $\hyH$,
        feature distribution $\cDtrue$, and 
        target $h^\star\in \hyH$.  
    Consider the corresponding instance of the List Decoding PIU Model in \cref{def:list} with $\sigma=\Omega(1)$ and $\ell=O(1)$ and parameter $q\geq 1$.
    \begin{itemize}
        \item[$\triangleright$] \textit{(Sample Complexity)}~~ 
            $\wt{O}\inparen{\eps^{-2q}\cdot \inparen{\vc{}(\hyH)+\log{\nfrac{1}{\delta}}}}$ independent positive samples are sufficient to $(\eps,\delta)$-PAC learn $\hyH$ with respect to $\cDtrue$.
        \item[$\triangleright$] \textit{(Computational Complexity)}~~
            Suppose $\cDtrue(\optset)\geq \Omega(1)$.
            Let $k\geq 1$ be such that $\hyH$ is $\zeta$-approximable by degree-$k$ polynomials in the $\nL_1$-norm with respect to $\cD_{i^\star}$ for $\zeta\leq \eps^{\,O(q)}$.
            Then, $n=\poly(d^k,\nfrac{1}{\zeta},\log{\nfrac{1}{\delta}})$ positive samples and $\poly(n)$ time is sufficient to $(\eps,\delta)$-PAC learn $\hyH$ with respect to $\cDtrue$.
    \end{itemize}
\end{inftheorem}
This theorem significantly extends our main result (\cref{infthm:main}), which is equivalent to the special case of $\ell=1$. Despite this, interestingly, minimal changes to the algorithm in our main result (\cref{infthm:main}) are sufficient to give the above result. The proof of \cref{infthm:list} appears in \cref{sec:intro:application:listDecoding} along with a formal statement.

\begin{remark}[Obtaining a List of Distributions From Adversarially Corrupted Data]\label{rem:list}
  In this remark, we show that if one is given adversarially corrupted unlabeled samples, then one can satisfy the guarantees of the above list-decoding model. Suppose that $\cDtrue$ is a Gaussian distribution $\cN(\muStar,\SigmaStar)$. Consider a data source that is $(1-\gamma)$-corrupted for some $\gamma\in (0,1]$: with probability $\gamma$, the source outputs a sample from $\cDtrue$, and, otherwise with probability $1-\gamma$, it outputs an arbitrary sample that may depend on the past samples and the knowledge of the learning algorithm. In \cref{sec:advCorruption}, we show that one can efficiently convert an $(1-\gamma)$-corrupted source for any $\gamma>0$, into one satisfying the requirements of the List Decoding PIU Model with $\ell= \poly(\nfrac{1}{\gamma})$. Moreover, under mild assumptions, we also extend this to the case where $\cDtrue$ is an exponential family distribution.
\end{remark} 



    \subsubsection{Truncated Estimation with Unknown Survival Set}\label{sec:intro:application:truncation}
        Our techniques have implications for efficient statistical inference from truncated samples, a problem with a long history of study in Statistics and Econometrics \cite{Galton1897,Pearson1902,PearsonLee1908,Lee1915,fisher31,maddala1983limited,Cohen91,hannon1999estimation,raschke2012inference} that has attracted considerable recent interest in Theoretical Computer Science \cite{daskalakis2018efficient, Kontonis2019EfficientTS, daskalakis2019computationally, trunc_regression_unknown_var, ons_switch_grad, lee2023learning, diakonikolas2024statistical,galanis2024discreteTruncated,lee2024unknown,nagarajan2020truncatedMixtureEM,fotakis2020booleanProductTruncated,nagarajan2023EMMixtureTruncation,tai2023mixtureCensored}.

        One important goal in truncated statistics is to estimate the parameters of a distribution when samples are observed only if they fall in some \textit{unknown} survival set $S^\star\subseteq\R^d$.
        A bit more formally, 
        there is a parametric family $\hyD=\inbrace{\cD(\theta)\colon \theta}$ with target parameter $\theta^\star$. 
        The goal is to find an estimate $\wh{\theta}$ of $\theta^\star$ given samples from the truncation\footnote{The truncation of a distribution $\cD$ to a set $S$ is the conditional distribution of $\cD$, conditioned on the event that a point lies in $S$. When $\cD$ has a density, the truncated density is given by $\frac{1}{\cD(S)} \cdot \ind\inparen{x \in S} \cdot \cD(x)$} of the target distribution $\cDtrue = \cD(\theta^\star)$ to the set $S^\star$ -- without \textit{any} access to $S^\star$.
        For this problem to be tractable, it is necessary to assume a lower bound on the mass of $S^\star$, specifically that $\cDtrue(S^\star)\geq \alpha$.
        The state-of-the-art algorithm for this problem \cite{lee2024unknown} works in three phases:
        \begin{enumerate}[itemsep=0pt]
            \item \textit{Phase 1:}~ First, it finds an unlabeled density $\cD$ close to $\cDtrue$ (in the sense of \cref{asmp:smoothness});
            \item \textit{Phase 2:}~ Then, it learns $S$ such that the symmetric difference $\cDtrue(S\triangle S^\star)\approx 0$; and 
            \item \textit{Phase 3:}~ Finally, it recovers $\theta^\star$ using  learned set $S$. 
        \end{enumerate}
        The bottleneck in their algorithm is Phase 2 which falls under the setting of \cref{infthm:main} and, substituting the algorithm in Phase 2 by the algorithm in \cref{infthm:main} gives us the following improvement upon \citet{lee2024unknown}'s result.
        
         
        \begin{infcorollary}[Truncated Estimation with Unknown Survival Set; see \cref{cor:truncation}]\label{infcor:truncation} 
            Let the family $\hyD=\inbrace{\cD(\theta)\colon \theta}$ satisfy the assumptions of \cite{lee2024unknown} and let $\alpha=\Omega(1)$. %
            Fix any $\eps, \delta \in (0, \nfrac{1}{2})$, and let $S^\star$ be $\poly(\eps)$-approximable by degree-$k$ polynomials in $\nL_1$-norm with respect to $\hyD$.
            There is an algorithm that given $n = \poly(d^k, \nfrac{1}{\eps}, \log{\nfrac{1}{\delta}})$ independent samples from $\cD(\theta^\star)$ truncated to $S^\star$ 
            outputs parameter $\wh{\theta}$ such that with probability $1-\delta$, 
            \[ \tv{\cD(\wh{\theta})}{\cD(\theta^\star)} \leq \eps\,.\]
            The algorithm runs in time $\poly(n)$. 
        \end{infcorollary} 
        We stress that before our work the only known algorithms in this setting are by \cite{Kontonis2019EfficientTS,lee2024unknown}, both of which require $S^\star$ to be approximable in $\nL_2$-norm.
        In contrast, the above result works with just an approximation in the $\nL_1$-norm, which is significantly weaker than an approximation in the $\nL_2$-norm.
        This change enables us to obtain the following new results from \cref{infcor:truncation}:
        \begin{center}
            \textit{The first efficient algorithm for truncated parameter estimation, when $S^\star$ is an arbitrary function of halfspaces}\footnote{
                Functions of $k$ halfspaces are a significant generalization of intersections of $k$ halfspaces and are of the form:
                $b(h_1,h_2,\dots,h_k)$ where $b$ is an arbitrary boolean function $b\colon \zo^k\to \zo$ and $h_1,h_2,\dots,h_k$ are $k$ halfspaces.
                In particular, they can capture complex \textit{non-convex} sets, while intersections of halfspaces are necessarily convex.
                } \textit{and $\hyD$ satisfies the assumptions of \cite{lee2024unknown}.}
        \end{center}
        Notably, their assumptions hold for Gaussians and product exponential distributions.


        
    \subsubsection{Detecting Truncation}\label{sec:intro:application:detection}
        Our next application is also from Truncated Statistics \cite{Cohen91,maddala1983limited}.
        In particular, to the problem of testing whether a given set of samples is \iid{} from a distribution $\cQ$ or \iid{} from the truncation $\cQ(S^\star)$ of $\cQ$ to some unknown set $S^\star$.
        
        A growing line of works in Theoretical Computer Science study this problem \cite{de2023detectingConvex,he2023testingjuntatruncation,de2024detecting}.
        It is not too hard to show (see the work of \citet*{de2024detecting}) that if the mass of $S^\star$ under $\cQ$ is very close to $1$ then this problem is hopeless.
        This is roughly because as the mass $\cQ(S^\star)$ approaches 1, the distributions $\cQ$ and $\cQ(S^\star)$ become indistinguishable. 
        Hence, one focuses on the case where $\cQ(S^\star)\leq 1-\beta$ for some known constant $\beta=\Omega(1)$.
        \citet*{de2024detecting} study this problem in the special case where $S^\star$ is a degree-$k$ polynomial threshold function (PTF) and give a $O(d^{k/2})$ sample and time algorithm that works when $\cQ$ satisfies the following:
        \begin{enumerate}[itemsep=0pt,label=(C\arabic*)]
            \item $\cQ$ is a hyper-contractive;
            \label[none]{asmp:detectingTruncation:hypercontractive}
            \item $\cQ$ is an \iid{} product distribution;\label[none]{asmp:detectingTruncation:product}
            \item $\cQ$ is \textit{known} in the following sense: we can \textit{exactly} compute an orthonormal basis of polynomials up to degree $k$ in time $O(d^{k/2})$.
            \label[none]{asmp:detectingTruncation:access}
        \end{enumerate}
        Our \cref{mainthm:algo} implies efficient algorithms for this problem for a broader set of truncation sets (beyond PTFs) and it also allows $\cQ$ to be a non-product distribution and does not require it to be known as long as one has sample access to $\cQ$ (relaxing \cref{asmp:detectingTruncation:product,asmp:detectingTruncation:access}).
        \begin{infcorollary}[see \cref{cor:detectingTruncation,cor:detectingTruncation:gen}]\label{infcor:detectingTruncation}
            Consider the setting described above with $\beta=\Omega(1)$.
            Suppose one has an efficient sampling oracle for $\cQ$.
            Fix any $\delta\in (0,\nfrac{1}{2})$ and $k\geq 1$ such that $S^\star$ is $O(1)$-approximable by degree-$k$ polynomials in the $\nL_1$-norm with respect to $\cQ$.
            There is an algorithm that given $n=d^{k}\cdot \polylog({\nfrac{1}{\delta}})$ samples from $\cQ$ or $\cQ(S^\star)$, with probability $1-\delta$, correctly detects whether the samples were from $\cQ$ or $\cQ(S^\star)$ in $\poly(n)$ time.
            
            Moreover, the algorithm also works (with the same guarantees) with an \emph{approximate} sampling oracle--that outputs samples from a distribution $\cR$ such that $\cQ$ is $(\sigma,q)$-smooth (for $q=O(1)$ and $\sigma=\Omega(1)$; see \cref{asmp:smoothness}) with respect to $\cR$. 
            (In particular, it is sufficient to have $\chidiv{\cR}{\cQ}=O(1)$; \cref{lem:closenessFromChiSquare}.)
        \end{infcorollary}
        This result, in particular, shows that detection of truncation is possible with only an \textit{approximate} sampling oracle to $\cQ$, which is significantly less restrictive than \cref{asmp:detectingTruncation:access}, and even when $\cQ$ is not a product distribution, which is much weaker than \cref{asmp:detectingTruncation:product}.
        The latter, in particular, makes progress on an open question of \cite{de2024detecting} asking for efficient truncation-detection algorithms that work for non-product distributions; see the associated STOC talk \cite{de2024detectingFOCStalk}.
        Further, since the above algorithm can handle truncation sets beyond PTFs, we get novel algorithms in the following case:
        \begin{center}
            \textit{$S^\star$ is any function of $O(1)$-halfspaces and $\cQ$ is \textit{any} log-concave distribution.}
            
        \end{center}
        Recall that functions of halfspaces can be non-convex.
        We stress that prior to our work no efficient algorithms for detecting truncation were known for non-convex hypothesis classes beyond PTFs, even when $\cQ$ was \textit{known} to be $\cN(0, I)$.
        Further, even for a single halfspace, no efficient algorithms were known to detect truncation when only sample access to $\cQ$ was available.

\subsection{Technical Overview}\label{sec:technicalOverview}
We initiate the study of learning in challenging and realistic settings where unlabeled samples are imperfect -- they can have covariate shifts and even adversarial corruptions.
Below we briefly explain the key ideas and challenges in the proofs of our main results. %


    \subsubsection{Challenges in Using Existing Approaches for PIU Learning}
        \label{sec:intro:technicalOverview:challenges}
        First, we discuss challenges in using existing techniques.

        \paragraph{{Issue I: ERM is insufficient}.}
            A standard approach to learning with positive and unlabeled samples is to reduce the problem to agnostic learning.
            In this reduction, the agnostic learning instance is created by imagining all unlabeled samples to be negative samples.
            The resulting feature distribution is a mixture of the form (for some weight $\gamma\in (0,1)$):
            \[
                \cM = \gamma \cPtrue + (1-\gamma) \cD\,.
                \yesnum\label{eq:overview:mixture}
            \]
            If the unlabeled samples are drawn from the underlying distribution $\cDtrue$ (\ie{}, if $\cD=\cDtrue$), then $\optset$ (the set of positive examples) is also the optimal solution for the agnostic learning problem and ERM is sufficient for learning.
            However, in general, when $\cD \neq \cDtrue$, the optimal hypothesis of the resulting problem can be very different from $\optset$ and, hence, ERM is not sufficient and a different approach is needed to even bound the sample complexity of PIU learning (see \cref{rem:necessityOfImproper}).
            
            
            


        \paragraph{Additional Challenges to Get Computational Efficiency.}
        To illustrate additional challenges in obtaining computationally efficient algorithms, let's suppose that $\cD$ satisfies the following strengthening of \cref{asmp:smoothness}: for all measurable sets $S$,
            \[
                \cDtrue(S)
                    \leq O(\cD(S)^{1/q})
                \qquadand
                \cDtrue(S)
                    \geq \Omega(\cD(S)^{q})
                \,,
                \yesnum\label{eq:bounded_chisquare}
            \]
            These requirements turn out to be sufficient to ensure that $\optset$ is (close to) the optimal solution of the agnostic learning problem created above and, hence, to ensure that ERM learns $\optset$.\footnote{See \cref{rem:densityLB} for a discussion of why the assumption in \eqref{eq:bounded_chisquare} is significantly stronger than \cref{asmp:smoothness}.}
            Thus, under the simplification in \cref{eq:bounded_chisquare}, the main challenge is developing an efficient algorithm to solve the agnostic learning problem created.
            However, agnostic learning is computationally challenging  \cite{guruswami2009hardness, feldman2006new, daniely2016complexity}.
            Almost all efficient algorithms for agnostic learning are based on the now-standard \lreg{} algorithm of \citet*{kalai2008agnostically}.
            As a result, all of these algorithms only work if the hypothesis class $\hyH$ is approximable by polynomials with respect to the feature distribution, say, $\cQ$.
            (The specific definition of polynomial approximation is not necessary to follow the remainder of the technical overview; the definition appears in \cref{sec:preliminaries}).
            Hence, for efficiently learning $\hyH$ one needs a bound on the degree of polynomials which approximate $\hyH$ with respect to $\cQ$.
            There are several challenges in obtaining these bounds.


            \paragraph{Issue II~~(Standard Polynomial-Approximation Guarantees Do Not Apply to $\cM$).} First, existing results on polynomial approximation are only known for special distributions $\cQ$ -- such as log-concave distributions or the uniform distribution over the Boolean hypercube.
            However, because distribution $\cPtrue$ and $\cD$ can be very different from each other, their mixture $\cM$ (\cref{eq:overview:mixture}) may not be log-concave or uniform (see \cref{fig:mixture}).\footnote{One might suspect that it is sufficient to have separate approximating polynomials $p_1(\cdot)$ and $p_2(\cdot)$ with respect to $\cPtrue$ and $\cD$ respectively.
            Unfortunately, this is not the case: for using \lreg{}, we require both of these polynomials to be the same, \ie{}, $p_1=p_2$.} 
            \begin{figure}[htb!]
                \centering
                \vspace{4mm}
                \includegraphics[width=0.75\linewidth]{figures/mixture2.pdf}
                \caption{Example of mixture distribution $\cM$ constructed by the distribution of positive examples $\cPtrue$ and the ``imperfect'' distribution of unlabeled examples $\cD$ (also see \cref{eq:overview:mixture}). 
                Here, $\cPtrue$ is a truncation of a Gaussian distribution to a convex set (an interval, namely, $\optset=[0,1]$) and $\cD$ is a Gaussian distribution.
                Nevertheless, $\cM$ is non-log-concave.
                }
                \label{fig:mixture}
            \end{figure}
            
            \paragraph{Issue III~~(Approximating Polynomial w.r.t.\ $\cD$ Can Be Poor for $\cPtrue$ {{and vice versa}}):}
                One idea to overcome this is to use the facts that (1) $\cPtrue$ is a truncation of $\cDtrue$ to a set of constant mass (namely, $\optset$; \cref{asmp:posFraction}) and (2) $\cDtrue$ and $\cD$ are close to each other (they satisfy \eqref{eq:bounded_chisquare}).
                Using this, one can show that if $p\colon\R^d\to \R$ $\eps$-approximates a hypothesis $h$ with respect to $\cD$ in $\nL_q$-norm for the constant $q$ from \eqref{eq:bounded_chisquare}, 
                then it is an $O(\eps)$-approximating polynomial with respect to $\cPtrue$ in $\nL_1$-norm.
                (The same also holds if we convert an approximating polynomial for $\cPtrue$ to one for $\cD$.)
                Hence, to ensure $\nL_1$-norm approximation with respect to both distributions, $p(\cdot)$ needs to approximate $h$ in $\nL_q$-norm for one of the two distributions.
                This is a very strong requirement, and in fact, to the best of our knowledge, no useful approximation results are known for $q \geq 3$.%
            
            \paragraph{Issue IV~~(Using Bridge Distributions Needs $\nL_2$-Approximation and May Not Exist):} 
                Another idea is to use the concept of a bridge distribution, which has appeared in some recent works (\eg{}, \cite{Kontonis2019EfficientTS,kalavasis2024transfer,lee2024unknown}).
                In some special cases, these works show that even though $\cD$ and $\cDtrue$ are far from each other -- say, they have a very large $\chi^2$-divergence -- one can construct a bridge distribution $\mu$ such that $\mu$ is close to both $\cD$ and $\cDtrue$ in the $\chi^2$-divergence (\ie{}, $\chidiv{\cD}{\mu},\chidiv{\cPtrue}{\mu}$ are small).
                (This construction crucially uses the fact that the $\chi^2$-divergence does not satisfy the triangle inequality.)
                If such a bridge distribution exists, then one can convert an $\nL_2$-approximating polynomial with respect to $\mu$  to an $\nL_1$-approximating polynomial with respect to \textit{both} $\cD$ and $\cPtrue$.
                This approach has a few downsides:
                First, it only works in restricted cases where a bridge distribution exists.
                Second, a bridge, even if exists, it might not admit approximating polynomials.
                Finally, even if the first two issues can be handled, this approach inherently require{s} an $\nL_2$-approximating polynomial to transfer an approximating polynomial with respect to $\mu$ to an approximating polynomial with respect to $\cD$ and $\cPtrue$ (using, \eg{}, the Cauchy--Schwarz inequality).
            In contrast, our result works with $\nL_1$-approximating polynomials with respect to just $\cD$.
                
            
                
        \subsubsection{Our Approach: Reduction to Constrained Learning} 
            Instead of trying to efficiently solve the above reduction to agnostic learning, we take a dramatically different approach:
            we reduce our problem to several \textit{constrained} learning problems.
            The starting point of this approach is practical heuristics used for learning from positive and unlabeled samples which, find a hypothesis $H$ that (when viewed as a set) contains the minimum number of unlabeled samples subject to selecting all positive samples \cite{bekker2020learning}.
            This can be formalized as the following program: given sets $P$ and $U$ of positive and unlabeled samples and a tolerance $\rho\geq 0$ %
            \[
                \argmin_{H\in \hyH}~~ \frac{\abs{H\cap U}}{\abs{U}}
                \,,\quad\text{ such that}\,,\quad 
                \frac{\abs{H\cap P}}{\abs{P}} \geq 1-\rho
                \,.
                \tag{\pERM{}}
            \]


        \paragraph{Sample Efficient Algorithm (see \cref{sec:sampleEfficiency}).}
            In general, the optimal solution of \pERM{} can be far from $\optset.$
            At the core, this is because \cref{asmp:smoothness} allows $\supp(\cDtrue)$ to be a strict subset of $\supp(\cD)$.
            To gain some intuition about why this introduces challenges, consider \cref{fig:intersection_halfspaces}, which shows that due to $\supp(\cDtrue)\subsetneq \supp(\cD)$, a low-VC-dimension hypothesis $H$ (say a halfspace), can appear like a high-VC-dimension halfspace (say intersections of $k$ halfspaces for large $k$) when provided samples from $\cD$.\footnote{In some of the earlier applications (\eg{}, learning from smoothened positive examples and the practical applications in \cref{sec:intro:realWorldApplications}) one cannot ensure that $\supp(\cD)=\supp(\cDtrue)$ and, so it is crucial to develop algorithms that rely solely on \cref{asmp:smoothness}.}
            To obtain a sample-efficient algorithm, our main insight is to solve many instances of the \pERM{} program. 
            Concretely, we solve $\nfrac{1}{\eps}$ instances of the program, where the instances are constructed iteratively and, roughly speaking, the $i$-th instance focuses on the region $B_{i-1}$ of the domain where the $(i-1)$-th instance is not guaranteed to perform ``well'' (by deleting all unlabeled examples which are not in $B_{i-1}$). 
            At the end we output the intersection of the hypotheses $H_1, H_2,\dots,H_{1/\eps}$ which are solutions to each instance.
            Since in practical applications one must often deal with $\supp(\cDtrue)\subsetneq \supp(\cD)$ (\cref{sec:intro:realWorldApplications}), we believe this approach of solving many instances of \pERM{} may also be of interest in practice.
            

        \begin{figure}[h!]
            \centering
            \begin{subfigure}[b]{0.3\textwidth}
                \includegraphics[width=0.85\textwidth]{figures/intersection_ex1.pdf}
                \caption{\textit{PIU Samples \\ \phantom{} }}
                \label{fig:sintersection_halfspaces_1}
            \end{subfigure}
            ~~~
            \begin{subfigure}[b]{0.32\textwidth}
                \includegraphics[width=\textwidth]{figures/intersection_ex2.pdf}
                \caption{\textit{PU Samples with $\optset$ as halfspace \\}}
                \label{fig:intersect_halfspaces_2}
            \end{subfigure}
            ~~~
            \begin{subfigure}[b]{0.32\textwidth}
                \includegraphics[width=\textwidth]{figures/intersection_ex3.pdf}
                \caption{\textit{PU Samples with $\optset$ as an intersection of six halfspaces}}
                \label{fig:intersect_halfspaces_3}
            \end{subfigure}
            \vspace{-3mm}
            \caption{
                {This figure illustrates one challenge in learning with PIU samples.
                Subfigure (a) presents the observed positive and (imperfect) unlabeled samples.
                Given these examples, it is impossible to distinguish between Scenarios (b) and (c). 
                Hence, in particular, due to imperfections in the unlabeled samples a hypothesis with low VC-dimension (\eg{}, a single halfspace in (b)) can appear like a high-VC-dimension hypotheis (\eg{}, intersections of many halfspaces in (c)).}
            } 
            \label{fig:intersection_halfspaces}
        \end{figure}

        \paragraph{Computationally Efficient Algorithm (see \cref{sec:computationalEfficiency}).}
            A natural idea is to convert a sample-efficient algorithm into a computationally efficient one by solving \pERM{} efficiently.
            However, one cannot use techniques for designing efficient learning algorithms in the realizable setting to solve this problem since the optimal value  of \pERM{} is generally bounded away from zero (so does not fit the standard template for efficient realizable learning algorithms that reduce the problem to checking feasibility of, \eg{}, a linear program).  
            At this point, given the challenges we discussed for efficiently implementing ERM (in the previous sub-section), efficiently implementing \pERM{} might appear hopeless: 
            instead of solving an unconstrained learning problem -- which is already hard -- we now have to solve a constrained learning problem -- that is seemingly harder.

            A key advantage of this approach is that it decouples the samples from $\cD$ and $\cPtrue$.
            Therefore, while solving this program, we avoid dealing with mixture distributions, which we have seen are challenging to deal with.
            We note that this is why \cref{infthm:main} only requires $\nL_1$-approximation with respect to $\cD$, which we have sample access to.
            This contrasts the attempts we discussed earlier in \cref{sec:intro:technicalOverview:challenges} that require a stronger approximation guarantee (in $\nL_q$-norm, for which no results are known when $q \geq 3$) or use the existence of some ``bridge'' distributions that are only known to exist for specific distribution families (\eg{}, Gaussians).  
            In contrast, our algorithm only needs an approximation guarantee in $\nL_1$-norm (regardless of the value of $q$ in \cref{asmp:smoothness}) and makes no further assumptions about the distribution family $\cD$ and $\cDtrue$ belong to (in particular, they can be non-Gaussian). 

            
            Next, we describe how we design an efficient approximation algorithm to solve (one instance of) \pERM{}. 
            This combined with the iterative sample-efficient algorithm described above is in-spirit sufficient to complete the proof of \cref{thm:main}.
            That said, compared to the analysis for sample complexity, we need to account for some additional challenges.
                First, we need to ensure that one can efficiently obtain samples from each of the PIU learning instances constructed (\cref{lem:main:rejectionSampling}). 
                Second, to solve each instance efficiently, we need to ensure that the $L_1$-approximation guarantee is maintained even though the $i$-th instance focuses on the subdomain $\R^d\setminus B_{i-1}$, which is determined by the $(i-1)$-th instance.
                Finally, since we only solve each \pERM{} instance approximately, it is possible that after discarding unlabeled samples within $B_{i-1}$, it no longer holds that $\supp(\cD)\supseteq \supp(\cPtrue)$, resulting in invalid PIU learning instances and we need to adapt the construction to avoid this (\cref{lem:main:assumptions}).

        \paragraph{An Approximation Algorithm for \pERM{} (see \cref{thm:constReg} and  \cref{sec:efficientPERM}).} 
            Up to this point, we have made several careful decisions to arrive at the program we want to solve; while avoiding several bottlenecks in other natural approaches. 
            Now, we are ready to introduce the final ingredients of our approach. 
            The final component of our proof is a constrained version of \lreg{} that ``approximately'' solves \pERM{}.
            Constrained versions of \lreg{} have appeared in other works, in particular for efficient reliable learning, which, roughly speaking, require classifiers to make no errors when they make a prediction but allows them to abstain from making a prediction on a fraction of the samples \cite{kanade2014independentRealiable,kalai2012reliable}.
            We provide a new constrained version of \lreg{} -- which we call \constrainedreg{} -- along with novel analysis showing that this algorithm is a pseudo-approximation algorithm for \pERM{} whenever $\hyH$ can be approximated by polynomials in $\nL_1$-norm with respect to $\cD$.
            
            At a high level, the algorithm operates as follows: it solves a constrained polynomial regression problem to find a polynomial $p(\cdot)$ and outputs a polynomial threshold function $\wh{h}(\cdot)=\mathds{1}\sinbrace{x \colon p(x)\geq t}$ constructed by thresholding $p(\cdot)$ at a suitable value $t$.
            The proof of correctness leverages ideas from the analysis of \lreg{}, but needs to account for some technicalities:
            First, we show that as long as $t$ is suitably bounded away from 1, then using standard uniform convergence arguments $\wh{h}$ is feasible for \pERM{}.
            This adjustment, however, is inconsistent with \lreg{} that selects the best threshold in $[0, 1]$.
            We show that \lreg{}'s analysis is robust to this change.
            Second, unlike \lreg{}, we aim to show that \constrainedreg{} outputs a hypothesis that is feasible for a constrained optimization problem (\pERM{}); and significant additional analysis is needed to establish feasibility and pseudo-optimality.
        
        \subsubsection{Learning in the List Decoding PIU Model}
            The List Decoding PIU Model (\cref{def:list}) is even more challenging than the previous model because we lack sample access to a distribution close to $\cDtrue$. A natural idea is to combine the list of distributions $\cD_1,\dots,\cD_\ell$ to approximate $\cDtrue$, for example via the uniform mixture $\cM$. This approach is promising since one can show that $(\cDtrue,\cM)$ satisfies \cref{asmp:smoothness} with $\sigma'=\sigma \ell^{-1/q}$ and $q'=q$: indeed, for any measurable set $S$, we have $\cM(S)^{1/q}\geq \ell^{-1/q}\cD_{i^\star}(S)^{1/q}\geq \sigma\ell^{-1/q}\cDtrue(S)$, where the last inequality holds as $(\cDtrue,\cD_{i^\star}$) satisfies \cref{asmp:smoothness}. While this suffices for bounding sample complexity, achieving computational efficiency is challenging because we lack guarantees on $\cD_i$ for $i\neq i^\star$, causing polynomial approximability of $\optset$ with respect to $\cM$ to fail. Other approaches that combine $\cD_1,\dots,\cD_\ell$ to yield a distribution $\cM$ close to $\cDtrue$ in the sense of \cref{asmp:smoothness} while retaining polynomial approximability also encounter the failure of polynomial approximation. 
            Instead of trying to combine the list of distributions, we run the algorithm from \cref{infthm:main} separately on each unlabeled distribution $\cD_i$ to obtain hypotheses $H_1,H_2,\dots,H_\ell$ and output the intersection of these hypotheses $\bigcap_i H_i$.
            
            
            
            
             
            
            
 
    \subsection{Practical Motivation for PIU Learning}
    \label{sec:realWorldApplications}
    \label{sec:intro:realWorldApplications} 
    \noindent While we have highlighted our results insofar as how they apply to other learning theory problems, %
    we next motivate the study of the PIU learning model itself by discussing real-world applications of PU learning where obtaining ``perfect'' unlabeled examples is difficult or even impossible and, hence, PU learning methods must grapple with \textit{imperfect} unlabeled data -- precisely the challenge that the PIU learning model attempts to formalize. %
    For further examples, see the surveys by \citet{elkan2008pu,plessis15convex} and a recent book \cite{jaskie2022positive} on PU learning. %
    
    \paragraph{Applications in Bioinformatics and Medicine.} %
    Consider \textit{disease-causing gene identification} -- a central task in Genetics and Molecular Biology with major healthcare implications \cite{mordelet2011prodige,yang2012positive}.
    Identification is very expensive due to its reliance on highly skilled human labor and precise experimental procedures. 
    Here, ML algorithms, if given sufficient and good quality data, can expedite discovery by flagging promising candidates for experimental validation.
    However, due to the high false negative rates in experimental methods, high-quality negative examples are scarce; a gene that appears benign under one condition may be disease-causing under another \cite{reviewPUBioformatics2021}.

    Similar issues affect a wide array of applications in Bioinformatics and Medicine -- from biological sequence classification, to drug interaction identification, to identifying undiagnosed diseases from health records \cite{elkan2008pu}.
    In response, starting from the seminal works of \cite{hur2006proteinProtein,elkan2008pu}, a substantial body of work designs and uses heuristics for PU learning; a partial list is as follows \cite{ming2016pupylation,pengyi2019adaSampling,zhao2008geneFunction,xiao2008niologicalSequence,yang2012positive,luigi2010geneRegulatory}.
    For an in-depth discussion on the challenges and promise of PU learning across several Bioinformatics applications, see \citet{reviewPUBioformatics2021}.
    


    \itparagraph{Reasons for Imperfect Unlabeled Examples: Source Mismatch.}
        The feature distribution of unlabeled samples often differs from that of the positive examples when drawn from different sources \cite{jaskie2022positive}.
        For instance, validated disease-causing genes (positive examples) come from specialized studies, while unlabeled examples may be sampled from all known genes \cite{mordelet2011prodige,reviewPUBioformatics2021}.
        Similarly, in medical applications, positive data might be collected from specialized and costly programs at a single hospital, whereas unlabeled data are aggregated nationwide, leading to discrepancies in feature distribution (\eg{}, demographics and measurement conditions) \cite{Hassanzadeh2018ClinicalDC}. 
        In both illustrative examples, the imperfect unlabeled distribution has a strictly larger support (\eg{}, its the set of all known genes) than the (perfect) unlabeled distribution (\eg{}, which are supported on only relevant genes), so the upper bound in \cref{asmp:smoothness} is natural.

        
        
        
    
  \paragraph{Applications in User-Specified Identification Tasks.}
      User-specified identification tasks -- also known as one-vs-rest classification, inlier-based outlier detection, and anomaly detection -- involve a user labeling a \textit{small} set of positive examples from a large pool of unlabeled data.
      Examples include automated face tagging, email classification, and remote sensing\footnote{Remote sensing identifies images with target concepts (\eg{}, roads, rivers) from Earth and planetary imagery.} \cite{plessis15convex,bekker2020learning,wenkai2011remoteSensingPU}.

      \itparagraph{Reasons for Imperfect Unlabeled Examples.}
      Here, reasons for covariate shift include mismatch in the time of collection of positive and unlabeled samples and corruptions in the unlabeled samples.
      In more detail, in some cases, unlabeled data is augmented with data generated by generative AI methods \cite{doi:10.1126/science.adi6000,rubin2004multiple,raghunathan2021syntheticData,jordon2022synthetic}, which
            can introduce complex or even adversarial shifts.
            Nevertheless, tools from the from the learning-with-corrupted-data literature can help identify a \textit{list} of unlabeled sources that ``approximates'' the true distribution (\cref{rem:list}); that turns out to be sufficient for learning (\cref{sec:intro:application:listDecoding}).
      Again in this setting, since unlabeled data is augmented, we expect the imperfect unlabeled samples to have a larger support than the perfect unlabeled samples -- making the bound in \cref{asmp:smoothness} natural.
      %
         


    \subsection{Takeaways and Open Problems}\label{sec:takeawaysOpenProblems}\label{sec:intro:takeawaysOpenProblems}
        We develop the first efficient algorithm for the problem of learning with positive and imperfect unlabeled (PIU) samples. 
        PIU learning not only seems like an important problem on its own as it reflects many practical settings (\cref{sec:realWorldApplications}), but it also has surprising connections to learning with positive examples and truncated statistics. 
        
        This unified view leads to efficient algorithms that work under significantly weaker assumptions than those required by existing algorithms in these applications (\cref{sec:intro:application}). 
        This applicability to several seemingly unrelated problems is an important aspect of our results, that we hope will inspire further research.
        To this end, we identify and present some open problems that arise from our work before discussing further related works, formal statements of our results, and proofs. 
        
        Our algorithm is a modification of the \lreg{} algorithm, and proving its correctness required significant new analysis.
        In light of this, our first open problem is:
        \begin{openproblem}
            \textit{Is there a black-box reduction from efficient agnostic learning to efficient PIU learning?}
        \end{openproblem}
        Such a reduction will, in particular, demonstrate that the computational complexity of PIU learning is at most as hard as that of agnostic learning.
         
        
        Next, while for a general hypothesis class, exponential dependence on $\nfrac{1}{\eps}$ is necessary for any SQ algorithm (see \cref{rem:SQlowerbound}), it may be possible to design polynomial-time algorithms for specific hypothesis classes and distribution families for $\cD$ and $\cDtrue$. 
        In this direction, we leave the following open problems whose counterparts in the realizable PAC learning model are already solved (by \citet{vempala2010convexGaussian}  for the first problem and by the folklore implementation of ERM via linear programming for the second problem \cite{mohri2018foundations}).
        \begin{openproblem}         
             Is there an algorithm that in the PIU learning model finds an $\wh{h}$ satisfying the guarantee of \cref{thm:main} in $\poly(\nfrac{d}{\eps})$  time for the following concept classes $\hyH$ and distributions $(\cDtrue,\cD)$?
            \begin{itemize}[itemsep=0pt]
                \item[$\triangleright$] $\hyH$ is the intersections of $O(1)$ halfspaces and $\sinparen{\cDtrue,\cD}$ are unknown Gaussian distributions; and 
                \item[$\triangleright$] $\hyH$ is a single halfspace and $\sinparen{\cDtrue,\cD}$ are arbitrary distributions.
            \end{itemize}
        \end{openproblem} 
        Even in the special case where $\cDtrue$ and $\cDgiven$ are Gaussian and $\hyH$ is the family of Halfspaces (a special case of Problem 2.1), this might require a different approach than solving \pERM{}, as we believe solving \pERM{} is \np{}-hard in this case.
        Indeed, in this special case, a $\poly(\nfrac{d}{\eps})$ algorithm is known, but it utilizes a very different moment-based approach \cite{lee2024unknown}. %
        
        Next, consider PIU learning model where smoothness holds with $q=1$.
        \cref{thm:sampleComplexity} shows that the sample complexity is at most $\wt{O}(\nfrac{\vc{}(\hyH)}{\eps^2})$.
        Since PIU learning is harder than realizable learning, we also know that at least $\wt{\Omega}(\nfrac{\vc{}(\hyH)}{\eps})$ samples are required.
        This raises the question:
        \begin{openproblem}
            What is the optimal sample complexity of PIU learning under smoothness with $q=1$?
        \end{openproblem}
        In this work, we study the standard uniform rates model.
        However, this can be overly pessimistic because it allows the worst-case data distribution to vary with the target accuracy $\varepsilon$, even though in practice the data distribution remains fixed as one increases the sample size to achieve higher accuracy (see \cite{bousquet2021universal}).
        Motivated by this observation  \citet*{bousquet2021universal} studied the learning curves in the \textit{universal rates model} for the PAC setting, and revealed a surprisingly different landscape of rates.
        Therefore, the following natural question arises: 
        \medskip

        \noindent \textbf{Open Problems 4 and 5.}~\emph{What is the characterization of universal rates for binary classification in the PU learning model?
            What is the characterization of the same in the PIU learning model?}
        

        
 
    \subsection{Related Work}\label{sec:intro:relatedWorks}
        As we indicated, our results are related to several learning models. We present a brief discussion of them below and present more discussion on the following and additional  works in \cref{sec:appendix:related}.
    
         \paragraph{Learning From Positive Samples Only.} \citet{natarajan1987learning} initiated work on learning with positive-only samples and showed that proper learners cannot learn even simple classes such as rectangles in 2-dimensions with positive-only samples; our \cref{thm:impossibilityPositiveSamples} extends this improper learners.
         On the positive side, under the assumption that the features are  Gaussian, \citet{Kontonis2019EfficientTS,lee2024unknown} present algorithms that learn binary concepts from positive samples assuming that the concepts have bounded \textit{Gaussian Surface Area} and the set of positive samples has at least a constant mass (\ie{}, \cref{asmp:posFraction} holds). 
        %
        Compared to these works, our results work for larger concept classes and extend beyond Gaussians.
        %

            
 

        \paragraph{Smoothened Analysis.} 
            \citet{spielman2004smooth} initiated the study of algorithms whose inputs are perturbed by Gaussian noise by showing that the simplex algorithm runs in polynomial time in this model. 
            %
            Since then, the \textit{smoothened analysis} model has found applications in many areas, often enabling significantly stronger theoretical guarantees than the worse-case model. 
            %
            For instance, in online learning, a line of works \cite{haghtalab2022efficientSmooth,haghtalab2020smooth,haghtalab2024smoothJACM,block2024performanceempiricalriskminimization,blanchard2024agnostic} showed than if the adversary is restricted to select strategies from smooth distributions then the lower bounds based on the Littlestone dimension can be bypassed and any VC class becomes learnable.
            In the same spirit \citet*{pmlr-v247-chandrasekaran24a} obtained a wide-array of faster PAC learning algorithms when the learner only has to compete the best classifier that is robust to Gaussian perturbations of its inputs.
            %
            We refer the reader to \cite[Chapter 13]{roughgarden2021beyond} and \cite{beier2004typical} for further discussion and applications outside of learning theory.
            
             

        \paragraph{Learning with Covariate Shift.}
            Most existing works in learning with covariate shift \cite{yishay_2009_domain_adapt,yishay_2009_multi_source_renyi,theory2010david,pmlr-v9-david10a,sugiyama2012machine,kalai2021arbitraryCS,pmlr-v75-kpotufe18a,pathak2022new,goel2024tolerantalgorithmslearningarbitrary,Weiss2016transfersurvey}
            have assumptions and how much shift the distribution can have, often related to \cref{asmp:smoothness}. For instance, the notion of a transfer exponent in \cite{hanneke_kpotufe_2019_transfer_exp,hanneke2023efficient,hanneke2024unifiedtheorytransferlearning} is equivalent to \cref{asmp:smoothness} for realizable classification. However, in existing works, one has access to negative examples, and the algorithmic solution is often to first efficiently learn on given samples and then transfer the error to the ``target'' distribution via a change of measure.
            This is not possible in our setting and it is unclear how to learn efficiently in the first place. This is our main contribution.

         
    
        \paragraph{Reliable Learning.} 
            Another related line of works studies reliable learning \cite{kalai2012reliable,kalai2021arbitraryCS,kanade2014independentRealiable,goel2017reliableRELU,durgin2019reliableCSPs,goldwasser2020beyondPurturbations,diakonikolas2024reliable}, a model introduced by \citet{kalai2012reliable}.
            This model captures scenarios where one type of error, \eg{}, false positives, is costlier than another type, \eg{}, false negatives. 
            Some works \cite{kalai2012reliable,kanade2014independentRealiable} on fully reliable learning (which requires the learner to make almost no errors but allows it to abstain) design constrained \lreg{} algorithms, that abstain on some examples. 
            While there are several differences between our work and theirs, a key difference is that they have access to both positive and negative samples and, hence, their constraint is defined with respect to the same distribution as their objective. 
            In our setting, however, we work without negative samples and our unlabeled samples have a different distribution than the ``true'' feature distribution; apart from significant differences in the analysis, this change also prohibits proper learning (see \cref{rem:necessityOfImproper}) and requires us to run multiple carefully constructed instances of the constrained version of \lreg{} we develop.


        \paragraph{List-Decodable Learning.}
            There is a growing array of list-decodable learning algorithms that continue to work even when a large fraction of their inputs are corrupted \cite{charikar2017learning,karmalkar2019list,Cherapanamjeri2020DecodableNearlyLinear,bakshi2021Subspace,Abhimanyu2023DecodableBatches,Diakonikolas2023DecodableCovariance,kothari2022decodableCovariance,Raghavendra2020Decoding,Diakonikolas2022DecodableSparseMean,Diakonikolas2021SQ_LB_Decoding,Diakonikolas2021PCA_Decodable}.
            Our results on the list-decoding PIU model (\cref{def:list}) enables us to utilize these algorithms to enable learning even in scenarios where a large fraction of the unlabeled samples are corrupted.
            We refer the reader to \cref{sec:appendix:related} for a more extensive discussion of list-decoding algorithms.
            
            
       
    \section{Preliminaries}\label{sec:preliminaries}
    
    \textbf{Notation.}\quad 
    Given a point $x\in \R^d$ and a set $S\subseteq\R^d$, we use $\mathds{1}_S(x)$ and $\ind\{x \in S\}$ to denote the indicator that $x\in S$.
    We use standard notation for vector and matrix norms:
    For a vector $z\in \R^d$ and $p\geq 1$, the $\nL_p$-norm of $z$ is $\norm{z}_p\coloneqq \inparen{\sum_i \abs{z_i}^p}^{\sfrac{1}{p}}$ and, taking limit $p\to\infty$ implies, $\norm{z}_\infty=\max_i \abs{z_i}$.
    Given a distribution $\cD$ over $\R^d$ and a measurable set $S\subseteq\R^d$, $\cD(S)$ denotes the mass of $S$ under $\cD$, \ie{}, $\Pr_{x\sim \cD}[x\in S]$.
    We overload this notation at times in applications in truncated statistics (in \cref{sec:intro:application,sec:application:unknownTruncation,sec:application:detectingTruncation}) to also mean the truncation of a distribution $\cD$ to the measurable set $S$, \ie{}, if $\d\cD$ is the density of $\cD$, $\cD(S)$ is the conditional distribution $\propto \ind_S(x) \cdot \d\cD(x)$. 
    The difference between the two would always be clear from context.

    \paragraph{Hypothesis Classes and PTFs.} 
    We use both Boolean functions (\eg{}, $h\colon \R^d \rightarrow \inbrace{0,1}$) and sets $(H \subseteq \R^d)$ to denote hypotheses, with the relation $H = \inbrace{x \in \R^d \colon h(x) = 1}$. 
    Our algorithm, like most efficient learning algorithms, to learn general hypothesis classes in non-realizable settings, relies on polynomial threshold functions (PTFs).
    Hence, PTFs frequently show up in our analysis.

    \begin{definition}[Polynomial Threshold Functions (PTF)]\label{def:ptfs:polynomials}
        Given $k\geq 1$, let $\hyP(k)$ be the collection of all polynomial threshold functions of degree-$k$.
        A polynomial threshold function (PTF) of degree $k$ is a function defined by a polynomial $p\colon \R^d \rightarrow \R$ of maximum degree $k$ and a threshold $t \in \R$ in the following sense:
        $H({p,t})\coloneqq \mathds{1}\sinbrace{p(x)\geq t}$. 
    \end{definition}

    \paragraph{Positive and Unlabeled Samples.}
        Given a hypothesis class $\hyH$ and a target function $\hstar$, we use $\optset$ -- which is the set corresponding to $\hstar$ -- to denote the set of positive examples.
        We use $\cPtrue$ to denote the distribution of positive covariates or features, which satisfies $\supp(\cPtrue)=\optset$.
        $\cDtrue$ denotes the distributions of all features (regardless of whether their label is positive or not).
        $\cPtrue$ and $\cDtrue$ are related: $\cPtrue$ is the truncation of $\cDtrue$ to $\optset$ -- the set of positive examples.

    \paragraph{Divergences Between Distributions.}
        In our applications, we often need bounds on the ``closeness'' of two distributions.
        For this, we recall some common notions. 
        Consider two distributions $\cP$ and $\cQ$.
        The total variation distance between $\cP$ and $\cQ$ is $\tv{\cP}{\cQ}=\inparen{\nfrac{1}{2}}\int_{x}\abs{\d \cP(x)-\d \cQ(x)}\d x$.
        The $\chi^2$ divergence from $\cP$ to $\cQ$ is  $\chidiv{\cP}{\cQ}
        = \int \inparen{\nfrac{\d\cP}{\d\cQ}}^2 \d \cQ - 1$.
        Similarly, the Rényi divergence of order \(r \geq 1\) from \(\cP\) to \(\cQ\) is defined as $\renyi{r}{\cP}{\cQ} = \frac{1}{r - 1} \log \int (\sfrac{\d \cP}{\d \cQ})^r \d \cQ$.


     
 

\section{PIU Learning Model and Main Results}\label{sec:model_and_result}
    In this section, we formally define the model of learning from positive and imperfect unlabeled samples (PIU learning) and state our main results. 

    \subsection{Learning with Positive and Imperfect Unlabeled Samples}\label{sec:model}
        %
        To begin with, we have a domain $\cX$ of covariates or features that is a subset of $\R^d$, either discrete or continuous, and a distribution $\cDtrue$ over the domain.
        %
        We are interested in binary classification and operate in the realizable setting, where for each feature $x$ there is a corresponding label $y(x)\in \zo$, that is realized by an unknown Boolean function $h^\star\colon \cX\to \zo$.
        $h^\star$ can be arbitrary and the only assumption we make is that there is a known hypothesis class $\hyH$ (which is a collection of Boolean functions $h\colon \cX\to \zo$) that contains $h^\star$.
        Of specific interest is the set of samples with a positive label and, for convenience, we use $\optset$ to denote it, \ie{}, 
        \changetag{Support of Positive Examples) \ (#1} 
        \[
            \phantom{ABCDEFGHIJK}
            \optset\coloneqq \inbrace{x\in \supp(\cD)\colon h^\star(x)=1}\,.
            \yesnum\label{eq:setOfPositiveExamples}
        \]
        \changetag{#1}
        
        \paragraph{Positive Examples.}
        As discussed in \cref{sec:intro}, in a number of applications, high-quality data sources only provide positive samples $x_1,x_2,\dots,x_n$ where each $x_i$ is drawn from $\cDtrue$ conditioned on lying in the set of positive examples, \ie{}, $x_i\in \optset$. 
        Of course, without loss of generality, one can also augment these samples with labels $(x_1,y_1),\dots,(x_n,y_n)$ where $y_1=\dots=y_n=1$.
        At first glance, these samples seem to contain all the necessary information: one can assume that any feature $x$ that we have not seen (in the high-quality data source) has a negative label.
        In low-dimensional settings ($d\ll n$), this is indeed the case, however, in high-dimensional settings most of the features remain unobserved and one needs information about the distribution $\cD$ outside $\optset$ to select a good hypothesis or classifier $h\in \hyH$.
        This intuition was formalized by \citet{natarajan1987learning}: Natarajan studied binary classification with \textit{only} positive samples (as our high-quality source provides) and showed that almost no hypothesis classes are learnable by any proper learner; indeed, even 2-dimensional halfspaces, cannot be learned from just positive examples.
        Extending this result, we show that this lower bound also applies to improper learners for many classes including two-dimensional halfspaces and one-dimensional intervals (see \cref{thm:impossibilityPositiveSamples}).

        \paragraph{Unlabeled Examples.}
            Thus, to enable learning in these applications, one turns to low-quality data sources (such as public medical datasets or the web).
            These low-quality sources provide samples $x_1,x_2,\dots,x_n$ drawn from some distribution $\cDgiven$ different from $\cDtrue$ (\ie{}, $\cDgiven\neq \cDtrue$).
            These samples are \emph{unlabeled} in the sense that we do not have any guarantee about their label, \ie{}, whether $x_i$ is in $\optset$ or not.
            Of course, if $\cDgiven$ has no relation to $\cDtrue$, then these samples provide us with little relevant information.
            Therefore, we need to assume some relation between $\cDgiven$ and $\cDtrue$ to have any hope for learning.
            In most of the paper, we consider the following relation between $\cDgiven$ and $\cDtrue$.
            \smoothness*
            \noindent 
            It is useful because it allows us to upper bound the amount of mass that the true distribution $\cDtrue$ assigns to any {(measurable)} set $S$ based on the mass $\cDgiven$ assigns to $S$.
            This generalizes the notions of smoothness considered by recent works on smoothened online learning \cite{haghtalab2020smooth,haghtalab2021smooth,haghtalab2022efficientSmooth,haghtalab2024smoothJACM,blanchard2024agnostic}; concretely, their notion of smoothness is equivalent to the special case $q=1$.


            \begin{remark}[Extensions to More General Notions of Smoothness]
                \label{rem:generalSmoothness}
                The techniques we use to bound the sample complexity and computational complexity of PIU learning extend beyond \cref{asmp:smoothness}.
                They suffice for PIU learning under the following weakening of \cref{asmp:smoothness}:
                there is a known rate function $r\colon \R_{\geq 0}\to \R_{\geq 0}$ satisfying $\lim_{m\to 0^+} r(m)=0$ such that the pair of distributions $\inparen{\cD, \cDtrue}$ satisfy
                \[
                    \text{for any measurable set $S$}\,,\quad 
                    r(\cDtrue(S))\leq \cD(S)\,.
                    \yesnum\label{eq:mostGeneral:smoothness}
                \]
                This captures \cref{asmp:smoothness} when $r(m)=(\sigma m)^{q}$.
                This additional generality also allows our framework to extend to settings where we can only ensure $\kl{\cDtrue}{\cD}=O(1)$ (see \cref{sec:klBoundsSatisfyGeneralSmoothness} for details).
                This is important as \cref{asmp:smoothness} does not hold in this setting for any finite $\sigma$ or $q$.
                Regarding results under \eqref{eq:mostGeneral:smoothness}, using our techniques one can deduce that the sample complexity of the PIU learning is $O\binparen{\poly{(r(\eps))}\cdot \inparen{\vc{}(\hyH)+\log{\nfrac{1}{\delta}}}}$ and the computational complexity is $\poly(d^k,\nfrac{1}{\zeta},\log{\nfrac{1}{\delta}})$ for $\zeta\leq r(\poly(\alpha,\eps))$.
            \end{remark}
            

        
        \paragraph{Main Problem.}
        Now, we are ready to state the main task we study in this work.
        \begin{problem}[PIU Learning]\label{prob:PUlearning}
            An instance of PIU Learning is specified by a hypothesis $h^\star\in \hyH$, that induces a set of positive samples $\optset$ (\cref{eq:setOfPositiveExamples}), a distribution $\cDtrue$ of features, and a (low-quality) distribution $\cDgiven$ that is close to $\cDtrue$ as specified by \cref{asmp:smoothness}.
            One has access to two streams of samples:
            \begin{enumerate}
                \item \emph{(High-quality Stream)}\quad Samples $x_1,x_2,\dots \sim \cPtrue$ where $\cPtrue$ is the truncation of $\cDtrue$ to $\optset;$ and 
                \item \emph{(Low-quality Stream)}\quad Samples $x_1,x_2,\dots \sim \cDgiven$. 
            \end{enumerate}
            The goal is to output a hypothesis $h\colon \cX\to \zo$ such that, with high probability, $h$ is close to $h^\star$ in the usual sense of PAC-learning: with probability $1-\delta$, $\Pr_{\cDtrue}\inparen{h^\star(x)\neq h(x)} \leq \eps.$
        \end{problem}
        Our computationally efficient algorithm  will also require the following mild assumption.
        \posFraction*
        \noindent This assumption requires that the fraction of positive examples have at least a constant mass under $\cDtrue$.
        This assumption, by now, is standard in the works on learning with positive and unlabeled examples (\eg{}, \cite{lee2024unknown}). %
        That said, this assumption can be dropped if one can estimate $\Pr_{\cDtrue}\inparen{h^\star(x)\neq h(x)}$; as if we can estimate $\Pr_{\cDtrue}\inparen{h^\star(x)\neq h(x)}$ and, in particular, check if it is at most $\eps$, as then we can simply output the all-zero hypothesis and only incur an $\eps$ error.
            
        \paragraph{Polynomial Approximation.}
            The other ingredient required by our efficient algorithm is a bound on the complexity of $\hyH$.
            The core notion of complexity of concept classes we consider is their approximability by polynomials, which was introduced by \citet*{kalai2008agnostically}.
            This notion of complexity is a very widely studied condition that is the backbone of most of the efficient learning algorithms that deal with agnostic noise \cite{kalai2008agnostically,kane2013learning,Blais:2010aa,10.1145/1806689.1806763,10.1007/978-3-642-15369-3_44,10.1145/2395116.2395118,5670941}. 
            %
            Formally, it is defined as follows.
            \begin{definition}[Polynomial Approximability]\label{def:poly_approx}
                Fix a concept $h\colon \R^d \to \inbrace{0,1}$, distribution $\cD$ over $\R^d$, numbers $k,p\geq 1$, and constant $\eps>0$.
                $h(\cdot)$ is said to be $\eps$-approximable by degree-$k$ polynomials in $\nL_p$-norm with respect to $\cD$ if there exists a degree-$k$ polynomial $f\colon \R^d\to\R$ such that 
                \[
                    \inparen{\Ex_{x\sim \cD}{\abs{f(x)-h(x)}^p}}^{\sfrac{1}{p}} < \eps\,.  
                \]
                (For sets, we simply replace $h(x)$ with the indicator of the set, \ie{}, $\ind_H(x)$.)
                Moreover, a hypothesis class $\hyH$ is said to be $\eps$-approximable by degree-$k$ polynomials in $\nL_p$-norm with respect to a family of distributions $\hyD$ if, for each $h\in \hyH$ and $\cD\in \hyD$, $h$ is  $\eps$-approximable by degree-$k$ polynomials in $\nL_p$-norm with respect to $\cD$. 
            \end{definition}
            \begin{table}[ht!]
            \centering
            \begin{tabular}{c c c}
             Hypothesis Class $\hyH$
                & Distribution Family & 
                    Upper bounds on $k$ \\
            \midrule
            Degree-$t$ Polynomial Threshold Functions  \cite{kane2011gsa} &
                Gaussian & 
                        $O(t^2/{\eps^4})$ \\
            Intersections of $t$ Halfspaces \cite{kalai2008agnostically} & 
                    Gaussian  & 
                        $O({\log t}/{\eps^4})$\\
            General Convex Sets \cite{ball1993gsaConvexSets} &
                 Gaussian & 
                    $O(\sqrt{d}/{\eps^4})$ \\
            \hline
            Arbitrary Functions of $t$ Halfspaces \cite{kane2013learning}
                & Log-Concave & {$\exp\sinparen{O\sinparen{t^4/\eps^4}}$}\\
            \hline
            Decision trees of size $s$
            \cite{klivans2023testable} 
                & Uniform on $\zo^d$ 
                    & ${O}(\log(s/\eps))$ \\
            Intersections of $t$ Halfspaces 
            \cite{gopalan_2010_fooling}
                & Uniform on $\zo^d$ 
                    & $\tilde{O}(t^6/\eps^2)$ \\
            Decision tree of halfspaces of size $s$, depth $t$
            \cite{gopalan_2010_fooling}
                & Uniform on $\zo^d$ 
                    & $\tilde{O}(s^2 t^6/\eps^2)$ \\
            \end{tabular}
            
                \caption{
                    \textit{
                        Upper bounds on the degree of polynomials required for $\eps$-approximating $\hyH$ in $\nL_1$-norm.
                    }
                }
                \label{tab:gsa:GaussianSurfaceArea}
        \end{table} 
        
            \noindent Our algorithms will always work whenever the target $h^\star$ (or equivalently, the set $\optset$ defined by $h^\star$) is approximable by polynomials with respect to  $\nL_1$-norm.
            It is worthwhile to note that if a hypothesis is approximable in $\nL_2$-norm, then it is also approximable in the $\nL_1$-norm because for any random variable $Z$, $\Ex[\abs{Z}]\leq \sqrt{\Ex[Z^2]}$.
            Hence, $\nL_1$-approximability is a weaker requirement than $\nL_2$-approximability.
            In fact, it can be significantly weaker: there are hypothesis classes and distributions for which $\nL_1$-approximation guarantees are known but \textit{no} $\nL_2$-approximation guarantees are known; we refer the reader to the work of \citet{kane2013learning} for an extensive discussion about this.
 
    \subsection{Main Results for PIU Learning}\label{sec:mainresult}
        In this section, we present our main results for PIU learning (\cref{prob:PUlearning}). %
        Our first result, \cref{thm:sampleComplexity}, bounds the sample complexity and we restate it below.  
        \sampleComplexity*
        \noindent As mentioned in \cref{sec:intro}, this result shows that VC dimension also characterizes PIU learning.
        The proof of \cref{thm:sampleComplexity} appears in \cref{sec:perm}.


        A recent work \cite{lee2024unknown} designed a robust version of \citet{denis1990positive}'s reduction from PU learning to agnostic learning that works when $\cD\neq \cDtrue$ but the pair $(\cD,\cDtrue)$ satisfy \cref{asmp:smoothness} and an additional assumption (see \cref{asmp:densityLB}).
        This results in a sample complexity of {$O(\vc{}(\hyH)\cdot \eps^{-2(1+q)})$}.
        The above result improves upon this sample complexity for all $q$ while only requiring \cref{asmp:smoothness}; see \cref{tab:sample}.

        \begin{remark}[Proper vs.\ Improper Learning]
            The learner in \cref{thm:sampleComplexity} is improper: it outputs a hypothesis that is an intersection of $(\nfrac{1}{\eps})$-hypotheses in our original class. 
            This is necessary as simple examples show that proper PIU learning is impossible (\cref{rem:necessityOfImproper}).
        \end{remark}


        
        
        \noindent Our next result bounds the computational complexity of PIU learning.
        \begin{theorem}[Computationally Efficient Algorithm]\label{thm:main}
            Let \cref{asmp:smoothness,asmp:posFraction} hold.
            Fix any $\eps,\delta\in (0,\nfrac{1}{2})$ and $k$ such that degree-$k$ polynomials $\zeta$-approximate $\hyH$ with respect to $\cDgiven$ in $\nL_1$-norm for $\zeta\leq O(\alpha^{8q})\cdot (\sigma\eps\alpha)^{5q^2}$.
            Then, there is an algorithm (\cref{alg:main}) that, given 
                any $\eps,\delta\in (0,1)$, 
                the constants $\sigma,q$ from \cref{asmp:smoothness},and 
                $n=\poly(d^k,\nfrac{1}{\zeta},\log{\nfrac{1}{\delta}})$ samples from each stream in PIU Learning model (\cref{prob:PUlearning}), 
            outputs a hypothesis ${h}\colon \cX\to\zo$ such that, with probability $1-\delta$, 
            \[
                \Pr_{\cDtrue}\inparen{{h}(x)\neq h^\star(x)}\leq \eps\,.
            \]
            The algorithm runs in time $\poly(n)$.
        \end{theorem}
        As mentioned in \cref{sec:intro}, \cref{thm:main} has several applications to learning from positive examples and truncated statistics, which we outline in \cref{sec:intro:application}. 
        The proof of \cref{thm:main} appears in \cref{sec:computationalEfficiency}.  
        
        To get a sense of the running time in \cref{thm:main}, see \cref{tab:gsa:GaussianSurfaceArea}, and for a comparison with prior works see \cref{tab:computational}.
        It is worth noting that the degree of polynomials $k$ required scales with $\nfrac{1}{\eps}$ and, typically, behaves as $k=\poly(\nfrac{1}{\eps})$.
        Therefore, the running time in \cref{thm:main} typically has an exponential dependence on $\nfrac{1}{\eps}$.
        This exponential dependence in the running time is necessary for any SQ-learning algorithm -- such as the one in \cref{thm:main} -- due to recent SQ lower bounds \cite{diakonikolas2024statistical} (\cref{rem:SQlowerbound}).
        \begin{remark}[SQ Lower Bound]\label{rem:SQlowerbound}
            \citet{diakonikolas2024statistical} provide an SQ-lower bound for the -- seemingly -- unrelated task of Gaussian parameter estimation with samples truncated to an unknown set.
            However, if we can avoid the exponential dependence on $\nfrac{1}{\eps}$ in \cref{thm:main} then, due to the recent results \cite{lee2024unknown}, we will also obtain a $\poly(\nfrac{d}{\eps})$ time algorithm for the problem studied by \citet{diakonikolas2024statistical} thereby contradicting their lower bound.
            See \cref{sec:application:unknownTruncation} for further discussion of parameter estimation from truncated samples and how our results relate to it.
        \end{remark}
        %
        \cref{thm:main} requires $\nL_1$-approximability of the target $h^\star$ with respect to $\cDgiven$ (which we have sample access to) and not with respect to $\cDtrue$ (which we do not have sample access to).
        In some cases, this enables one to use the sample access to $\cDgiven$ to check that the $\nL_1$-approximation requirement indeed holds (see \cref{rem:testing}).
    
        \begin{remark}[Testing Approximability]
            \label{rem:testing}
            \cref{thm:main} requires the hypothesis class $\hyH$ to be approximable by polynomials with respect to $\cDgiven$ in $\nL_1$-norm.
            Many guarantees on polynomial approximability are known with respect to specific distribution families such as Gaussians.
            One way to test this guarantee is to test whether the provided samples indeed come from a Gaussian distribution.
            However, most existing distributional testers have prohibitively large running time -- exponential in $d$ \cite{rubinfeld2023testing}.
            Instead, if $\hyH$ satisfies a slightly stronger property that it has sandwiching polynomials with respect to the underlying distribution, then one can use Theorem 3.2 of \citet{gollakota2023momentMatching} that gives a tester that runs in $\poly(d^{k},\nfrac{1}{\eps})$ time. %
            In particular, the tester comes with the following guarantees.
            \begin{enumerate}[itemsep=0pt]
                \item {(Completeness)}\quad It says \textsf{Yes} if samples are drawn from distribution $\cDgiven$;
                \item {(Soundness)}\quad If this tester says \textsf{Yes} on samples $S$, then with high probability, $\eps$-$\nL_1$-approximating polynomials exist with respect to the empirical distribution $\wh{\cD}$ over the samples $S$.\footnote{
                    To be precise, their result actually shows the existence of $\nL_1$-sandwiching polynomials, which is sufficient for us since $\nL_1$-sandwiching (whose definition is not relevant here) is a stronger requirement than $\nL_1$-approximation.
                }
            \end{enumerate}
            Now, we can run the algorithm in \cref{thm:main} with respect to the empirical distribution $\wh{\cD}$.
            To use \cref{thm:main}, we need to verify that $\cDtrue$ is smooth with respect to $\cDtrue$ in the sense of \cref{asmp:smoothness}.
            By standard VC theory for sufficiently large samples uniform convergence holds with respect to degree-$k$ polynomials, yielding the following additive-relaxation version of \cref{asmp:smoothness}:
            for any measurable set $T$, $\cDtrue(T)\leq (\nfrac{1}{\sigma}) \, \inparen{\smash{\wh{\cD}}(T)+\eps}^{1/q}$.
            This weaker guarantee is sufficient to execute our proof as we only use assumption \cref{asmp:smoothness} \mbox{for sets $T$ whose mass is bounded away from $0$.}
        \end{remark}

        
        \noindent A final remark is that our results (on both sample and computational complexity) can be adapted to work with more general notions of smoothness (see \cref{rem:generalSmoothness}).

        \paragraph{Optimal Sample Complexity for $q=1$.}
        Next, we improve the sample complexity in the important $q=1$ case if the following additional assumption holds. %
        \begin{restatable}[Density Lower Bound]{assumption}{lowerBound}\label{asmp:densityLB}
            There are known constants $\sigma\in (0,1]$, such that, for any measurable set $S$, $\cDtrue(S)\geq 
                \sigma
                \cdot \cDgiven(S).$
        \end{restatable}
        \vspace{-5mm}
        \begin{remark}\label{rem:densityLB}
            While \cref{asmp:densityLB} is similar to \cref{asmp:smoothness}, it is substantially stronger: 
            \eg{}, when the (imperfect) unlabeled samples are gathered from an unreliable source, they are bound to contain some elements not in support of the true unlabeled distribution $\cDtrue$, and such elements are enough to violate \cref{asmp:densityLB}.
        \end{remark} 
        With \cref{asmp:densityLB}, we obtain the following sample complexity.
        \begin{restatable}[Optimal Sample Complexity with \cref{asmp:densityLB}]{theorem}{optimalSampleComplexity}
            \label{thm:sampleComplexity:bothSided}
            Suppose \cref{asmp:smoothness,asmp:densityLB} hold with $q=1$. 
            Fix any $\eps,\delta\in (0,\nfrac{1}{2})$.
            There is an algorithm that, given $\eps,\delta$ and $n=\wt{O}\binparen{\inparen{\nfrac{1}{\eps\sigma}}^{-1}\cdot \inparen{\vc{}(\hyH)+\log{\nfrac{1}{\delta}}}}$ independent samples from $\cPtrue$ and $\cDgiven$, outputs a hypothesis $\wh{h}$ such that, with probability $1-\delta$,
            \[ 
                \Pr_{\cDtrue}\binparen{\wh{h}(x)\neq h^\star(x)}\leq \eps\,.
            \]
        \end{restatable}
        Note that this sample complexity is optimal for PIU learning (up to logarithmic factors), since it is optimal for realizable PAC learning which is an easier problem.

        
         
        

        
        \begin{table}[tbh!]
            \centering
            \vspace{4mm}
            \midsepremove{}
            \begin{tabular}{c | C{1.75cm}|C{1.75cm}|C{1.9cm}| C{1.75cm} | C{1.75cm} | C{1.75cm} }
                  & \multicolumn{6}{c}{Sample Complexity}\\
                 \cmidrule(r){2-7}
                  & \multicolumn{3}{c|}{Smoothness and density lower bound}
                  & \multicolumn{3}{c}{Smoothness (\cref{asmp:smoothness})}\\
                  & \multicolumn{3}{c|}{( \cref{asmp:smoothness,asmp:densityLB})} \\
                 \cmidrule(r){2-7}
                 & $q=1$ & $q=2$ & $q>2$ & $q=1$ & $q=2$ & $q>2$ \\
                 \midrule 
                 \cite{denis1990positive}
                 & $\eps^{-1}$ & --- & --- & --- & --- & --- \\
                 \cite{lee2024unknown}
                 & $\eps^{-4}$ & $\eps^{-6}$ &  $\eps^{-2(1+q)}$ & --- & --- & --- \\
                 This Work
                 & $\eps^{-1}$  & {{$\eps^{-4}$}}    & {$\eps^{-2q}$}    & $\eps^{-2}$ & $\eps^{-4}$  & $\eps^{-2q}$  \\
            \end{tabular}
            \midsepdefault{}
            \caption{\textit{Sample Complexity and Comparisons to Existing Work on Learning in the PIU Learning Model from \cref{thm:sampleComplexity,thm:sampleComplexity:bothSided}. Logarithmic dependence on $\eps$ and a linear factor of $\vc{(\hyH)}$ are hidden.}
            }
            \label{tab:sample}
        \end{table}
        




\addtocontents{toc}{\protect\setcounter{tocdepth}{2}}

\section{Proofs of Main Results for PIU Learning (\cref{thm:main,thm:sampleComplexity})}\label{sec:mainProof} 
    In this section, we prove our main results on the sample complexity (\cref{thm:sampleComplexity}) and computational complexity of PIU learning (\cref{thm:main}).

    
    The starting point of our approach is practical heuristics used for learning from positive and unlabeled samples.
    At a high level, these heuristics find a hypothesis $H$ containing the minimum number of unlabeled samples subject to selecting all positive samples \cite{bekker2020learning}.
    We formalize this idea as the following constrained empirical risk minimization (ERM) problem, which we call \textit{Pessimistic ERM:}
    given tolerance $\rho\geq 0$, sets of positive and unlabeled samples $P$ and $U$ respectively, and a hypothesis class $\hyH$, the \pERM{} problem is  
    \[
        \argmin_{H\in \hyH}~~ \frac{\abs{H\cap U}}{\abs{U}}
        \,,\quad\text{ such that}\,,\quad 
        \frac{\abs{H\cap P}}{\abs{P}} \geq 1-\rho
        \,.
        \tag{\pERM{}}
    \]


    \paragraph{Outline of this section.}
        First, in \cref{sec:perm}, we prove \cref{thm:sampleComplexity} using an algorithm that solves many instances of \pERM{} and aggregates the solutions.
        Second, in \cref{sec:efficientPERM}, we reduce computationally efficient PIU learning to developing a computationally efficient approximation algorithm for \pERM{}.
        This proves \cref{thm:main} up to designing an efficient approximation algorithm for \pERM{}, which we do in \cref{thm:constReg}.




    


    

    
    
    \subsection{Proof of \cref{thm:sampleComplexity} (Sample Efficient Algorithm)}
        \label{sec:perm}
        \label{sec:sampleEfficiency}
    In this section, we prove \cref{thm:sampleComplexity} which bounds the sample complexity of PIU learning.
    Our algorithm is an extension of the \pERM{} and solves $\nfrac{1}{\eps}$ carefully-constructed instances of \pERM{}; we believe this approach could be useful in practical settings where \pERM{} is used (\cref{sec:realWorldApplications}).
    

    
    Our proof relies on the following standard second-order uniform convergence bound, which gives sharper uniform convergence rates for hypotheses with small error.
    \begin{theorem}[Second-Order Uniform Convergence \protect{\cite[Section 5.1.2]{Boucheron_Bousquet_Lugosi_2005}}]
        \label{thm:unifConvergence}
        Let $\hyF$ be any concept class on $\R^d$ with finite VC-dimension, \ie{}, $\vc{(\hyF)}<\infty$, and let $\cD$ be any distribution over $\R^d$.
        For any confidence level $\delta\in (0,\nfrac{1}{2})$ and any fixed set $F^\star\subseteq \R^d$, with probability $1-\delta$, over the draw $X\sim \cD^n$, every $F\in \hyF$,
        \begin{align*}
            \cD\sinparen{F\triangle F^\star}
                &\leq 
                    D_n\sinparen{F\triangle F^\star}
                    + \sqrt{
                        D_n\sinparen{F\triangle F^\star}
                        \cdot 
                        s(\hyF,n,\delta)
                    }
                    + 
                    4 s(\hyF,n,\delta)\,,
                \yesnum\label{eq:sampleComplexity:1}\\
                D_n\sinparen{F\triangle F^\star}
                &\leq 
                    \cD\sinparen{F\triangle F^\star}
                    + \sqrt{
                        \cD\sinparen{F\triangle F^\star}
                        \cdot 
                        s(\hyF,n,\delta)
                    }
                    + 
                    4 s(\hyF,n,\delta)\,,
                \yesnum\label{eq:sampleComplexity:2}
        \end{align*}
        where $s(\hyF,n,\delta)$ and $D_n\sinparen{F\triangle F^\star}$ are defined as follows 
        \[
            s(\hyF,n,\delta) \coloneqq \frac{2\vc{(\hyF)}\log{(n+1)}+\log{\nfrac{4}{\delta}}}{n}
            \qquadand
            D_n\sinparen{F\triangle F^\star}
            \coloneqq 
            \frac{\abs{X\cap \sinparen{F\triangle F^\star}}}{\abs{X}}\,.
        \]
    \end{theorem}
    \begin{remark}[Necessity of Improper Learners for PIU Learning]
        \label{rem:necessityOfImproper}
    To illustrate that improper learners are necessary, consider the class of one-dimensional halfspaces (\ie{}, thresholds with a ``direction''). Define
        \[
        h_{\leq 1} = \{ x \in \R : x \leq 1 \} \qquadand h_{\geq -1} = \{ x \in \R : x \geq -1 \}\,.
        \]
        Let the two ``true" unlabeled distributions be 
        \[
            \cDtrue_1 = \unif[-2,1] \qquadand \cDtrue_{-1} = \unif[-1,2]\,.
        \]
        We now consider two PIU learning instances defined by (a): unlabeled distribution $\cDtrue_1$ with optimal hypothesis $h_{\leq 1}$ and (b): unlabeled distribution $\cDtrue_{-1}$ with optimal hypothesis $h_{\geq -1}$.
        Note that both instances yield the same positive sample distribution, $\cPtrue = \unif[-1,1]$.
        Let the imperfect unlabeled distribution for both instances be $\cDgiven = \unif[-2,2]$; which satisfy \cref{asmp:smoothness} with parameters $(\sigma,q)=(\nfrac{3}{4},1)$ in both instances.
        Since any deterministic learner cannot distinguish between these two instances, if it is proper (\ie{}, its output is restricted to one-dimensional halfspaces), it must incur a constant error on at least one of the instances. The same holds with probability $\nfrac{1}{2}$ for randomized learners.
    \end{remark}
    \subsubsection{Proof of \cref{thm:sampleComplexity}}
        
        \begin{proof}[Proof of \cref{thm:sampleComplexity}]
        The algorithm is presented in \cref{alg:sampleComplexity}.
        It solves $T\coloneqq\nfrac{1}{\eps}$ carefully selected instances of \pERM{} to obtain hypotheses $H_1,H_2,\dots,H_T$ and outputs $\bigcap_{i=1}^{T} H_i$.
        The instances of \pERM{} are constructed iteratively, where $i$-th problem removes all unlabeled samples not in $H_1 \cap H_2 \cap\dots \cap H_{i-1}.$
        \begin{algorithm}[htb!]
        \caption{\textsf{Iterative-\pERM{}} (from \cref{thm:sampleComplexity})}
        \begin{algorithmic}[1]
        \Procedure{\rm \textsf{Iterative-\pERM{}}}{$\eps,\delta,\sigma,q,\hyH$}
            \vspace{2mm}
            \State Set $n=\wt{O}\binparen{\inparen{\nfrac{1}{\eps\sigma}}^{2q}\cdot \inparen{\vc{}(\hyH)+\log{\nfrac{1}{\delta}}}}$
            \State Obtain $n$ positive samples $P$ and $n$ (imperfect) unlabeled samples $U$
            \vspace{2mm}
            \State Set $T=\nfrac{1}{\eps}$
            \vspace{2mm}
            \For{$1\leq i\leq T$}
                \State Define the survival set $S_i=H_1\cap H_2\cap\dots H_{i-1}$ if $i>1$ and, otherwise, $S_i=\R^d$
                \State Compute the solution $H_i$ of \pERM{} specified by tolerance $\rho=0$, positive
                \item[] \phantom{....}~~~~~~~ samples $P$, unlabeled samples $U\cap S_i$, and hypothesis class $\hyH$
            \EndFor{}
            \vspace{2mm}
            \State \textbf{return} $H\coloneqq H_1\cap H_2\cap\dots\cap H_T.$
        \EndProcedure
        \end{algorithmic}
        \label{alg:sampleComplexity}
        \end{algorithm} 
        \noindent We divide the proof of correctness into two parts.
        
        \paragraph{Step A (Bounding the number of false negatives):}
            Each of $H_1,H_2,\dots,H_T$ is feasible for the \pERM{} constraint with tolerance $\rho=0$ and positive samples $P$ with $\abs{P}=n$.
            We claim that this implies that that with probability $1-\delta$, for each $1\leq i\leq k$, $\cDtrue\inparen{\optset\setminus H_i}\leq O\inparen{\eps^{2q}}\leq O(\eps^2)$ (as $q\geq 1$).
            We proved this claim below and for now proceed with the proof.
            Since $H\coloneqq H_1\cap H_2\cap\dots\cap H_T$, an immediate consequence of the above is that with probability $1-\delta$, 
            \[
                \cDtrue\inparen{\optset\setminus H}
                \leq \sum_i \cDtrue\inparen{\optset\setminus H_i}
                \leq O\inparen{\frac{\eps^2}{\eps}}
                = O(\eps)\,.
            \]

            \noindent \textit{Proof of Claim.}~~
            Next, we prove our earlier claim.
            Let $S\in\hyH$ be any solution of \pERM{} with $\rho=0$.
            By feasibility, it satisfies $\nfrac{\abs{S \cap P}}{\abs{P}}=1$ and, hence, $\nfrac{\abs{ P\setminus S}}{\abs{P}}=0$.
            Due to the choice of the sample size $n$ and \cref{thm:unifConvergence}, with probability $1-\inparen{\nfrac{\delta}{2}}$, for each $R\in \hyH\cup \inbrace{G^c\colon G\in \hyH}$,\footnote{Here, we use the facts that $\vc{(\inbrace{G^c\colon G\in \hyH})}=\vc{(\hyH)}$ and $\vc{(\hyH_1\cup\hyH_2)}\leq {\vc{(\hyH_1)}+\vc{(\hyH_2)}+1}$ \cite{mohri2018foundations}.}
            \[
                \cPtrue(R) 
                \leq 
                \frac{\abs{R\cap P}}{\abs{P}}
                +
                \sqrt{\frac{\abs{R\cap P}}{\abs{P}}\cdot \eps}
                + \eps\,.
            \]
            In particular, applying this to $R=S^c$ (note that $S^c\in\{G^c: G\in \hyH\}$), and using $\nfrac{\abs{ P\cap S^c}}{\abs{P}}=\nfrac{\abs{ P\setminus S}}{\abs{P}}=0$, implies that 
            \[
                \cPtrue(S^c) \leq \eps\,.
            \]
            Finally, since $\cDtrue( R )\leq \cPtrue( R )$ for any $ R \subseteq \optset$, it follows that 
            \[
                \cDtrue(\optset\setminus  S)
                \leq \cPtrue(\optset\setminus  S )
                = \cPtrue(S^c)
                \leq  \eps\,.
            \]  
        
        \paragraph{Step B (Bounding the number of false positives):}
            Next, we bound the amount of false positives.
            We claim that 
            \[
                \abs{U\setminus H_1}
                \geq 
                \abs{(U\cap H_1)\setminus H_2}
                \geq 
                \abs{(U\cap H_1\cap H_2)\setminus H_3}
                \geq 
                \dots 
                \geq 
                \abs{
                    (U\cap H_1\cap \dots \cap H_{T-1})
                    \setminus H_{T}
                }\,.
                \yesnum\label{eq:oneSided:monotone}
            \]
            To prove this, consider any $1\leq i\leq k-1$ and toward a contradiction suppose that 
            \[
                \abs{(U\cap H_1\cap \dots \cap H_{i-1})\setminus H_i}
                <
                \abs{(U\cap H_1\cap \dots \cap H_{i})\setminus H_{i+1}}\,.
                \yesnum\label{eq:onesided:algo:contradiction}
            \]
            Hence, in particular,
            \begin{align*}
                \abs{(U\cap H_1\cap \dots \cap H_{i-1})\setminus H_{i+1}}
                ~\geq~ \abs{(U\cap H_1\cap \dots \cap H_{i})\setminus H_{i+1}}
                ~\Stackrel{\eqref{eq:onesided:algo:contradiction}}{>}~ \abs{(U\cap H_1\cap \dots \cap H_{i-1})\setminus H_i}\,.
            \end{align*}
            This contradicts the fact $H_i$ is the optimal solution of \pERM{} defined by tolerance $\rho=0$, positive samples $P$, and unlabeled samples $U\cap H_1\cap \dots \cap H_{i-1}$.
            Therefore, due to the contradiction, our supposition in \cref{eq:onesided:algo:contradiction} must be wrong and, hence, \cref{eq:oneSided:monotone} is true.

            Next, we will use the above claim to conclude that $\abs{
                    (U\cap H_1\cap H_2\cap \dots \cap H_{T-1})
                    \setminus H_{T}
                }
                \leq \eps\abs{U}$.
            Toward this, observe that the following collection of sets
            \[
                \inbrace{(U\cap H_1\cap H_2\cap \dots \cap H_{j})
                    \setminus H_{j+1} \mid 1\leq j\leq T}\,,
            \]
            is disjoint and, hence, $\sum_i \abs{(U\cap H_1\cap H_2\cap \dots \cap H_{i})}\leq \abs{U}$.
            This, along with \cref{eq:oneSided:monotone} implies that 
            \[
                \abs{
                    (U\cap H_1\cap H_2\cap \dots \cap H_{T-1})
                    \setminus H_{T}
                }
                ~~\leq~~ \frac{\abs{U}}{T}
                ~~~\Stackrel{(T=\nfrac{1}{\eps})}{=}~~~ \eps\abs{U}\,.
                \yesnum\label{eq:oneSided:upperBound}
            \]
            \noindent Now, we are ready to prove the claim.
                    Consider the $T$-th instance of \pERM{} constructed in the above algorithm, \ie{}, the one defined by unlabeled samples $U\cap H_1\cap H_2\cap \dots \cap H_{T-1}$.
                    Since $\optset\supseteq P$, $\optset$ is feasible for this instance.
                    But since $H_{T}$ is the optimal solution of instance {$T$},
                    \[
                        \abs{U\cap H_1\cap H_2\cap \dots \cap H_{T-1} \cap \optset}
                        \geq 
                        \abs{U\cap H_1\cap H_2\cap \dots \cap H_{T-1} \cap H_{T}}\,.
                    \]
                    Taking complements implies 
                    \[
                        \abs{U\cap H_1\cap H_2\cap \dots \cap H_{T-1} \setminus \optset}
                        \leq 
                        \abs{U\cap H_1\cap H_2\cap \dots \cap H_{T-1} \setminus H_{T}}
                        ~~\Stackrel{\eqref{eq:oneSided:upperBound}}{\leq}~~ 
                            \eps \abs{U}\,.
                    \]
                    Since $H\subseteq H_1\cap H_2\cap \dots \cap H_{T}$, the above in particular implies that
                    \[
                        \abs{U\cap H \setminus \optset}
                        \leq 
                        \abs{U\cap H_1\cap H_2\cap \dots \cap H_{T-1} \setminus \optset}
                        \leq \eps\abs{U}\,.
    \yesnum\label{eq:oneSided:empiricalUB}
                    \]
                    Now, second-order uniform convergence (\cref{thm:unifConvergence}) with respect to $U$ and the following hypothesis class 
                    \[
                        \hyH_{\cap T} = \inbrace{G_1\cap G_2\cap \dots\cap G_{T}\setminus\optset\colon G_1,G_2,\dots,G_{T} \in \hyH}\,,
                    \]
                    implies that 
                    \[
                        \cD(H\setminus \optset) \leq \eps + \wt{O}\inparen{\sqrt{
                            \frac{\abs{U\cap H\setminus \optset}}{\abs{U}}\cdot 
                            \frac{\vc{}\inparen{\hyH_{\cap T}}+\log{\nfrac{1}{\delta}}}{n}}}
                        + \wt{O}\inparen{{\frac{\vc{}\inparen{\hyH_{\cap T}}+\log{\nfrac{1}{\delta}}}{n}}}\,.
                    \]
                    Substituting the value of $n$, using that $\vc{}\inparen{\hyH_{\cap T}}=\vc{{(\hyH)}}\cdot O(T\log{T})$, and $\nfrac{\abs{U\cap H\setminus\optset}}{\abs{U}}\leq \eps$ (from \cref{eq:oneSided:empiricalUB}) implies that 
                    \[
                        \cD(H\setminus \optset) \leq \eps + \wt{O}\inparen{\sqrt{
                            \eps\cdot \frac{\vc{}(\hyH) T +\log{\nfrac{1}{\delta}}}{n}}
                        }
                        + \wt{O}\inparen{
                            \frac{\vc{}(\hyH) T+\log{\nfrac{1}{\delta}}}{n}
                        }
                        \leq O(\eps)
                        \,.
                    \]
                    Therefore, \cref{asmp:smoothness}, implies 
                    \[
                        \cDtrue(H\setminus \optset)
                        \leq O(\nfrac{1}{\sigma})\cdot \eps^{1/q}\,,
                    \]
                    completing the proof.

    \end{proof}
    
    \subsection{Proof of \cref{thm:main} (Computationally Efficient Algorithm)}
        \label{sec:computationalEfficiency}
        In this section, we prove \cref{thm:main}.
        The key step in the proof is to design the following pseudo-approximation algorithm for \pERM{}.
        \begin{restatable}[]{theorem}{constReg}
            \label{thm:constReg}
            Suppose \cref{asmp:smoothness,asmp:posFraction} hold.
            Fix any $\delta\in (0,\nfrac{1}{2})$, $k$ such that degree-$k$ polynomials $\zeta$-approximate $\hyH$ with respect to $\cDgiven$ in $\nL_1$-norm for $0<\zeta\leq O\sinparen{{\alpha\sigma}{}}^{2q}$, and $n=\poly\inparen{d^k, \nfrac{1}{\zeta}, \log{\nfrac{1}{\delta}}}$.
            There exists a polynomial time algorithm (\cref{alg:constReg}) that given $n$ independent samples from the PIU learning model (\cref{prob:PUlearning}), outputs a degree-$k$ PTF $H$, s.t., the following hold with probability $1-\delta$
            \begin{enumerate}
                \item $\cDtrue\sinparen{\optset \setminus H}\leq O\sinparen{\nfrac{1}{\alpha\sigma}} \cdot \zeta^{\sfrac{1}{(4q)}}$; and 
                \item $\cD(H)\leq \opt_{\rm PERM} +
                    O\sinparen{\nfrac{1}{\alpha\sigma}}\cdot {\zeta^{\sfrac{1}{(4q)}}}+O(\delta)$. 
            \end{enumerate}
            \mbox{Where $\opt_{\rm PERM}\coloneqq \min_{H\in \hyH \,:\, \cPtrue(H)=1} \cD(H)$ is the optimal value of \pERM{} with tolerance $\rho=0$.}\footnote{Note that \pERM{} is defined with respect to samples and, hence, is a random variable. $\opt_{\rm PERM}(\rho)$ is the limiting optimal value as the number of samples goes to infinity.}
        \end{restatable}
        In the remainder of this section, we prove \cref{thm:main} assuming \cref{thm:constReg}, and then present the proof of \cref{thm:constReg} in the next section.

        \paragraph{Outline of Proof of \cref{thm:main} Assuming \cref{thm:constReg}.}
        The proof of \cref{thm:constReg} follows the strategy of the sample complexity proof (\cref{sec:perm}). 
        We iteratively construct instances of \pERM{} and define the final output as the intersection of the (approximate) solutions obtained at each iteration. 
        Compared to the sample complexity setting, several additional challenges arise since here we only find an approximate solution of the \pERM{} instances. 
        In particular, we must ensure that:
        \begin{enumerate}
            \item The hypothesis class $\hyH$ remains well-approximated by low-degree polynomials with respect to the updated unlabeled distributions.
            \item The smoothness condition in \cref{asmp:smoothness} holds at every iteration.
        \end{enumerate}
        To address the former challenge, we include a check that the remaining region has at least a constant fraction of the mass. 
        To address the latter challenge, we also truncate the distribution of positive samples in addition to the distribution of unlabeled samples, and account for this in the analysis.
        Moreover, we truncate the distribution of positive samples and incorporate this modification into our analysis, ensuring that the surviving (\ie{}, non-truncated) region has sufficient mass to bound the runtime of obtaining samples from the truncated distributions via rejection sampling.
 

    \begin{figure}[h!]
        \centering
        \begin{tikzpicture}[node distance=2.5cm, scale=0.83, transform shape]
        \node[myNode, line width=2pt, text width=9cm] (correctness) at (0,-2) 
            {{Computationally Efficient Algorithm (\cref{alg:main}) \\[1mm] \cref{thm:main}}};


         \node[myNode, line width=1pt, text width=4cm](claim45) at (-7.5, -4.2)
            {{Guarantee on the solution of each \pERM{} instance \\[1mm] \cref{lem:main:chaining}}};

        \node[myNode, line width=2pt, text width=7cm] (thmConstReg) at (0, -4.2)
            {{Efficient Approximation Algorithm for \pERM{} (\cref{alg:constReg})\\[1mm] \cref{thm:constReg}}};

        \node[myNode, line width=1pt, text width=2.2cm](lemRejSamp) at (7.5, -4.2)
            {\cref{lem:main:rejectionSampling}};


        \node[myNode, line width=1pt, text width=4cm](lemNew) at (7.5, -6.2)
            {{{Efficient sampling from created instances}}};


        \draw[-stealth, line width=0.5mm](lemNew) -- (lemRejSamp);
        

        \node[myNodeFlex, line width=1pt, text width=4.5cm](lem44) at (-7.5, -7)
            {{Assumptions required by \cref{thm:constReg} hold in each iteration (\cref{lem:main:assumptions})}};

        \node[myNodeFlex, line width=1pt, text width=2.25cm](r2) at (-3.5, -6.5)
            {{Feasibility \\ [1mm] \cref{result:constReg:feasibility}}};

        \node[myNodeFlex, line width=1pt, text width=2.25cm](r1) at (0, -6.5)
        {{{Runtime} \\[1mm] \cref{result:constReg:polyTimeAlg}}};

        \node[myNodeFlex, line width=1pt, text width=2.25cm](r3) at (3.5, -6.5)
            {{Optimality \\ [1mm] \cref{result:constReg:optimality}}};

        \draw[-stealth, line width=0.5mm] (claim45) -- (correctness); 
        \draw[-stealth, line width=0.5mm] (lem44) -- (claim45); 
        \draw[-stealth, line width=0.5mm] (thmConstReg) -- (claim45); 
        \draw[-stealth, line width=0.5mm] (thmConstReg) -- (lemRejSamp); 
        \draw[-stealth, line width=0.5mm] (thmConstReg) -- (correctness); 
        \draw[-stealth, line width=0.5mm] (lemRejSamp) -- (correctness); 
        

        \draw[-stealth, line width=0.5mm](r1) -- (thmConstReg);
        \draw[-stealth, line width=0.5mm](r2) -- (thmConstReg);
        \draw[-stealth, line width=0.5mm](r3) -- (thmConstReg);

        \node[myNode, line width=1pt, text width=4cm] (lem48) at (-3.5, -9)
            {{Repetitions are feasible \\[1mm] \cref{lem:constReg:feasiblity}}};

        \node[myNodeFlex, line width=1pt, text width=4.25cm] (lem411) at (1, -9)
            {{Boosting via repetition \\ [1mm] \cref{lem:constReg:boost}}};
        \node[myNodeFlex, line width=1pt, text width=4.5cm] (lem410) at (6, -9)
            {{\mbox{Optimal's expected value} \\ [1mm] \cref{lem:constReg:expectedValue}}};

        \node[myNode, line width=1pt, text width=3.75cm](lem47) at (-5.5, -11.2)
            {{Repetitions satisfy uniform convergence \\ [1mm] \cref{lem:constReg:unifConvergenceHoldsWHP}}};
        \node[myNode, line width=1pt, text width=3.5cm](lem49) at (-1.5, -11.2)
            {{Repetitions feasible over samples \\ [1mm] \cref{lem:constReg:feasibility:unconditional}}};

        \draw[-stealth, line width=0.5mm](lem48) -- (r2);
        \draw[-stealth, line width=0.5mm](lem47) -- (lem48);
        \draw[-stealth, line width=0.5mm](lem49) -- (lem48);

        \node[myNodeNarrow, line width=1pt, text width=2.2cm](lem412) at (3.25, -11.2)
            {{\cref{lem:constReg:expectedValue:upperBound}}};
        \node[myNodeNarrow, line width=1pt, text width=2.2cm](lem413) at (6, -11.2)
            {{\cref{lem:constReg:expectedValue:feasibility}}};
        \node[myNodeNarrow, line width=1pt, text width=2.2cm](lem414) at (8.75, -11.2)
            {{\cref{lem:constReg:expectedValue:ub2}}};

        \draw[-stealth, line width=0.5mm](lem411) -- (r3);
        \draw[-stealth, line width=0.5mm](lem410) -- (r3);
        \draw[-stealth, line width=0.5mm](lem412) -- (lem410);
        \draw[-stealth, line width=0.5mm](lem413) -- (lem410);
        \draw[-stealth, line width=0.5mm](lem414) -- (lem410);
        
        \end{tikzpicture}
        \caption{\textit{Outline of the proof of \cref{thm:main}.} 
        The proof of \cref{thm:main} follows a structure analogous to the proof of the sample complexity of PIU learning (\cref{thm:sampleComplexity}): 
        we construct a sequence of instances of \pERM{} and output the intersection of all the (approximate) solutions $H_1,H_2,\dots$ to the instances created.
        To approximately solve each instance created, we use an approximation algorithm for \pERM{} (constructed in \cref{thm:constReg}).
        This is the main new technical ingredient in the proof of \cref{thm:main} compared to \cref{thm:sampleComplexity}.
        }
        \label{fig:outline:compComplexity}
    \end{figure}



    \paragraph{Proof of \cref{thm:main}.}
        We use the algorithm in \cref{alg:main}.
        First, we rescale $\eps$ and $\delta$ to ensure that $\eps\leq \nfrac{\alpha}{6}$ and $\delta\leq \sigma^2\eps^{2q}$. %
        Due to the scaling in $\eps$, the relation between $\zeta$ and $\eps$ becomes $\zeta\leq O\binparen{(\sigma\eps)^{5q^2}\alpha^{8q}}$ where $\eps$ is the rescaled accuracy parameter, and we will use this updated relation in our analysis.
        Due to the scaling in $\delta$, the running time increases by logarithmic factors in $\max\inbrace{\nfrac{1}{\sigma},\nfrac{1}{\eps^q}}$.
        The following constant shows up in our analysis:
        \[
            \gamma \coloneqq \sigma \eps^q\,.
        \]
        \begin{algorithm}[htb!]
        \caption{\textsf{Iterative-}\constrainedreg{} (from \cref{thm:main})}
        \begin{algorithmic}[1]
        \Procedure{\rm \textsf{Iterative-}\constrainedreg{}}{$\eps,\delta,\sigma,q,\alpha,k$}
            \vspace{2mm}
            \State {If $\eps \geq \alpha$, set $\eps=\nfrac{\alpha}{6}$}
            \State Set the accuracy parameters {$\zeta=O((\sigma\eps)^{5q^2}\alpha^{8q})$} and $\gamma=\sigma^{-1}\eps^{-q}$
            \State Set the number of iterations as $T=\nfrac{1}{\gamma}$
            \vspace{2mm}
            \vspace{2mm}
            \For{$1\leq i\leq T$}
                \State Define the survival set $S_i=H_1\cap H_2\cap\dots H_{i-1}$ if $i>1$ and, otherwise, $S_i=\R^d$
                \vspace{3mm}
                \item[] \phantom{.}\qquad~~ \textit{\#~~Step 1 (Ensure Low-Degree Approximation Exists)}
                \State Obtain $m=\wt{O}((\nfrac{1}{\gamma^3}) \cdot (d^k+\log{\nfrac{1}{\delta}}))$ samples $U$ from $\cDgiven$ and $P$ from $\cPtrue$
                \If{$i > 1$ and either $\abs{U\cap S_i}\leq \gamma \abs{U}$ or $\abs{P\cap S_i}\leq (\nfrac{\alpha}{2})\abs{P}$}\label{step:alg:main:massCheck}
                    \State Break out of the for-loop
                \EndIf{}
                \vspace{3mm}
                \item[] \phantom{.}\qquad~~ \textit{\#~~Step 2 (Update Sample Distributions)}
                \State Define $\cDgiven_i,\cPtrue_i,\cDtrue_i$ as the truncations of $\cDgiven,\cPtrue,\cDtrue$ to $S$
                \State Construct sampling oracles for $\cDgiven_i$ and $\cPtrue_i$ via rejection sampling from $\cDgiven$ and $\cPtrue$
                \vspace{3mm}
                \item[] \phantom{.}\qquad~~ \textit{\#~~Step 3 (Solve $i$-th \pERM{} Instance)}
                \State $H_i~\longleftarrow$~\hyperref[alg:constReg]{{\rm \textsf{Boosted-}}\constrainedreg{}}($\zeta,\delta,\sigma,q,k$) with
                sample access to $\cD_i$ and $\cPtrue_i$
                \item[] \phantom{.}\qquad~~ \textit{\#~~Here, \hyperref[alg:constReg]{{\rm \textsf{Boosted-}}\constrainedreg{}} is the algorithm referenced in \cref{thm:constReg}}
            \EndFor{}
            \vspace{2mm}
            \State \textbf{return} $H\coloneqq H_1\cap H_2\cap\dots\cap H_\ell$ where $\ell$ is the last iteration of the for-loop completed
        \EndProcedure
        \end{algorithmic}
        \label{alg:main}
        \end{algorithm} 
        
        \paragraph{Assumptions Required by \cref{thm:constReg} Hold.}
            We begin by proving that the assumptions required by \cref{thm:constReg} hold in each iteration where the for-loop executes.
            \begin{lemma}
                \label{lem:main:assumptions}
                In each iteration $i$ where the for-loop in \cref{alg:main} executes, the following hold:
                \begin{enumerate}
                    \item Degree-$k$ polynomials $(\nfrac{\zeta}{\gamma})$-approximate $\hyH$ with respect to  $\cDgiven_i$ in $\nL_1$-norm. %
                    \item $(\cDtrue_i,\cPtrue_i,\cDgiven)$ is a valid PIU learning instance with constants $\sigma'\geq \nfrac{\alpha^2 \sigma}{3}$ and $q'=q$.
                \end{enumerate}
            \end{lemma}
            \begin{proof}[Proof of \cref{lem:main:assumptions}]
                Due to the guarantee of \cref{thm:constReg}, each of $H_1,H_2,\dots,H_T$ is a degree-$k$ PTF over $d$ variables.
                Since $S_i$ is an intersection of at most $T=\nfrac{1}{\gamma}$ sets from $H_1,H_2,\dots,H_T$, $S_i$ belongs to the family of the intersections of (at most) $\nfrac{1}{\gamma}$ degree-$k$ PTFs over $d$-variables, which has a VC dimension at most $\wt{O}(d^k/\gamma)$.
                Now, since the for-loop executes in the $i$-th iteration, the check in Step~\ref{step:alg:main:massCheck} ensures that with probability $1-\delta$, 
                \[
                    \cD(S_i)\geq \Omega({\gamma})
                    \quadand
                    \cPtrue(S_i)\geq \frac{\alpha}{2}-\gamma \geq \frac{\alpha}{3}\,,
                    \yesnum\label{eq:main:massGuarantee}
                \]
                due to uniform convergence over the class of (at most) $\nfrac{1}{\gamma}$ degree-$k$ PTFs over $d$-variables and $\gamma\leq \eps \leq \nfrac{\alpha}{6}$.
                Further, since degree-$k$ polynomials $\zeta$-approximate $\hyH$ with respect to  $\cDgiven$ in $\nL_1$-norm, for every $H'\in \hyH$, there is a degree-$k$ polynomial $p'\colon \R^d\to \R$ such that 
                \[
                    \Ex_{\cDgiven}\insquare{\abs{p'(x) - \mathds{1}\sinbrace{x\in H'}}}\leq \zeta\,.
                \]
                Since $\cD(S_i)\geq \Omega(\gamma)$ and $\cDgiven_i$ is the truncation of $\cD$ to $S_i$, a change of measure, implies 
                \[
                    \Ex_{\cDgiven_i}\insquare{\abs{p'(x) - \mathds{1}\sinbrace{x\in H'}}}\leq O\inparen{\frac{\zeta}{\gamma}}\,.
                \]
                Since this holds for each $H'\in \hyH$, the first result follows.

                To prove the second claim, it suffices to show that $\cPtrue_i$ is the truncation of $\cDtrue_i$ to the set $\optset$ (which holds by construction) and that $\cDgiven_i$ and $\cDtrue_i$ satisfy \cref{asmp:smoothness} with the specified constants.
                To prove the latter, consider any measurable set $T\subseteq \R^d$.
                Observe that due to \cref{asmp:smoothness}
                \[
                    \cDtrue_i(T)
                    = \frac{\cDtrue(T\cap S_i)}{\cDtrue(S_i)}
                    \leq \frac{\cD(T\cap S_i)^{1/q}}{\sigma\cdot \cDtrue(S_i)}
                    \leq \frac{3\cdot \cD(T\cap S_i)^{1/q}}{\sigma \alpha^2}
                    \,.
                \]
                Where in the last inequality we used \cref{eq:main:massGuarantee} and the fact that $\cDtrue(S_i)\geq \cDtrue(\optset)\cPtrue(S_i)\geq \alpha\cPtrue(S_i)$ (due to \cref{asmp:posFraction}).
                We can continue the above chain of inequalities as follows 
                \[
                    \cDtrue_i(T)
                    \leq \frac{3}{\sigma\alpha^2} \cdot 
                        \inparen{\frac{\cD(T\cap S_i)}{\cD(S_i)}}^{1/q}
                    = \frac{3}{\sigma\alpha^2}\cdot \cDgiven_i(T)^{1/q}\,.
                \]
            \end{proof}
        
        \paragraph{Guarantee of \cref{thm:constReg}.}
        %
        To show correctness, we have to show that with probability $1-\delta$, $\cDtrue(\optset\triangle H)\leq \eps$. 
        Since the required assumptions hold (\cref{lem:main:assumptions}), for each iteration $1\leq i\leq T$ executed, \cref{thm:constReg} implies the following guarantees:
        \begin{align*}
            \cDtrue_i(\optset\setminus H_i) 
                \leq O\inparen{\frac{1}{\sigma\alpha^2}}\cdot \inparen{\frac{\zeta}{\gamma}}^{1/(4q)}
                \quadand
            \cDgiven_i(H_i) 
                &\leq \opt_{\rm PERM}(i) + O\inparen{\frac{1}{\sigma\alpha^2}}\cdot \inparen{\frac{\zeta}{\gamma}}^{1/(4q)} + O(\delta)\,.
        \end{align*}
        Where $\opt_{\rm PERM}(i)\coloneqq \min_{H\in \hyH\,\colon\,\cPtrue_i(H)=1}\cDgiven_i(H)$ is the optimal value of the $i$-th instance of \pERM{} in \cref{alg:main}.
        Since we ensured that $\zeta\leq O((\sigma\eps)^{5q^2}\alpha^{8q})\leq  \gamma^{1+4q}\cdot (\sigma\alpha^2)^{4q}$ and $\delta\leq \gamma^2$, the above simplify to 
        \begin{align*}
            \cDtrue_i(\optset\setminus H_i) 
                \leq \gamma^2
                \quadand
            \cDgiven_i(H_i) 
                &\leq \opt_{\rm PERM}(i) + \gamma^2\,.
                \yesnum\label{eq:main:constRegGuarantee}
        \end{align*} 
        in the remainder of the proof, we condition on the event that the above inequalities hold. 

        \paragraph{Correctness.}
        Now, we are ready to complete the proof of correctness of \cref{alg:main}.
            We consider two cases depending on whether the for-loop exits.

            \itparagraph{Case A (For-loop breaks at iteration $i$):}
            First, suppose that the if statement first broke at step $i$, then the check in Step~\ref{step:alg:main:massCheck} (along with uniform convergence; as discussed in \cref{lem:main:assumptions}) implies that $\cD(S_i)=\cD(H_1\cap H_2\cap\dots H_{i-1})\leq O(\gamma)$ and $H=H_1\cap H_2\cap\dots H_{i-1}$.
            Therefore, $\cD(H\setminus \optset)\leq \cD(H)\leq O(\gamma)$. 
            Applying \cref{asmp:smoothness} yields $\cDtrue(H \setminus \optset) \leq \sigma^{-1}\cdot \gamma^{1/q}\leq \eps$.
            Further, since all previous iterations $1,2,\dots,i-1$, the for-loop executed, \cref{eq:main:constRegGuarantee} holds and implies that 
            $\cDtrue(\optset\setminus H)\leq \sum_{j=1}^{i-1} \cDtrue(\optset\setminus H_i)\leq (i-1)\gamma^2\leq \gamma\leq \eps$ where we used that $i\leq T=\nfrac{1}{\gamma}$.
            Combining these bounds implies that 
            \[
                \cDtrue(\optset\triangle H)\leq 2\eps\,,
            \]
            completing the proof of correctness in this case.

            \itparagraph{Case B (For-loop never breaks):}
            Now, we consider the other case, where the break condition is never satisfied and, hence, all $T$ iterations of the for loop take place.
            \begin{align*}
                \cDtrue(\optset\setminus H)~~
                &\leq~~ \sum_i \cDtrue\sinparen{\optset\cap H_1 \cap H_2 \cap \dots \cap H_{i-1}\setminus H_i}\\
                &\leq~~ \sum_i \cDtrue\sinparen{\optset\setminus H_i}\\
                &\Stackrel{\eqref{eq:main:constRegGuarantee}}{\leq}~~ 
                    \gamma\,. 
                    \tagnum{since $T=\nfrac{1}{\gamma}$}{eq:computational:caseB}
            \end{align*}
            To finish the proof of correctness in Case B, we need to bound $\cDtrue(H \setminus \optset)$. 
            For this, we will use the following claim to bound $\cD(H \setminus \optset)$, and then apply \cref{asmp:smoothness}.
            \begin{claim}
                \label{lem:main:chaining}
                It holds that 
                \begin{align*}
                    \cDgiven(\R^d\setminus H_1)
                    &\geq \cD(H_1\setminus H_2) - \gamma^2\\
                    &\geq \cD(H_1\cap H_2\setminus H_3) - 2\gamma^2\\
                    &~~\vdots\\
                    &\geq \cD(H_1\cap H_2\cap \dots H_{T-1}\setminus H_T) - (T-1)\gamma^2\,.
                \end{align*}
            \end{claim}
            \begin{proof}[Proof of \cref{lem:main:chaining}]
                Toward a contradiction, suppose that for some index $i$
                \[
                    \cD(H_1\cap H_2\cap \dots \cap H_i\setminus H_{i+1})
                    > \cD(H_1\cap H_2\cap \dots \cap H_{i-1}\setminus H_{i}) + \gamma^2\,.
                \]
                Replacing the intersection on the LHS from $H_1\cap H_2\cap \dots \cap H_i$ to $H_1\cap H_2\cap \dots \cap H_{i-1}$, implies
                \[
                    \cD(H_1\cap H_2\cap \dots \cap H_{i-1}\setminus H_{i+1})
                    > \cD(H_1\cap H_2\cap \dots \cap H_{i-1}\setminus H_{i}) + \gamma^2\,.
                \]
                Equivalently, 
                \[
                    \cD_{H_1\cap H_2\cap \dots \cap H_{i-1}}(\R^d\setminus H_{i+1})
                    > \cD_{H_1\cap H_2\cap \dots \cap H_{i-1}}(\R^d\setminus H_{i}) 
                        + \frac{\gamma^2}{\cD(H_1\cap H_2\cap \dots \cap H_{i-1})}\,.
                \]
                Since $\cD(H_1\cap H_2\cap \dots \cap H_{i-1})\geq \Omega(\gamma)$ due to \cref{lem:main:assumptions}, the above implies that 
                \[
                    \cD_{H_1\cap H_2\cap \dots \cap H_{i-1}}(\R^d\setminus H_{i+1})
                    > \cD_{H_1\cap H_2\cap \dots \cap H_{i-1}}(\R^d\setminus H_{i}) 
                        + O(\gamma)\,.
                \]
                Taking complements, we get 
                \[
                    \cD_{H_1\cap H_2\cap \dots \cap H_{i-1}}(H_{i}) 
                    > \cD_{H_1\cap H_2\cap \dots \cap H_{i-1}}(H_{i+1})
                        + O(\gamma)\,.
                \]
                This contradicts \cref{eq:main:constRegGuarantee} (at the $i$-th step) since $H_{i+1}$ is feasible for the $i$-th instance of \pERM{} constructed by \cref{alg:main} and, hence, $\cD_{H_1\cap H_2\cap \dots \cap H_{i-1}}(H_{i+1})\geq \opt_{\rm PERM}(i)$.
            \end{proof}
            Now we will use \cref{lem:main:chaining} to bound $\cD(H\setminus\optset)$. 
            Since the sets $\inbrace{H_1\cap H_2\cap\dots \cap H_{i-1}\setminus H_i}_i$ are disjoint, it follows that 
            \[
                T\cdot \cD(H_1\cap H_2\cap \dots \cap H_{T-1} \setminus H_T)
                - T(T-1)\gamma^2
                \leq \sum_i \cDgiven(H_1 \cap H_2 \cap \dots \cap H_{i-1}\setminus H_i)
                \leq 1\,.
            \]
            Therefore, 
            \begin{align*}
                \cD(H_1\cap H_2\cap \dots \cap H_{T-1}\setminus H_T)
                \leq \frac{1}{T} + (T-1)\gamma^2
                \leq 2\gamma\,.
                \yesnum\label{eq:main:chainingGuarantee}
            \end{align*}%
            Apart from this, as $\cD_{H_1\cap H_2\cap \dots \cap H_{T-1}}(H_T)\leq \opt_{\rm PERM}(T)+\gamma^2$ and $\optset$ is feasible for the $T$-th instance of \pERM{} (so $\cD_{H_1\cap H_2\cap \dots \cap H_{T-1}}(\optset)\geq \opt_{\rm PERM}(T)$), the following also holds
            \[
                \cD_{H_1\cap H_2\cap \dots \cap H_{T-1}}(H_T)
                \leq 
                \cD_{H_1\cap H_2\cap \dots \cap H_{T-1}}(\optset)
                + \gamma^2\,.
            \]
            and consequently, 
            \[
                \cD(H_1\cap H_2\cap \dots \cap H_{T-1}\setminus \optset)
                \leq 
                \cD(H_1\cap H_2\cap \dots \cap H_{T-1}\setminus H_T)
                + \gamma^2\,.
                \yesnum\label{eq:main:optimalityGuaranteeTth}
            \]
            Hence,
            \[
                \cD(H \setminus \optset)
                \leq 
                \cD(H_1\cap H_2\cap \dots \cap H_{T-1}\setminus \optset)
                ~~~\Stackrel{\eqref{eq:main:chainingGuarantee},~\eqref{eq:main:optimalityGuaranteeTth}}{\leq}~~~  
                \gamma^2 + 2\gamma
                \leq 3\gamma\,.
            \]
        Applying \cref{asmp:smoothness}, we have $\cDtrue(H \setminus \optset) \leq \sigma^{-1} \cdot \cD(H \setminus \optset)^{1/q} \leq \sigma^{-1} \cdot (3\gamma)^{1/q}\leq 3\eps$. 
        Combining this with \eqref{eq:computational:caseB} implies that 
        \[
            \cDtrue(\optset\triangle H)\leq 2\eps\,,
        \]
        completing the proof of correctness in Case B.
        
    

    \paragraph{Computational Efficiency.}
        Finally, it remains to show that \cref{alg:main} runs in sample polynomial time.
        
        \begin{lemma}\label{lem:main:rejectionSampling}
            Consider the setting of \cref{thm:main}. \cref{alg:main} runs in time $\poly(n)$ with probability $1-\delta$ where $n=\poly(d^k,\nfrac{1}{\zeta},\log{\nfrac{1}{\delta}})$.
        \end{lemma}
        \begin{proof}[Proof of \cref{lem:main:rejectionSampling}]
            The only computationally non-trivial step in \cref{alg:main} is the call to boosted constraint regression and the construction of and sampling from distributions $\cDgiven_i$ and $\cPtrue_i$ via rejection sampling.
            Since the assumptions of \cref{thm:sampleComplexity} hold (see \cref{lem:main:assumptions}), it promises that the call to boosted constraint regression is efficient and runs in time $\poly(d^k,\nfrac{1}{\zeta},\log{\nfrac{1}{\delta}})$. 
            Hence, it remains to bound the time of rejection sampling.
            
            For this, it suffices to show a lower bound on the mass of $\cD(S_i)$ and $\cPtrue(S_i)$ for each $i$ so that it is not too small, since the inverse of this quantity is the expected number of samples from $\cD$ and $\cPtrue$ one would need to see a single sample in $S_i$. 
            We already proved this in \cref{eq:main:massGuarantee}. 
            In particular, the bounds in \cref{eq:main:massGuarantee} ensure that each call to the sampling oracle in \cref{alg:main}, with probability $1-\eta$, runs in time $O((\nfrac{1}{\eps})\log{\nfrac{1}{\eta}})$ (for any $\eta\in(0,1)$). 
            Selecting $\eta$ to be a suitably small multiple of $\delta$ and taking a union bound over $\poly(d^k,\nfrac{1}{\zeta},\log{\nfrac{1}{\delta}})$ calls implies the claim with only a logarithmic increase in the running time.

        \end{proof}
        
        


         
    \subsection{Proof of \cref{thm:constReg} (Efficient Approximation Algorithm for \pERM{})}
        \label{sec:constReg}\label{sec:efficientPERM}
        In this section, we present a pseudo-approximation algorithm for \pERM{} -- \cref{alg:constReg}.
        \cref{alg:constReg} comes with the guarantees in \cref{thm:constReg}, which we restate below.
        \constReg*
        \noindent The algorithm in \cref{thm:constReg} is summarized in \cref{alg:constReg}.

        \begin{algorithm}[htb!]
        \caption{\textsf{Boosted-}\constrainedreg{}}
        \begin{algorithmic}[1]
        \Procedure{{\rm \textsf{Boosted-}}\constrainedreg{}}{$\zeta,\delta,\sigma,q,\alpha,k$}
            \State Initialize variable $i=1$

            \vspace{2mm}
            
            \State\textit{\#~~Subroutine A (Obtain Samples)}
                \State Obtain sets $P$ and $U$ of $\wt{\Omega}\inparen{
                    \inparen{\nfrac{1}{\zeta^2}}\cdot
                        \inparen{d^k+\log{\nfrac{1}{\delta}}}
                }$ positive and unlabeled samples \label[step]{step:constReg:getSamples} 
                \item[]\quad~ respectively from the PIU learning model (\cref{prob:PUlearning}) 
                \vspace{4mm}
            \State \textit{\#~~Subroutine B (Solve Constrained Regression)}
            \State Find the optimal solution $p_i(\cdot)$ of the following program for $\rho=\frac{2}{\alpha\sigma} \cdot \zeta^{1/(2q)}$ \label[step]{step:constReg:findP} 
            \changetag{Constrained Regression) \ (#1}
            \[
                \min_{{\rm deg}(p)\leq k} \frac{\sum_{x\in U} \abs{{p(x)}}}{\abs{U}}
                \,,~~
                \text{s.t.}\,,~~
                \frac{\sum_{x\in P} \min\inbrace{p(x), 1}}{\abs{P}}     \geq 1  -  \rho
                \,.
                \yesnum\label[program]{eq:constReg:prog}
            \]
            \changetag{#1}
            \State \textit{\#~~This program can be solved by formulating it as a linear program over monomials of degree at}
            \State \textit{\#~~most $k$ and using any polynomial-time linear program solver (see \cref{sec:module:linearProgram} for details)}
            \State Let $H_i=\mathds{1}\sinbrace{p_i(\cdot)\geq t_i}$ where $0\leq t_i\leq 1-\sqrt{\rho}$ is chosen to minimize $\sum_{x\in U}  \mathds{1}\sinbrace{p_i(x)\geq t_i}$\label[step]{step:constReg:selectThreshold}  
            \vspace{4mm}
            \State \textit{\#~~Subroutine C (Boost Performance)}
            \State Repeat the Subroutines A and B for $T=O\sinparen{\inparen{\nfrac{1}{\sqrt{\zeta}}}\cdot \log{\nfrac{1}{\delta}}}$ times to get $H_1,H_2,\dots,H_T$ 
            \State Select fresh sets $P'$ and $U'$ of $\wt{O}\sinparen{\inparen{\nfrac{1}{\zeta}}\cdot \log{\nfrac{1}{\delta}}}$ samples from the PIU model (\cref{prob:PUlearning})\label[step]{step:constReg:freshSamples}
            \State \textbf{return} $H_i$ that minimizes $\abs{H_i\cap U'}$ among $1\leq i\leq T$ 
        \EndProcedure
        \end{algorithmic}
        \label{alg:constReg}
        \end{algorithm} 

        \medskip 
        \paragraph{Notation and Conventions.}
        Before proceeding to the proof, we explain some of the conventions we follow in this section.
        First of all, recall that for each $\ell\geq 1$, $\hyP(\ell)$ is the set of all degree-$\ell$ PTFs over $\R^d$.
        Define $H_{\opt}\in \hyP(k)$ as the following hypothesis 
        \[
            H_{\opt} \coloneqq 
                \min_{H\in \hyP(k)} \cD(H)\,, 
                \quadtext{such that\,,}
                    \cPtrue(H) = 1
                    \,.
                    \yesnum\label{eq:constReg:optimalSolution}
        \]
        The construct $\cP(H_{\opt})$ ensures that $H_{\opt}\supseteq \optset\cap \supp(\cDtrue)$ (except on a measure 0 set) but, in general, when $\supp(\cDtrue)\neq \R^d$, $H_{\opt}$ may be very different from $\optset$.
        Finally, by definition
        \[
            \opt_{\rm PERM} = \cD(H_{\opt})\,.
            \yesnum\label{eq:constReg:expectedValue:defOPT}
        \]
        Next, as a convention, we fix $p_{\opt}\colon \R^d\to \R$ to denote the polynomial which $\zeta$-approximates $H_{\opt}$ with respect to $\cDgiven$ (for $\zeta\leq O\sinparen{{\alpha\sigma}{}}^{2q}$ as specified in \cref{thm:constReg}).
        In other words, $p_{\opt}(\cdot)$ is a polynomial satisfying 
        \[
            \Ex_{\cDgiven}\abs{p_{\opt}(x) - \mathds{1}\sinbrace{x\in H_{\opt}}}\leq \zeta\,.
        \]
        

        \paragraph{Outline of Proof of \cref{thm:constReg}.}
        Now we are ready to prove \cref{thm:constReg}.
        For this, we need to prove the following results 
        \begin{enumerate}[label=(R\arabic*)]
            \item \cref{alg:constReg} is a polynomial-time algorithm;
                \label[result]{result:constReg:polyTimeAlg}
            \item $H$ satisfies $\cDtrue(\optset\setminus H)\leq O\sinparen{\nfrac{1}{(\alpha\sigma)}} \cdot \zeta^{1/(4q)}$ with probability $1-\delta$;
                \label[result]{result:constReg:feasibility}
            \item $H$ satisfies that $\cD(H)\leq \opt_{\rm PERM} + {O(\nfrac{1}{(\alpha\sigma)})\, \zeta^{1/(4q)}} + O(\delta)$ with probability $1-\delta$
            
                \label[result]{result:constReg:optimality}
        \end{enumerate} 
        We present the proofs of these results in the order above. 




        \subsubsection{Proof of \cref{result:constReg:polyTimeAlg}}
            The only computationally non-trivial step in \cref{alg:constReg} is Step~\ref{step:constReg:findP} which solves a constrained optimization problem over polynomials.
            This can be done in polynomial time by casting the problem as a linear program.
            We present the resulting linear program in \cref{sec:module:linearProgram}.
            
        \subsubsection{Proof of \cref{result:constReg:feasibility}}
            Now that we know that \cref{alg:constReg} can be implemented in polynomial time, we proceed to show that the hypothesis $H$ output by it satisfies $\cDtrue(\optset\setminus H)$ is small.
            To show this, we will show that each hypothesis $H_1, H_2, \dots, H_T$ constructed by \cref{alg:constReg} is feasible for \pERM{} with a small $\rho$ which will turn out to be sufficient due to uniform convergence and the constraint in \pERM{}.
            Define the following events 
            \begin{enumerate}[label=(E)] %
                \item[] Let $\evE_i$ (for $1\leq i\leq T$) be the event that the samples in the $i$-th repetition of Subroutine A satisfy $\zeta$-uniform convergence with respect to degree-$k$ PTFs.\label[event]{def:constReg:eventEi}
            \end{enumerate}
            We will show that conditioned on all of the events $\inbrace{\evE_i\colon 1\leq i\leq T}$ happening, each of $H_1,\dots,H_T$, will satisfy $\cDtrue(\optset\setminus H_i)\leq O\sinparen{\nfrac{1}{(\alpha\sigma)}}\cdot \zeta^{1/(4q)}$.
            This will be sufficient to prove \cref{result:constReg:feasibility} since $H$ is selected from among these hypotheses and these events hold with a high probability: 
            \begin{lemma}[Uniform Convergence in Each Repetition]\label{lem:constReg:unifConvergenceHoldsWHP}
                With probability $1-\delta$,  all of the events $\inbrace{\evE_i\colon 1\leq i\leq T}$ hold. 
            \end{lemma}
            \cref{lem:constReg:unifConvergenceHoldsWHP} follows from standard uniform convergence bounds and its proof appears in \cref{sec:proofof:lem:constReg:unifConvergenceHoldsWHP}.
            Next, we show that conditioned on the above set of events happening, feasibility holds.
            \begin{lemma}[{Feasibility for \pERM{} instance in \cref{thm:constReg}}]\label{lem:constReg:feasiblity}
                Conditioned on all of the events $\inbrace{\evE_i\colon 1\leq i\leq T}$ happening, for each $1\leq i\leq T$, $\cDtrue(\optset\setminus H_i)\leq 2\sqrt{2}\sqrt{\nfrac{1}{(\alpha\sigma)} }\cdot \zeta^{1/(4q)}$.
            \end{lemma}
            For each $1\leq i\leq T$, let $P_i$ and $U_i$ be the samples drawn in Subroutine A in the $i$-th repetition.
            To prove \cref{lem:constReg:feasiblity}, we first show that $H_i$ is feasible for \pERM{} over samples $P_i$ and $U_i$.

            \begin{lemma}[{Feasibility for \pERM{} Instances in Each Repetition}]\label{lem:constReg:feasibility:unconditional}
                For each $i$, $H_i$ is feasible for \pERM{} over samples $P_i$ and $U_i$ with tolerance $\rho=
                    \sqrt{\nfrac{2}{(\alpha\sigma)} }\cdot \zeta^{1/(4q)}
                $.
            \end{lemma}
            It is worth noting that unlike \cref{lem:constReg:feasiblity}, this result holds \textit{unconditionally}, in other words, it holds for all draws of $P_i$ and $U_i$.
            The first reason for this is that, for each $P_i$ and $U_i$, $H_i$ is a PTF constructed from thresholding a polynomial $p_i$ that is feasible for \cref{eq:constReg:prog} (Constrained Regression).
            However, this by itself is not sufficient: as the proof below shows it is important that the threshold $t_i$ (which determines $H_i$ as $\mathds{1}\sinbrace{p_i(\cdot)\geq t_i}$) satisfies $t_i\leq 1-\sqrt{\rho}$.
            
            \begin{proof}[Proof of \cref{lem:constReg:feasibility:unconditional}]
                As mentioned above, $H_i$ is the PTF $\mathds{1}\sinbrace{p_i(\cdot)\geq t_i}$ where the polynomial $p_i$ is feasible for \cref{eq:constReg:prog} (Constrained Regression) and the threshold $t_i$ satisfies $0\leq t_i \leq 1- \sqrt{\rho}$.
                Since $p_i$ is feasible for \cref{eq:constReg:prog} (Constrained Regression) with positive samples $P_i$, 
                \[
                    \sum_{x\in P_i} \frac{\min\inbrace{p(x), 1}}{\abs{P_i}}
                    \geq 1-\rho\,.
                \]
                Let $\unif(P_i)$ be the uniform distribution over elements of $P_i$.
                The result follows by applying Markov's inequality on $\max\inbrace{1-p(x),0}$:
                since $1-\min\inbrace{p(x),1}=\max\inbrace{1-p(x),0}$, the above inequality implies that $\Ex_{\unif(P_i)}[\max\inbrace{1-p(x),0}]\leq \rho$.
                Hence, by Markov's inequality, for each $\eta\in (0,1)$, 
                \[
                    \Pr_{x\sim \unif(P_i)}[\max\inbrace{1-p(x),0} > \eta]\leq \frac{\rho}{\eta}\,.
                    \yesnum\label{eq:constReg:feasibility:markov}
                \]
                Substituting $\eta=\sqrt{\rho}$, it follows that 
                \[
                    \frac{\sum_{x\in P_i} \mathbb{I}\insquare{   x\in H_i } }{\abs{P_i}}
                    =
                    \sum_{x\in P_i} 
                    \frac{\mathbb{I}\insquare{
                        p(x) \geq t_i
                    }}{\abs{P_i}}
                    ~~~~~~\Stackrel{t_i \leq 1-\sqrt{\rho}}{\geq}~~~~~~
                    \sum_{x\in P_i} 
                    \frac{\mathbb{I}\insquare{
                        p(x) > 1 - \sqrt{\rho}
                    }}{\abs{P_i}}
                    ~~\Stackrel{\eqref{eq:constReg:feasibility:markov}}{\geq}~~
                    {1-\sqrt{\rho}}\,.
                \]
                The result follows since $\rho$ is set to be $
                    \inparen{\nfrac{2}{(\alpha\sigma)} }\cdot \zeta^{1/(2q)}
                    $ in \cref{alg:constReg}. 
            \end{proof}
            Now, \cref{lem:constReg:feasiblity} follows by appropriately using uniform convergence as we explain below.
            \begin{proof}[Proof of \cref{lem:constReg:feasiblity}]
                Condition on all of the above events happening. 
                Fix any $1\leq i\leq T$ and $\gamma=\sqrt{\nfrac{2}{(\alpha\sigma)} }\cdot \zeta^{1/(4q)}$.
                Since (1) $H_i$ is feasible for \pERM{} over samples $P_i$ and $U_i$ with $\rho={\gamma}$ (\cref{lem:constReg:feasibility:unconditional}), and (2) $P_i$ and $U_i$ satisfy $\zeta$-uniform-convergence with respect to degree-$k$ PTFs due to event $\evE_i$, 
                we get that $\cPtrue(H_i)\geq 1-\zeta-\gamma.$
                Therefore, $\cDtrue(\optset\setminus H_i)\leq \cPtrue(\optset\setminus H_i)\leq \zeta + \gamma.$
                Since $\zeta\leq \gamma\leq 1$ and $\gamma=\sqrt{\nfrac{2}{(\alpha\sigma)} }\cdot \zeta^{1/(4q)}$, it follows that $\cDtrue(\optset\setminus H_i)\leq 2\sqrt{2}\sqrt{\nfrac{1}{(\alpha\sigma)}}\cdot \zeta^{1/(4q)}.$
            \end{proof}
            
            
        \subsubsection{Proof of \cref{result:constReg:optimality}}
            Recall the definition of $H_{\opt}$ from \cref{eq:constReg:optimalSolution} and that $\opt_{\rm PERM} = \cD(H_{\opt})$ (from \cref{eq:constReg:expectedValue:defOPT}).
            We divide the proof of \cref{result:constReg:optimality} into two parts.
            The first part shows that the \textit{expectation} of the objective value attained by each hypothesis $H_1,H_2,\dots,H_T$ from \cref{alg:constReg} is close to $\cD(H_{\opt})$ (\cref{lem:constReg:expectedValue}).
            %
            The second part shows that the selected $H$ has an objective value close to $\cD(H_{\opt})$ with high probability (\cref{lem:constReg:boost}).

            

        \paragraph{Part 1 (Expected Value of Optimal).}
            \begin{restatable}[Expected Value of Optimal]{lemma}{constrainedRegExpValue}\label{lem:constReg:expectedValue}
                For each $1\leq i\leq T$, $H_i$ satisfies
                \[
                    \Ex_{U_i\sim \cDgiven}\insquare{
                        \frac{\abs{H_i\cap U_i}}{\abs{U_i}} 
                    }
                    \leq 
                    \cD(H_{\opt})
                    + O\inparen{\sqrt{\frac{1}{\alpha\sigma}}} \cdot \zeta^{1/(4q)}
                    + O(\delta)
                    \,. 
                \]
            \end{restatable}   

            \paragraph{Outline of the Proof of \cref{lem:constReg:expectedValue} (\cref{sec:proofof:lem:constReg:expectedValue}).}
            Next, we briefly outline the proof of \cref{lem:constReg:expectedValue} and present the complete proof in the next section (\cref{sec:proofof:lem:constReg:expectedValue}).
            This proof is the part of our analysis that is inspired by the analysis of \lreg{} by \citet*{kalai2008agnostically}.
            The proof of \cref{lem:constReg:expectedValue} follows the ideas in the analysis of \lreg{}, but needs to account for some technicalities:
            First, unlike \lreg{} which selects the best threshold in $[0,1]$ (\cref{alg:l1reg}), \cref{alg:constReg} has to select a threshold bounded away from 1 to ensure that the output hypothesis will be feasible for \pERM{} (which is required to establish \cref{result:constReg:feasibility} and, in particular, to prove \cref{lem:constReg:feasibility:unconditional}).
            In the proof, we observe that this only changes the expected value by a small amount.
            Second, the analysis of \lreg{} relies on the polynomial $p_{\opt}$ (which approximates $H_{\opt}$) being feasible for the corresponding optimization problem.
            For \lreg{}, this always holds since the optimization problem solved by \lreg{} is unconstrained (\cref{alg:l1reg}).
            For \cref{alg:constReg} this is not always the case.
            But, we show that $p_{\opt}$ is feasible with high probability, and carry out our proof conditioned on this event.
            This conditioning does not skew the expectations significantly as the involved random variables are suitably bounded. 

        
        \paragraph{Part 2 (Boosting via Repetition).}
        Next, using a standard argument one can show that selecting the ``best performing'' hypothesis from $H_1,H_2,\dots,H_T$ (as evaluated on a fresh set of samples) has an objective value close to optimal with probability $1-\delta$.
        The proof of \cref{lem:constReg:boost} is standard and is provided in \cref{sec:proofof:lem:constReg:boost} for completeness.
        \begin{lemma}[Boosting via Repetition]\label{lem:constReg:boost}
            The hypothesis $H$ output by \cref{alg:constReg}, with probability $1-\delta$, satisfies that 
            \[
                \cD(H)
                \leq 
                \cD(H_{\opt})
                +  
                O\inparen{\sqrt{\frac{1}{\alpha\sigma}}}
                \cdot \zeta^{1/(4q)}
                +
                    O(\delta)
                \,.
            \]
        \end{lemma}

        
        \subsubsection{Proof of \cref{lem:constReg:expectedValue} {(Expected Value of Optimal)}}\label{sec:proofof:lem:constReg:expectedValue}
            In this section, we prove \cref{lem:constReg:expectedValue}, which we restate below.
            \constrainedRegExpValue*
            \noindent Fix any $1\leq i\leq T$.
            We divide the proof of \cref{lem:constReg:expectedValue} into the following lemmas.
            \begin{lemma}\label{lem:constReg:expectedValue:upperBound} 
                It holds that 
                $\nfrac{\abs{U_i\cap H_i}}{\abs{U_i}}\leq \frac{1}{1-\sqrt{\rho}} \cdot \sum_{x\in U_i}  \abs{p_i(x)},$ where $p_i(\cdot)$ is the polynomial constructed in the $i$-th iteration of Subroutine B in \cref{alg:constReg}.
            \end{lemma}
            Recall that $p_{\opt}$ is a polynomial satisfying 
            \[
                \Ex_{{\cDgiven}}\abs{p_{\opt}(x) - \mathds{1}\sinbrace{x\in H_{\opt}}}\leq \zeta\,.
                \yesnum\label{eq:lem:constReg:pstarDef}
            \]
            Let $\evF$ be the event that $p_{\opt}(\cdot)$ is feasible for \cref{eq:constReg:prog} in all iterations $1\leq i\leq T$.  
        \begin{lemma}\label{lem:constReg:expectedValue:feasibility}  
            It holds that $\Pr[\evF]\geq 1-\delta$. 
        \end{lemma}
        Next, we show that conditioned on $\evF$, the expected value of $\sum_{x\in U_i} \abs{p_i(x)} \leq \cDgiven(H_{\opt})$.
        \begin{lemma}\label{lem:constReg:expectedValue:ub2}  
            $\Ex_{U_i\sim \cDgiven}\insquare{
                \inparen{\nfrac{1}{\abs{U_i}}}\sum_{x\in U_i} \abs{p_i(x)}\mid \evF
            } \leq \cDgiven(H_{\opt}) + \zeta +4\delta.$
        \end{lemma}
        Now, we are ready to complete the proof of \cref{lem:constReg:expectedValue}.
        \begin{proof}[Proof of \cref{lem:constReg:expectedValue}]
            Since $\nfrac{\abs{U_i\cap H_i}}{\abs{U_i}}\in [0,1]$, \cref{lem:constReg:expectedValue:feasibility} implies that 
            \[
                \Ex\insquare{
                    \frac{\abs{U_i\cap H_i}}{\abs{U_i}}
                }
                - \Ex\insquare{
                    \frac{\abs{U_i\cap H_i}}{\abs{U_i}}\given \evF
                }
                \leq \delta\,.
            \]
            Further, \cref{lem:constReg:expectedValue:upperBound,lem:constReg:expectedValue:ub2} imply that 
            \[
                \Ex\insquare{
                    \frac{\abs{U_i\cap H_i}}{\abs{U_i}}\given \evF
                }\leq \frac{1}{1-\sqrt{\rho}}\Ex_{U_i\sim \cDgiven}\insquare{
                \sum_{x\in U_i} \abs{p_i(x)}
                    \given \evF
                } \leq
                    \frac{1}{1-\sqrt{\rho}}
                    \cdot \inparen{
                        \cDgiven(H_{\opt})
                        + \zeta
                        + 4\delta
                    }\,.
            \]
            Further, since $\rho=\inparen{\nfrac{2}{\alpha\sigma}} \cdot \zeta^{1/(2q)}\leq \nfrac{1}{4}$ and $\nfrac{1}{(1-z)}\leq 1+2z$ for all $z\in [0,\nfrac{1}{2}]$, 
            \[
                \Ex\insquare{
                    \frac{\abs{U_i\cap H_i}}{\abs{U_i}}\given \evF
                }
                \leq
                {
                    \cDgiven(H_{\opt})
                    + \zeta
                    + 4\delta
                }
                + 2\sqrt{\rho}\inparen{
                    \cDgiven(H_{\opt})
                    + \zeta
                    + 4\delta
                }\,.
            \]
            Finally, substituting the value of $\rho$ and using that $\cD(H_{\opt})+\zeta+4\delta\leq 6$ implies that
            \[ 
                \Ex\insquare{
                    \frac{\abs{U_i\cap H_i}}{\abs{U_i}}\given \evF
                }
                \leq
                {
                    \cDgiven(H_{\opt})
                    +
                    12 
                    \sqrt{\frac{2}{\alpha\sigma}} \cdot \zeta^{1/(4q)}
                    + 4\delta
                    + \zeta
                }
                \,. 
            \] 
            Chaining the above two inequalities implies \cref{lem:constReg:expectedValue}.
        \end{proof}
        In the remainder of this section, we prove \cref{lem:constReg:expectedValue:upperBound,lem:constReg:expectedValue:feasibility,lem:constReg:expectedValue:ub2}. %

        
        \begin{proof}[Proof of \cref{lem:constReg:expectedValue:upperBound}]
            Recall that $H_i=\mathds{1}\sinbrace{p_i(x)\geq t_i}$.
            Since \cref{alg:constReg} selects the threshold $t_i\in [0, 1-\sqrt{\rho}]$, that minimizes $\nfrac{\abs{U_i\cap H_i}}{\abs{U_i}}$, it follows that 
            \[
                \frac{\abs{U_i\cap H_i}}{\abs{U_i}}
                \leq \Ex_{t\sim {\rm Unif}[0, 1-\sqrt{\rho}]} 
                \frac{
                    \sum_{x\in U} \mathbb{I}\insquare{0\leq t\leq p_i(x)}
                }{\abs{U_i}}
                \leq  \frac{1}{{1 - \sqrt{\rho}}}\cdot 
                \frac{
                    \sum_{x\in U_i} \abs{p_i(x)}
                }{\abs{U_i}}
                \,.
                \yesnum\label{eq:constRegAnalysis1}
            \]
        \end{proof}
        \begin{proof}[Proof of \cref{lem:constReg:expectedValue:feasibility}]
            Since $\Ex_{\cDgiven}\abs{\mathds{1}\sinbrace{x\in H_{\opt}} - p_{\opt}(x)}\leq \zeta$, for any $\eta\in (0,1)$
            \[
                \Pr_{x\sim \cDgiven}\insquare{
                    p_{\opt}(x) < 1-\eta
                    ~~\text{and}~~
                    x\in H_{\opt}
                }
                \leq \frac{\zeta}{\eta}
                \,.  
            \]
            Therefore, 
            \[
                \Pr_{x\sim \cDgiven}\insquare{
                    p_{\opt}(x) < 1-\eta
                    ~~\text{and}~~
                    x\in H_{\opt}\cap \optset
                }
                \leq \frac{\zeta}{\eta}
                \,,
            \]
            and now applying \cref{asmp:smoothness} implies that  
            \[
                \Pr_{x\sim \cDtrue}\insquare{
                    p_{\opt}(x) < 1-\eta
                    ~~\text{and}~~
                    x\in H_{\opt} \cap \optset
                }
                \leq \frac{1}{\sigma}\cdot \inparen{\frac{\zeta}{\eta}}^{\sfrac{1}{q}}
                \,.  
            \]
            Since $\cPtrue(H_{\opt})=1$, we also know that 
            \[
                \Pr_{x\sim \cDtrue}\insquare{
                    p_{\opt}(x) < 1-\eta
                    ~~\text{and}~~
                    x\in \optset\setminus H_{\opt} 
                }
                \leq 
                \cDtrue\inparen{\optset\setminus H_{\opt}} = 0
            \]
            Combining the above two, yields 
            \[
                \Pr_{x\sim \cDtrue}\insquare{
                    p_{\opt}(x) < 1-\eta
                    ~~\text{and}~~
                    x\in \optset
                }
                \leq \frac{1}{\sigma}\cdot \inparen{\frac{\zeta}{\eta}}^{\sfrac{1}{q}}
            \]
            In other words,
            \[
                \Pr_{x\sim \cDtrue}\insquare{
                    p_{\opt}(x) < 1-\eta
                    \mid x\in \optset
                }\cdot \cDtrue(\optset)
                \leq \frac{1}{\sigma}\cdot \inparen{\frac{\zeta}{\eta}}^{\sfrac{1}{q}}
                \,.  
            \]
            Now, as $\cDtrue(\optset)\geq  \alpha $ and $\Pr_{x\sim \cDtrue}\insquare{
                    p_{\opt}(x) < 1-\eta
                    \mid x\in \optset
                }=\Pr_{x\sim \cPtrue}\insquare{
                    p_{\opt}(x) < 1-\eta
                    \mid x\in \optset
                }$,
            \[
                \Pr_{x\sim \cPtrue}\insquare{
                    p_{\opt}(x) < 1-\eta
                }
                \leq {
                     \frac{1}{\alpha\sigma}  \cdot \inparen{\frac{\zeta}{\eta}}^{\sfrac{1}{q}}
                }
                \,.  
            \]
            Therefore, we can lower bound the expectation of $\min\inbrace{p_{\opt},1}$ with respect to $\cPtrue$ as
            \[
                \Ex_{\cPtrue}\insquare{\min\inbrace{p_{\opt}, 1}}
                \geq 
                    (1-\eta)\cdot \inparen{
                        1 - {  \frac{1}{\alpha\sigma}   \cdot \inparen{\frac{\zeta}{\eta}}^{\sfrac{1}{q}}}
                    }
                    \geq 1 - \eta - {  \frac{1}{\alpha\sigma}  \cdot \inparen{\frac{\zeta}{\eta}}^{\sfrac{1}{q}}}\,.
                \]
                Substituting $\eta=\zeta^{1/(1+q)}$ implies that 
            \[
                \Ex_{\cPtrue}\insquare{\min\inbrace{p^\star(x), 1}}
                    \geq 1 -  \frac{1}{\alpha\sigma}  \cdot {
                        \zeta^{\sfrac{1}{(1+q)}}
                    }\,.
            \] 
            The result follows
            by uniform convergence over $\cPtrue$ with respect to the hypothesis $\mathds{1}\binbrace{p^\star(x)\geq 1-\zeta^{1/(1+q)}}$ since $\rho$ is set to be $\inparen{\nfrac{2}{(\alpha\sigma)}} \cdot \zeta^{1/(2q)} \geq \inparen{\nfrac{1}{(\alpha\sigma)}} \cdot \zeta^{1/(1+q)}$ in \cref{alg:constReg}.
        \end{proof}
        \begin{proof}[Proof of \cref{lem:constReg:expectedValue:ub2}]
            Recall that $\evF$ is the event that $p_{\opt}(\cdot)$ is feasible for \cref{eq:constReg:prog} (Constrained Regression) defined by samples $P_i$ and $U_i$ drawn in the $i$-th repetition.
            Further, since $p_i(\cdot)$ is optimal for \cref{eq:constReg:prog} (Constrained Regression), conditioned on event $\evF$,
            \[
                \sum_{U_i} \abs{p_i(x)} \leq \sum_{U_i} \abs{p_{\opt}(x)}\,.
            \]
            Hence, in particular, 
            \[
                \Ex_{U_i\sim \cDgiven}\insquare{
                    \frac{
                        \sum_{U_i} \abs{p_i(x)}
                    }{\abs{U_i}}
                \given \evF}
                \leq 
                \Ex_{U_i\sim \cDgiven}\insquare{
                    \frac{
                        \sum_{U_i} \abs{p_{\opt}(x)}
                    }{\abs{U_i}}
                \given \evF}\,.
                \yesnum\label{eq:constReg:bound1}
            \]
            Next, we upper bound the right-hand side of the above inequality.
            By the triangle inequality,
            \begin{align*}
                \hspace{-5mm}
                \Ex_{U_i\sim \cDgiven}\insquare{ \frac{
                        \sum_{U_i} \abs{p_{\opt}(x)}
                    }{\abs{U_i}} \given \evF}
                &\Stackrel{}{\leq}
                \Ex_{U_i\sim \cDgiven}\insquare{
                    \frac{
                        \sum_{U_i}\mathds{1}\sinbrace{H_{\opt}\in x} 
                    }{\abs{U_i}}
                \given \evF}
                + \Ex_{U_i\sim \cDgiven}\insquare{
                    \frac{\sum_{U_i}\abs{\mathds{1}\sinbrace{H_{\opt}\in x} - p_{\opt}(x)}}{\abs{U_i}}
                    \given \evF}\,.
                \yesnum\label{eq:constReg:bound2}
            \end{align*} 
            Next, we upper bound each term on the right-hand side.
            \begin{enumerate}
                \item First, since $\nfrac{
                        \sum_{U_i}\mathds{1}\sinbrace{H_{\opt}\in x} 
                    }{\abs{U_i}}$ is bounded by $0$ and $1$ and $\evF$ happens with probability $1-\delta$, 
                \[
                    \Ex_{U_i\sim \cDgiven}\insquare{
                        \frac{
                            \sum_{U_i}\mathds{1}\sinbrace{H_{\opt}\in x} 
                        }{\abs{U_i}}
                    \given \evF}
                    \leq 
                    \Ex_{U_i\sim \cDgiven}\insquare{
                        \frac{
                            \sum_{U_i}\mathds{1}\sinbrace{H_{\opt}\in x} 
                        }{\abs{U_i}}} + \frac{\delta}{1-\delta}
                    =\cDgiven\inparen{H_{\opt}} + \frac{\delta}{1-\delta}\,.
                \]
                \item Second, since $\nfrac{\sum_{U_i}\abs{\mathds{1}\sinbrace{H_{\opt}\in x} - p_{\opt}(x)}}{\abs{U_i}}$ is bounded between $0$ and $1$ and $\evF$ happens with probability $1-\delta$, 
                \[
                    \Ex_{U_i\sim \cDgiven}\insquare{
                        \frac{
                            \sum_{U_i}\abs{\mathds{1}\sinbrace{H_{\opt}\in x} - p_{\opt}(x)}
                        }{\abs{U_i}}
                    \given \evF}
                    \leq 
                    \Ex_{U_i\sim \cDgiven}\insquare{
                        \frac{
                            \sum_{U_i}\abs{\mathds{1}\sinbrace{H_{\opt}\in x} - p_{\opt}(x)}
                        }{\abs{U_i}}} + \frac{\delta}{1-\delta}\,.
                \]
                Further, recall that $p_{\opt}(\cdot)$ $\zeta$-approximates $H_{\opt}$ with respect to $\cDgiven$, \ie{}, \cref{eq:lem:constReg:pstarDef} holds, which implies that
                \[
                    \Ex_{U_i\sim \cDgiven}\insquare{
                        \frac{
                            \sum_{U_i}\abs{\mathds{1}\sinbrace{H_{\opt}\in x} - p_{\opt}(x)}
                        }{\abs{U_i}}}
                    \leq 
                    \zeta\,.
                \]
            \end{enumerate}
            Finally, substituting the above two bounds into \cref{eq:constReg:bound2} and chaining the resulting inequality with  \cref{eq:constReg:bound1} implies that 
            \[
                \Ex_{\cDgiven\sim U_i}\insquare{ \frac{\sum_{x\in U} \abs{p_i(x)} }{\abs{U}} 
                    \given \evF
                }
                \leq
                \cDgiven\inparen{H_{\opt}}
                +\zeta
                +\frac{2\delta}{1-\delta}\,.
            \]
            The result follows by observing that $\delta\in (0,\nfrac{1}{2})$ and for any $z\in (0,\nfrac{1}{2})$, $\nfrac{z}{(1-z)}\leq 2z$.
        \end{proof}


        
    
        
        

\section{Applications to Other Learning Theory Settings}\label{sec:application}
    In this section, we present the formal statements of the results from \cref{sec:intro:application}.
    \subsection{Learning with Smooth Positive Samples}\label{sec:application:smooth}
        The following result formalizes \cref{infcor:smooth}.
         \begin{corollary}\label{cor:smooth_classification}
                
                Consider a PAC learning instance with      
                    class $\hyH$,
                    feature distribution $\cDtrue$ that is $(\sigma,q)$-smooth (see \cref{asmp:smoothness}) with respect to a known reference $\cD$ that one can efficiently sample from, and 
                    target $h^\star\in \hyH$.  
                \begin{itemize}
                    \item[$\triangleright$] \textit{(Sample Complexity)} 
                        $\wt{O}\inparen{{\eps^{-2q}}\cdot \inparen{\vc{}(\hyH) + \log{\nfrac{1}{\delta}}}}$ independent positive samples are sufficient to $(\eps,\delta)$-PAC learn $\hyH$ with respect to $\cDtrue$.
                    \item[$\triangleright$] \textit{(Computational Complexity)}
                        Suppose $\cDtrue(\optset)\geq \alpha$.
                        Let $k$ be such that degree-$k$ polynomials $\zeta$-approximate $\hyH$ with respect to  $\cDgiven$ in $\nL_1$-norm for $\zeta\leq \inparen{{\alpha^2\sigma}}^{2q}\cdot \eps^{1+4q}$.
                        Then, $n=\poly(d^k,\nfrac{1}{\eps},\log{\nfrac{1}{\delta}})$ positive samples and $\poly(n)$ time is sufficient to $(\eps,\delta)$-PAC learn $\hyH$ with respect to $\cDtrue$.
                \end{itemize}
            \end{corollary}
            \begin{proof}[Proof of \cref{cor:smooth_classification}]
            The setting in \cref{cor:smooth_classification} is a restatement of the standard PIU learning model, and, hence, the result follows from \cref{thm:main,thm:sampleComplexity}.
        \end{proof}
        
        
         
        
    \subsection{Learning with Positive Samples and Lists of Unlabeled Samples}\label{sec:application:listDecoding}
        The following result formalizes \cref{infthm:list}.
        \begin{theorem}\label{thm:listDecoding}
            Consider the List Decoding PIU Model (\cref{def:list}) and $\eps,\delta\in (0,\nfrac{1}{2})$.
            \begin{itemize}
                \item[$\triangleright$] (Sample Complexity)~~ There is an algorithm that, given $\eps$ and 
                $\tilde{O}\binparen{\binparen{\ell\cdot \sinparen{\nfrac{\ell}{(\eps\sigma)}}^{2q} + {\eps^{-3}}}\cdot \sinparen{\vc{(\hyH)} + \log \nfrac{1}{\delta}}}$ independent samples from $\cDtrue(H^\star)$, outputs a hypothesis $\wh{h}$, such that, with probability $1-\delta$,
                \[
                    \Pr_{\cDtrue}\binparen{\wh{h}(x) \neq h^\star(x)} \leq \eps\,.
                \]
                \item[$\triangleright$] (Computational Complexity)~~ 
                    Suppose $\cDtrue(H^\star) \geq \alpha$. 
                    Let $k \geq 1$ be such that $\hyH$ is $\zeta$-approximable by degree-$k$ polynomials in the $\nL_1$-norm with respect to $\cDgiven_{i^\star}$ for the index $i^\star$ in \cref{def:list} and $\zeta\leq O\sinparen{{\alpha^3\sigma}{}}^{2q}\cdot \sinparen{\nfrac{\eps}{\ell}}^{1+4q}$.
                    There is a polynomial time algorithm that given $n=\poly(d^k, \nfrac{1}{\zeta}, \nfrac{\ell}{\eps}, \log \nfrac{1}{\delta})$ samples from $\cDtrue(H^\star)$ and $\cDgiven_1,\cDgiven_2,\dots,\cDgiven_\ell$, outputs a hypothesis $\wh{h}$ such that, with probability $1-\delta$,
                    \[
                        \Pr_{\cDtrue}\binparen{\wh{h}(x) \neq h^\star(x)} \leq \eps\,.
                    \]
                    The algorithm runs in $\poly(n)$ time.
            \end{itemize}
        \end{theorem} 
        Thus, efficient learning is possible under the List Decoding PIU Model.
        See \cref{sec:intro:application:listDecoding} for a discussion of the above results. %
        We stress that the list-decoding model is very general and, in fact, also allows for some cases where 99\% of the unlabeled examples have been \textit{adversarially} corrupted (see \cref{sec:advCorruption}).
            In the remainder of this section, we prove \cref{thm:listDecoding}.
            \begin{proof}[Proof of \cref{thm:listDecoding}]
                First, we prove the computational complexity result. 
                The proof of the sample complexity result is analogous and is presented at the end.

                \paragraph{Computational Complexity.}
                    The algorithm in \cref{thm:main} is not directly applicable in the List Decoding PIU Model because we do not know the index $i^\star$ of the distribution that is ``close'' to $\cDtrue$.
                    However, the simple extension of the algorithm in \cref{thm:main} described in \cref{alg:listDecoding} works.
                    Since each call to \cref{alg:main} uses $n=\poly(d^k, \nfrac{1}{\zeta}, \nfrac{\ell}{\eps}, \log \nfrac{1}{\delta})$  samples and runs in $\poly(n)$ time, the sample complexity and running time claimed in the theorem hold.
                    It remains to prove the correctness of \cref{alg:listDecoding}.
                    Toward this, we claim that, with probability $1-\delta$, the hypothesis $H^{(1)},H^{(2)},\dots,H^{(\ell)}$ satisfy the following properties.
                    \begin{enumerate}[label=(P\arabic*)]
                        \item For each $1\leq i\leq \ell$ that is included in the set $S$ satisfies, $\cDtrue\sinparen{\optset\setminus H^{(i)}}\leq \nfrac{\eps\alpha}{\ell}$.
                        \label{property:listDecoding:proof:1}
                        \item $\cDtrue\sinparen{\optset\triangle H^{(i^\star)}}\leq \nfrac{\eps}{\ell}$.
                        \label{property:listDecoding:proof:2}
                        \item $i^\star\in S$.
                        \label{property:listDecoding:proof:3}
                    \end{enumerate}
                    First, we complete the proof of correctness, assuming that the above claim is true.
                    Condition on the above event happening.
                    First, observe that 
                    \[
                        \cDtrue\inparen{\optset \triangle \wh{H}}
                        = \cDtrue\inparen{\optset \triangle \bigcap\nolimits_{i\in S} H^{(i)}}
                        = 
                        \underbrace{\cDtrue\inparen{\optset \setminus \bigcap\nolimits_{i\in S} H^{(i)}}}_{\text{Term 1: False-Negative Rate}}
                        + 
                       \underbrace{ \cDtrue\inparen{\bigcap\nolimits_{i\in S} H^{(i)} \setminus \optset}}_{\text{Term 2: False-Positive Rate}}\,.
                        \yesnum\label{eq:listDecoding:1}
                    \]
                    These terms can be bounded as follows
                    \begin{align*}
                        \cDtrue\inparen{\optset~\setminus~ \bigcap\nolimits_{i\in S} H^{(i)}}
                     &=\sum_{i\in S} \cDtrue \inparen{\optset\setminus H^{(i)}}
                     ~~\Stackrel{{\rm \ref{property:listDecoding:proof:1}}}{\leq}~~ \ell \cdot \frac{\eps}{\ell}
                     = \eps\,,
                        \yesnum\label{eq:listDecoding:2}\\
                     \cDtrue\inparen{\bigcap\nolimits_{i\in S} H^{(i)}~\setminus~\optset}
                        &~~\Stackrel{\rm\ref{property:listDecoding:proof:3}}{\leq}~~ \cDtrue\inparen{H^{(i^\star)}~\setminus~\optset}
                        ~~\Stackrel{\rm\ref{property:listDecoding:proof:2}}{\leq}~~
                            \eps\,.\yesnum\label{eq:listDecoding:3}
                    \end{align*} 
                The result follows by rescaling $\eps$ to $\nfrac{\eps}{2}$.

                It remains to prove that Properties~\ref{property:listDecoding:proof:1}, \ref{property:listDecoding:proof:2}, and \ref{property:listDecoding:proof:3} hold with probability $1-\delta$.
                Property~\ref{property:listDecoding:proof:2} follows immediately from \cref{thm:main} since the distribution $\cD_{i^\star}$ satisfies the conditions required by \cref{thm:main} by definition of the list decoding model.
                To prove Properties~\ref{property:listDecoding:proof:1} and \ref{property:listDecoding:proof:3}, let $\hyP_{\cap 1/\eps}(k)$ denote the hypothesis class of intersections of $(\nfrac{1}{\eps})$ degree-$k$ PTFs which has VC dimension $\wt{O}\sinparen{\nfrac{d^k}{\eps}}$; hence, in particular, $O(\nfrac{\eps}{\ell})$-uniform convergence holds with respect to the set of samples $P$ constructed in \cref{alg:listDecoding} and the class $\hyP_{\cap 1/\eps}(k)$.
                Moreover, by the construction in \cref{alg:main}, $H^{(1)},\dots,H^{(\ell)}\in \hyP_{\cap 1/\eps}(k)$. 
                Now, Property~\ref{property:listDecoding:proof:1} holds since, for each $i\in S$, we have $\nfrac{\abs{H^{(i)}\cap P}}{\abs{P}}\leq \nfrac{\eps}{(2\ell)}$ which by $O(\nfrac{\eps}{\ell})$-uniform convergence, implies that $\cDtrue(\optset \setminus H^{(i)})\leq \cPtrue(\optset \setminus H^{(i)})\leq \nfrac{\eps}{\ell}$.
                Similarly, Property~\ref{property:listDecoding:proof:3} holds since, from Property~\cref{property:listDecoding:proof:2}, we know that $\cDtrue(\optset\setminus H^{(i^\star)})\leq \nfrac{\eps\alpha}{\ell}$ and, as $\cD(\optset)\geq \alpha$, it follows that $\cPtrue(\optset\setminus H^{(i^\star)})\leq \nfrac{\eps}{\ell}$: furthermore, by $O(\nfrac{\eps}{\ell})$-uniform convergence over $\cPtrue$, we have $\nfrac{\abs{H^{(i^\star)} \cap P}}{\abs{P}}\leq O(\nfrac{\eps}{\ell})$.
                    

            \begin{algorithm}[ht!]
            \caption{Efficient Algorithm to Learn from Positive Samples and List of Unlabeled Samples}\label{alg:listDecoding}
            \begin{algorithmic}
                \Require $\eps, \delta,  \alpha , q, \sigma, k, \ell$ %
                \Ensure Sample access to $\cPtrue=\cDtrue(\optset)$ and each of $\cD_1,\cD_2,\dots,\cD_\ell$
                \State Let $P$ be a set of $m=\wt{O}\inparen{\eps^{-2}\cdot \inparen{\nfrac{\vc(\hyH)}{\eps} + \log{\nfrac{1}{\delta}}}}$ independent samples from $\cPtrue$
                \State Initialize $S$ as the empty set $\emptyset$
                \For{$1\leq i \leq \ell$}
                    \State Obtain $H^{(i)}$ from \cref{alg:main} with positive sample distribution $\cPtrue$ and (imperfect) 
                    \item[] \phantom{.}~~~~ unlabeled distribution $\cDgiven_i$ and parameters: accuracy $\nfrac{\eps}{\ell}$, confidence $\nfrac{\delta}{\ell}$, smoothness 
                    \item[] \phantom{.}~~~~  parameters $(\sigma,q)$, lower bound on mass $\alpha$, and number $k$
                    \vspace{2mm}
                    \State \textbf{If} $\sabs{H^{(i)}\cap P}\geq 1-\nfrac{\eps}{(2\ell)}$, \textbf{then,} add index $i$ to $S$
                    \State \textit{$\#$~~~Guarantees:~~$H^{(i)}$ and $S$ satisfy Properties~\ref{property:listDecoding:proof:1}, \ref{property:listDecoding:proof:2}, and \ref{property:listDecoding:proof:3} with probability $1-\delta$} 
                \EndFor
                \State \Return $\wh{H}\coloneqq \bigcap_{i\in S} H^{(i)}$
            \end{algorithmic}
            \end{algorithm}
 

                \paragraph{Sample Complexity.}
                For the sample complexity result, we use the variant of \cref{alg:listDecoding}, where each call to \cref{alg:main} is replaced by a call to \cref{alg:sampleComplexity}.
                Further, the set of samples $P$ is only of size $\wt{O}\inparen{\eps^{-2}\cdot \inparen{\nfrac{\vc(\hyH)}{\eps} + \log{\nfrac{1}{\delta}}}}$.
                The claimed sample complexity follows from the sample complexity guarantee of \cref{alg:sampleComplexity} in \cref{thm:sampleComplexity} and the specified size of $P$.
                It remains to prove correctness.
                If Properties~\ref{property:listDecoding:proof:1}, \ref{property:listDecoding:proof:2}, and \ref{property:listDecoding:proof:3}  hold, then correctness follows from \cref{eq:listDecoding:1,eq:listDecoding:2,eq:listDecoding:3}.
                So it suffices to prove that Properties~\ref{property:listDecoding:proof:1}, \ref{property:listDecoding:proof:2}, and \ref{property:listDecoding:proof:3}  hold with probability $1-\delta$.
                This proof is identical to the one presented above for computational complexity except that the class $\hyP_{\cap1/\eps}(k)$ is replaced by $\hyH$ and noting that that this improves the factor of $\vc{}(\hyP_{\cap1/\eps}(k))=\wt{O}(\nfrac{d^k}{\eps})$ to $\vc{}(\hyH)$.
            \end{proof}
            
            



    
    \subsection{Truncated Estimation with Unknown Survival Set}  
    \label{sec:application:unknownTruncation}
    In this section, we present the formal statement and proof of \cref{infcor:truncation}, which gives an improved algorithm from truncated estimation with an unknown survival set.

    The most general result for this problem is by \citet{lee2024unknown} and works in the following setting. The unlabeled distribution comes from an exponential family with polynomial sufficient statistics, which is a parametric family of the form
    \[ \cE(x; \theta) \propto h(x) \cdot e^{ \inangle{\theta, t(x)}}\,, \]
    where $t(x)$ is an $m$-dimensional vector whose components are polynomials in $x$.
    The set of parameters $\theta$ for which the density of $\cE(\theta)$ is well defined is called the \textit{natural parameter space} $\overline{\Theta}\subseteq\R^m$.
    Suppose that the ``true'' distribution $\cDtrue$ is described by an unknown target parameter $\theta^\star$ in $\overline{\Theta}$, \ie{}, $\cDtrue=\cE(\theta^\star)$. 
    Given samples from $\cE(\theta^\star)$ truncated to unknown survival set $S^\star \in \hyS$ that is guaranteed to have mass at least $\alpha>0$ under $\cE(\theta^\star)$, \ie{}, $\cE(\theta^\star)(S^\star)\geq \alpha$, the goal is to recover an estimate of $\theta^\star$.
    We begin by describing the assumptions required by existing algorithms \cite{lee2024unknown}. 
    \begin{assumption}\label{asmp:exponentialFamily}
        There is a subset $\Theta$ of the natural parameter space and $\Lambda\geq \lambda > 0$ and $\eta>0$ such that:
        \begin{enumerate}
            \item \textbf{(Interiority)}           
                $\theta^\star$ is in the $\eta$-relative-interior $\Theta(\eta)$ of $\Theta$.
                \label{unknowntrunc:int}
            \item \textbf{(Projection Oracle)} There is an efficient projection oracle to the $\eta$-relative-interior of $\Theta$.
            \label{unknowntrunc:proj}
            \item \textbf{(Bounded Fisher-Information)} For any $\theta\in \Theta$,
            \label{unknowntrunc:fisher}
            $\lambda I \preceq \cov_{x\sim \cE(\theta)}[t(x)]\preceq \Lambda I$.
            \item \textit{\textbf{(Starting Point)} There is an algorithm to find $\theta_0\in \Theta(\eta)$, s.t., $\norm{\theta^\star-\theta_0}\leq \poly(\nfrac{1}{\alpha})$ efficiently.}
            \label{unknowntrunc:st}
        \end{enumerate}
    \end{assumption}
    In addition to the above assumptions, \citet{lee2024unknown} also require two additional assumptions, which we are able to relax.
    \begin{enumerate}[start=5]
            \item \textit{\textbf{($\chi^2$-Bridge)} For each $\theta_1,\theta_2\in \Theta$ there is a $\theta$, s.t., $\chidiv{\cE(\theta_1)}{\cE(\theta)}, 
            \chidiv{\cE(\theta_2)}{\cE(\theta)}
            \leq  
            e^{\poly\inparen{\sfrac{\Lambda}{\lambda}}}$.}
            \label{unknowntrunc:chisquarebridge}
            \item \mbox{\textit{\textbf{($\nL_2$-Approximable)}} $\hyS$ is $\poly(\eps)$-approximable by degree-$k$ polynomials in $\nL_2$-norm w.r.t. $\cE$.}
    \end{enumerate}
    They require these additional assumptions to use \lreg{} as a black box, by showing that $\nL_2$-approximability with respect to the family $\cE$ gives $\nL_1$-approximability with respect to a mixture of a distribution from $\cE$ and the truncation of $\cE(\theta^\star)$ to $S^\star$. 
    We are able to avoid them by using the algorithm from our main result on PIU learning (\cref{thm:main}).
    
    \begin{corollary}[Truncated Estimation with Unknown Survival Set]\label{cor:truncation}
        Suppose \cref{asmp:exponentialFamily} (\cref{unknowntrunc:int,unknowntrunc:proj,unknowntrunc:fisher,unknowntrunc:st}) holds
        and  one has an efficient sampling oracle for $\cE$.\footnote{By this we mean that given any $\theta \in \Theta$, one can draw an independent sample $x$ from a distribution $\cQ$ satisfying $\tv{\cQ}{\cE(\theta)}\leq \delta$ in $\poly(\nfrac{md}{\delta})$ time.} Fix $\eps, \delta \in (0, \nfrac{1}{2})$ and $k$ such that degree-$k$ polynomials $\eps^{D}$-approximate $S^\star$ with respect to family $\cE$ in $\nL_1$-norm for $D=\poly(\nfrac{\Lambda}{\alpha \eta \lambda})$
        and $n = \poly(d^k, \eps^{-D}, \log \nfrac{1}{\delta})$. There is an algorithm that given $\eps, \delta$, constants $(\lambda,\Lambda,\deg(t(\cdot)),\eta,\alpha)$, and $n$ independent samples from $\cE(\theta^\star)$ truncated to $S^\star$, in $\poly(n)$ time, outputs an estimate $\theta$ such that with probability $1-\delta$,
        \[
            \tv{\cE(\theta)}{\cE(\theta^\star)} \leq \eps\,.
        \]
    \end{corollary}
    See \cref{sec:intro:application:truncation} for a discussion of the implications of the above result and \cref{tab:unknown_trunc} for a comparison of our results with existing ones.
    
    \begin{table}[ht!]
        \centering
        \small 
        \begin{tabular}{p{3cm} | l |  C{1.5cm} | C{2.4cm} | C{1.75cm}}
             Distribution 
                    & Hypothesis Class & \cite{Kontonis2019EfficientTS} & \cite{lee2024unknown} & \mbox{This Work}  \\
             \hline
             \multirow{3}{*}{Gaussian $\cN(\mu,\Sigma)$} & Finite GSA & \mbox{\hspace{-2mm}\textit{Diagonal $\Sigma$}\hspace{-2mm}}  & \tick{} &\tick{} \\
             & Arbitrary functions of $O(1)$ halfspaces & \cross{} & \cross{} &\tick{}\\
             & Approximable by polynomials in $\nL_1$-norm & \cross{} & \cross{} & \tick{} \\ 
             \hline
             \multirow{3}{3.5cm}{Log-concave exponential family satisfying \cref{asmp:exponentialFamily}} 
                & Approximable by polynomials in $\nL_2$-norm  
                    & \cross{} 
                        & \textit{\mbox{\hspace{-2mm}Only if $\chi^2$-Bridge\hspace{-2mm}} exists;} see \eqref{unknowntrunc:chisquarebridge}
                            & \tick{} \\ 
              & Arbitrary functions of $O(1)$ halfspaces & \cross{} & \cross{} & \tick{} \\ 
              & Approximable by polynomials in $\nL_1$-norm & \cross{}  & \cross{} & \tick{} \\  
        \end{tabular}
        \caption{
        \textit{Results for Efficient Estimation from Samples Truncated to Unknown Survival Set.} 
            The table includes a \tick{} if there exists an efficient algorithm in the corresponding setting, in the sense that the algorithm runs in {$d^{\poly(1/\eps)}$ time,} and \cross{} otherwise.
        }
        \label{tab:unknown_trunc}
    \end{table}


    
    \newcommand{\customSeparator}{$\phantom{\frac{\frac{1}{1}}{2}}$}
    \begin{table}[tbh]
        \centering
        \midsepremove{}
        \small 
        \begin{tabular}{ C{1.6cm} | C{3.1cm}| C{3.1cm}|C{5.4cm}|C{2.1cm}}
             & Distribution & Access & Hypothesis Class & {\mbox{Computational}} Complexity \\
             \midrule 
             \cite{de2023detectingConvex}  & $\cN(0, I)$ & Known & Convex Sets \customSeparator{} & $\smash{O({d}/\alpha^2)}$ \\
             \hline
             \cite{de2024detecting} & Hypercontractive \mbox{and \iid{} Product \textsuperscript{($\star$)}} & {\mbox{\textit{Exact} orthonormal\textsuperscript{($\star$)\hspace{-3mm}}} \mbox{basis is computable}} & Degree-$k$ PTFs \customSeparator{}
              & $\smash{O(d^{k/2}/\alpha^2)}$  \\
             \hline
             {\multirow{8}{2cm}{This Work}} 
                &\multirow{3}{*}{Gaussian}  
                & \multicolumn{1}{c|}{\multirow{3}{2.8cm}{Approximate sample access\textsuperscript{($\vartriangle$)}}} 
                    &  Convex Sets  
                        & $\smash{{d^{\sqrt{d}\cdot \poly(\sfrac{1}{\alpha})}}}$\\
             && 
                & Degree-$k$ PTF \customSeparator{} & 
                    $\smash{d^{k\cdot \poly(\sfrac{1}{\alpha})}}$\\
             && 
                & 
                    Intersection of $t$ Halfspaces \customSeparator{} & 
                        \hspace{-2mm}$\smash{d^{(\log t)\cdot \poly(\sfrac{1}{\alpha})}}$\hspace{-6mm} \\
            \cline{2-5}
             & \multirow{3}{2.8cm}{~~Log-Concave\textsuperscript{($\dagger$)}} & \multicolumn{1}{c|}{\multirow{3}{2.8cm}{Approximate sample access \textsuperscript{($\vartriangle$)}}}  & 
                    \customSeparator{}Arbitrary functions of $O(1)$ halfspaces & 
                        \customSeparator{}\vspace{-4mm} $d^{\exp\inparen{\poly\inparen{1/\alpha}}}$ \\
                && & 
                    $\poly(\alpha)$-approximable by degree-$k$ polynomials in $\nL_1$-norm & 
                        $d^{k\cdot \poly\inparen{1/\alpha}}$ \\
             \cline{2-5}
             & \multirow{3}{2.8cm}{\mbox{Uniform on $\zo^d$}} & \multicolumn{1}{c|}{\multirow{3}{2.8cm}{Approximate sample access \textsuperscript{($\vartriangle$)}}}  & 
                    Decision trees of size $s$ & 
                       \vspace{-3mm} $d^{(\log s) \cdot \poly(\sfrac{1}{\alpha})}$ \\
                && & 
                    Intersection of $t$ Halfspaces & 
                        $d^{t^6\cdot \poly(\sfrac{1}{\alpha})}$ \\
                && & 
                    Decision tree of halfspaces of size $s$, depth $t$ &
                        $d^{s^2t^6\cdot\poly(\sfrac{1}{\alpha})}$ \\
        \end{tabular}
        \midsepdefault{}
        \caption{\textit{Computational Complexity of Detecting Truncation} 
        \newline
        \vrule height 0.2pt width 15cm\newline
            {\footnotesize \white{.}\quad  ($\star$)~~ The requirements for the distribution in \cite{de2024detecting} are specified in \cref{asmp:detectingTruncation:hypercontractive,asmp:detectingTruncation:product} and access in \cref{asmp:detectingTruncation:access}.}\\[-1mm]
            {\footnotesize \white{.}\quad ($\dagger$)~~ Here, while we do not make explicit assumptions on the distribution, we require approximability by}\\[-1mm]
            {\footnotesize \white{.}\quad ~~~~~~~polynomials with respect to the distribution, which is typically only available for log-concave distributions or}\\[-1mm]
            {\footnotesize \white{.}\quad ~~~~~~~ the uniform distribution on $\zo^d$.}\\[-1mm]
            {\footnotesize \white{.}\quad  ($\vartriangle$)~~ One has sample access to $\cR$ that is close to $\cQ$ as explained in \cref{cor:detectingTruncation:gen}.
            }
        }
        \label{tab:detect}
    \end{table}
    
    
    
    \begin{proof}[Proof of \cref{cor:truncation}]
        This result follows by combining \cref{thm:main} with the reduction in Corollary 6.13 in the arXiv version of \cite{lee2024unknown}.
        The reduction converts the estimation problem in \cref{cor:truncation} to the following learning problem in $\poly(\nfrac{d}{(\eps\alpha)})$ time:
            given samples from a distribution $\cD=\cE(\theta_0)$ that is close to $\cDtrue=\cE(\theta^\star)$ in a sense stronger than that of \cref{asmp:smoothness} with $\nfrac{1}{\sigma}=q=C$ for $C=\poly\inparen{\nfrac{\Lambda}{(\alpha\eta\lambda)}}$, find a set $S$ such that $\cDtrue(S\triangle S^\star)\leq \eps^{O(C)}$.
        This precisely fits into the framework of PIU learning with $\optset=S^\star$ and, moreover, since in the model of \citet{lee2024unknown}, $\cDtrue(S^\star)\geq \alpha$ which ensures \cref{asmp:posFraction} holds allowing us to apply \cref{thm:main} directly.
    \end{proof}
    \subsection{Detecting Truncation}  
    \label{sec:application:detectingTruncation} 
        In this section, we present the formal statement and proof of \cref{infcor:detectingTruncation}. 
        The following two results formalize \cref{infcor:detectingTruncation} in the cases where we have samples from the underlying distribution $\cQ$ or only from a distribution $\cR$ close to it, respectively.
        \begin{corollary}[Detecting Truncation with Exact Samples]\label{cor:detectingTruncation}
            Fix constant $\beta\in (0,1)$, set $S^\star\in \hyS \subseteq \R^d$, and distribution $\cQ$ over $\R^d$, such that, $\cQ(S^\star)\leq 1-\beta$.
            The following holds. 
                    Let $k$ be such that degree-$k$ polynomials $\zeta$-approximate $\hyS$ in $L_1$-norm with respect to $\cQ$ for $\zeta\leq \beta^4 \eps^{5}$.
                    There is an algorithm that given $\beta,\delta,k$ and $n=\poly(d^k,\nfrac{1}{\zeta},\log{\nfrac{1}{\delta}})$ samples (a) $\cD\in\inbrace{\cQ,\cQ(S^\star)}$ and (b) $\cQ$, with probability $1-\delta$, correctly detects whether $\cD=\cQ$ or $\cQ=\cQ(S^\star)$.
                    The algorithm runs in $\poly(n)$ time.
        \end{corollary}
        The next corollary strictly generalizes \cref{cor:detectingTruncation}.
        \begin{corollary}[Detecting Truncation with Approximate Samples]
            \label{cor:detectingTruncation:gen}
            Consider the setting in \cref{cor:detectingTruncation}.
            Let $\cR$ be a distribution such that $\cQ$ is $(\sigma,q)$-smooth with respect to $\cR$ (\ie{}, for any measurable set $S$, $\cQ(S)\leq (\nfrac{1}{\sigma})\cdot \cR(S)^{1/q}$).
                    Let $k$ be such that degree-$k$ polynomials $\zeta$-approximate $\hyS$ in $L_1$-norm with respect to $\cQ$ for $\zeta\leq ({\sigma\alpha^2}{})^{2q}\cdot \eps^{1+4q}$.
                    There is an algorithm that given $\beta,\delta,k$ and $n$ samples (a) $\cD\in\inbrace{\cQ,\cQ(S^\star)}$ and (b) $\cR$, with probability $1-\delta$, correctly detects whether $\cD=\cQ$ or $\cQ=\cQ(S^\star)$.
                    The algorithm runs in $\poly(n)$ time.
        \end{corollary}
        See \cref{tab:detect} for a summary of results and prior work. 

    
    
    \begin{proof}[Proof of \cref{cor:detectingTruncation}]
        \cref{cor:detectingTruncation} follows from \cref{cor:detectingTruncation:gen} by substituting $\sigma=q=1$.
    \end{proof}
    \begin{proof}[Proof of \cref{cor:detectingTruncation:gen}]
        In this application, the samples are observed from the ``truncated'' distribution $\cQ(T)$ where $T=\R^d$ if there is no truncation and $T=S^\star$ if there is a truncation.
        Further, since $\cQ(S^\star)\leq 1-\beta$.
        To distinguish the two cases, it suffices to learn $T$ up to $O(\beta)$ accuracy:
        this is because, for any set $T'$ satisfying $\cQ(T\triangle T') < \nfrac{\beta}{2}$, if $\cQ(T') > 1-\nfrac{\beta}{2}$, then it must be that $\cQ(T) > 1-\beta$ which is only possible if $T=\R^d$ (since $\cQ(S^\star)\leq 1-\beta$).
        Now, the result follows from the sample and time complexity of learning in \cref{thm:main}.
    \end{proof}

\ifanon
\else 
\section*{Acknowledgments}
    We thank Alkis Kalavasis and anonymous reviewers for helpful comments on a draft of this paper which helped improve its presentation.
    Jane H.\ Lee was supported by a Graduate Fellowship for STEM Diversity sponsored by the U.S.\ National Security Agency (NSA). 
\fi 




\newpage
\printbibliography
\newpage

\appendix 
\addtocontents{toc}{\protect\setcounter{tocdepth}{1}}

\section{Additional Related Works}\label{sec:appendix:related}
        

        
       Our problem setting and results have connections to many learning models, for which we present the following non-exhaustive list. Some of these have been mentioned in the main body (in \cref{sec:intro:relatedWorks}), but we expand on some of these here as well as mention additional ones not mentioned before.

        One important class of problems between realizable and agnostic learning model the outcomes as $y_i = h^{\star}(x_i) + \eta_i$, where $\eta_i$ is a structured noise, \eg{}, see \cite{angluin1988learning, decatur1997cpcn, blum1998polynomial, feldman2006new, natarajan2013learning}.  In this paper, we consider a particular aspect of this structured noise that has many practical instantiations. In particular, we consider the case where there is significant asymmetry between positive and negative samples, and the quality of the data is very poor for one of the two types of samples -- say negative samples. (See \cref{sec:realWorldApplications} for further examples.)
    
         \paragraph{Learning From Positive Samples Only (Extension of Discussion in \cref{sec:intro:relatedWorks}).} 
         Apart from the line of works discussed in \cref{sec:intro:relatedWorks}, another line of work \cite{frieze1996linearTransformations,Nguyen2009learningParallelepiped,Rademacher2009hardnessOfLearningConvex,anderson2013simplices,de2015satisfying,canonne2020satisfying,he2023robustMoments} studies learning-from-positive-samples in a setting where the distribution of the underlying examples $\cDtrue$ is explicitly known to the learner.
         For instance, $\cDtrue$ may be known to be uniform or known to be Gaussian with a known mean and covariance.
         In this model, these works provide algorithms and hardness results for learning specific hypothesis classes such as simplices, degree-2 PTFs, and parallelpipes. 
         Compared to these works, we consider the setting where $\cDtrue$ is unknown (we see samples from it) and more general hypothesis classes (any function approximable by $L_1$-degree polynomials).

         \paragraph{{(Language) Learning from Positive Samples Only.}} 
         Starting from the seminal works of \citet{gold1967language,angluin1979finding,angluin1988learning}, a huge body of works tackles the language identification problem: given (positive) examples from a target language $K \in \mathcal{L} = \{L_1, L_2, \dots\}$, identify the index $i$ such that $K = L_i$. The challenge is that only samples from within $K$ are available; with both positive and negative examples, identification is feasible, but without negatives, strong impossibility results hold. Recent work extends these ideas to language generation -- where the goal is to output an unseen string $x\in K$  
        [\plaincite{kleinberg2024language}; \plaincite{li2024generationlenslearningtheory}; \plaincite{kalavasis2025limitslanguagegenerationtradeoffs}; \plaincite{charikar2024exploringfacetslanguagegenerationV2}; \plaincite{kalavasis2024characterizationslanguagegenerationbreadth}; \plaincite{raman2025generationnoisyexamples}].
        While all of these works also deal with issues of ``learning'' with positive examples, they do not consider computational and sample complexity, leading to techniques that differ markedly from those in our work.
            
        \paragraph{Learning from Positive and (Perfect) Unlabeled Samples.} 
        \citet{denis1990positive} initiated the study from positive and unlabeled samples, showing that any class that is PAC-learnable is also PAC-learnable from positive and unlabeled samples (in the realizable setting). To show this, they reduced the (realizable) PAC-learnable from positive and unlabeled samples to learning in the CPCN model with noise bounded away from $\nfrac{1}{2}$.
       Since then, a growing number of works in Applied Machine Learning design efficient heuristics for learning with positive and unlabeled samples, \eg{}, \cite{elkan2008pu,mordelet2011prodige,plessis15convex,sakai2019covariate, bekker2020learning,reviewPUBioformatics2021}; we refer the reader to the survey of \citet{bekker2020learning} for an overview. 
        The heuristics in this line of work inspire the optimization program (\pERM{}) at the core of our approach.
 
            
             

        \paragraph{{Regression} with Covariate Shift.}
            As discussed in \cref{sec:intro:relatedWorks}, several works on PAC learning with covariate shift provide sufficient guarantees on the source and traget distributinos, under which a classifier trained on the source distribution ``transfers'' to the target distribution, \ie{}, it has low error on the target distribution.
            A recent work \citet{kalavasis2024transfer} showed how to transfer beyond the usual assumptions made by prior works for the regression task.
            However, their techniques fail for binary classification which is the focus of our study.


         
             

        \paragraph{List-Decodable Learning {(Extension of Discussion in \cref{sec:intro:relatedWorks})}.}
            Almost all statistical tasks become information-theoretically impossible when a large fraction of samples is adversarially corrupted.
            The simplest example is perhaps Gaussian mean estimation: if an adversary can corrupt $(1-\gamma)$-fraction of the samples (for $\gamma\in (0,1)$), then the corrupted samples can contain $\nfrac{1}{\gamma}$-clusters that have a large separation, any of which could correspond to the true distribution.
            To circumvent these impossibility results, \citet{charikar2017learning} proposed a relaxed notion of learning, termed \textit{list-decodable learning}, where the learner is allowed to output a list of estimates one of which must be close to the correct parameter.
            Since the formulation of this problem, there has been a growing line of works designing list-decodable algorithms for several fundamental statistical problems  \cite{charikar2017learning,karmalkar2019list,Cherapanamjeri2020DecodableNearlyLinear,bakshi2021Subspace,Abhimanyu2023DecodableBatches,Diakonikolas2023DecodableCovariance,kothari2022decodableCovariance,Raghavendra2020Decoding,Diakonikolas2022DecodableSparseMean,Diakonikolas2021SQ_LB_Decoding,Diakonikolas2021PCA_Decodable}.
            As we mentioned in \cref{sec:intro:relatedWorks}, these (and future) works can be combined with our techniques to learn in the stronger list-decodable PIU learning model.

\section{Additional Preliminaries and the \lreg{} Algorithm}
    In this section, we first present the definitions of concept classes mentioned in the main body and then present the \lreg{} algorithm of \citet*{kalai2008agnostically} (see \cref{alg:l1reg}).
    
    Recall that a halfspace is a Boolean function $h\colon \R^d \rightarrow \{0, 1\}$ of the form $h(x) = \mathds{1}\sinbrace{w^\top x + t \geq 0}$ for some vector $w \in \R^d$ and threshold $t\in \R$. 
    An intersection of $k$ halfspaces $h_1,h_2,\dots,h_k$ is the Boolean function formed by taking the logical AND of the halfspaces, \ie{}, $h_1(x)\cdot h_2(x)\cdot \dots\cdot h_k(x)$.
    Functions of halfspaces are a significant generalization of intersections of halfspaces.
    \begin{definition}[Functions of Halfspaces]
        The concept class of functions of halfspaces consists of all Boolean functions of the form $b(h_1,h_2,\dots,h_k)$ where $b\colon \zo^k\to \zo$ is an arbitrary boolean function and $h_1,h_2,\dots,h_k$ are any $k$ halfspaces.
    \end{definition}
    Next, we formally define the class of decision trees of a specific size and/or depth.
    \begin{definition}[Decision Tree]
        A decision tree is a rooted binary tree in which each internal node is labeled with a variable $x_i$ (for $1\leq i\leq d$) and has two children. 
        Each leaf is labeled with an output from $\{0, 1\}$.
    \end{definition} 
    A decision tree computes a Boolean function in the following standard way: Given an input $x \in \{0,1\}^d$, the value of the function on $x$ is the output in the leaf reached by the path that starts at the root and goes left or right at each internal node according to whether the variable’s value in $x$ is 0 or 1, respectively.
    The \textit{size} of a decision tree is the number of leaves of the tree. The \textit{depth} of a node in a decision tree is the number of edges in the path from the root to the node.
    The \textit{depth of the tree} is the maximum depth of one of its leaves.
    Observe that a decision tree of depth $d$ has size at most $2^d$.
    This naturally defines the following concept classes.
    \begin{itemize}
        \item The family of \textit{decision trees of depth at most $t$} over $d$ variables; %
        \item The family of \textit{decision trees of size at most $s$} over $d$ variables.
    \end{itemize}
    When the number of variables $d$ is clear from context or not important, we simply express the above classes as decision trees of depth $t$ and of size $s$ respectively.
    Further, by taking the intersections of these classes, one can also define the concept class of decision trees of size at most $s$ and depth at most $t$.
    
    Decision trees of halfspaces are a generalization of decision trees, where each node $i$ is labeled with a halfspace $h_{i_j}$ instead of a single variable $x_{i_j}$.
    \begin{definition}[Decision Tree of Halfspaces]
        A decision tree of halfspaces is a rooted binary tree in which each internal node $i$ is labeled with a halfspace $h_{i_j}$ (for $1\leq i\leq d$) and has two children. 
        Each leaf is labeled with an output from $\{0, 1\}$.
    \end{definition} 
    The computation, on a decision tree of halfspaces. proceeds in a similar fashion as above: the only difference is that one goes to the left child from node $i$ if $h_{i_j}(x)=1$ and goes to the right child, otherwise, when $h_{i_j}(x)=0$.
    This naturally gives rise to the following concept class
    \begin{itemize}
        \item The family of \textit{decision trees of halfspaces of size at most $s$ and depth at most $t$} over $d$ variables. 
    \end{itemize}
    Observe that the above class can also be written as a sum of at most $s$ intersections of $\ell$ halfspaces.
    \begin{algorithm}[h!]
        \caption{\lreg{} \cite{kalai2008agnostically}}
        \begin{algorithmic}[1]
        \Procedure{{\rm\lreg{}}}{$\eps,\delta,k$}
                \vspace{2mm}
                \State \textit{\#~~Construct Samples}
                    \State Let $n=\wt{\Omega}\inparen{\sfrac{\inparen{d^k+\log\inparen{\nfrac{1}{\delta}}}}{\eps}}$ 
                    \State Draw $n$ samples $D=\inparen{(x_i,y_i)}_i$ from $\cD$
                    \vspace{2mm}
                \State \textit{\#~~Solve Regression}
                \State Find a polynomial $p$ of degree at most $k$ which minimizes $\inparen{\nfrac{1}{n}}\sum_{i}\abs{y_i-p(x_i)}$.
                \State \textit{\#~~This can be done by expanding all examples to monomials of degree at most $d$ and then}
                \State \textit{\#~~performing $\nL_1$-linear-regression using any polynomial time LP-solver.}
                \State  $H=\inbrace{x\colon p(x)\geq t}$ where $0\leq t\leq 1$ minimizes the error of the hypothesis on $D$
                \vspace{2mm}
                \State \textit{\#~~Boost Performance}
                \State Repeat the above steps $T=O\inparen{\inparen{\nfrac{1}{\eps}}\log\inparen{\nfrac{1}{\delta}}}$ times to get $H_1,H_2,\dots,H_T$
                \State Evaluate $H_1,H_2,\dots,H_T$ on a fresh set of $O\inparen{\inparen{\nfrac{1}{\eps^2}}\log\inparen{\nfrac{1}{\delta}}}$ samples
                \State \textbf{return} the best-performing hypothesis on $D$
            \EndProcedure
        \end{algorithmic}
        \label{alg:l1reg}
    \end{algorithm}

    

\section{Proof of \cref{thm:sampleComplexity:bothSided} (Optimal Sample Complexity with \cref{asmp:densityLB})} 
    In this section, we prove \cref{thm:sampleComplexity:bothSided}, which obtains the optimal sample complexity for PIU learning when $q=1$; when an additional assumption (\cref{asmp:densityLB}) holds. %
    \optimalSampleComplexity*
    \noindent We refer the reader to \cref{tab:sample} for a comparison of this result with those of prior works. 

    We also prove the following more general result, which allows the hypothesis to be agnostic.
    \cref{thm:sampleComplexity:bothSided} can be deduced from this result by setting $\opt=0$.\footnote{To be precise, we also need to check that there are no issues created in the proof due to $\alpha$ being $0$. This is a subtle point and turns out to not create an issue because \cref{thm:sampleComplexity:bothSided} works under the ``strict'' realizability assumption that there exists an hypothesis $\optset\in \hyH$ which agrees with the data on the whole domain (not just a set of measure 1). 
    Our computationally efficient algorithm does not require this.
    }
    \begin{theorem}\label{new:thm:perm:robust}
        Suppose \cref{asmp:posFraction,asmp:smoothness} hold, allow $\optset\not\in \hyH$, and define $\opt\coloneqq \inf_{H\in \hyH} \cDtrue(H\triangle \optset)$.
        Fix any parameters $\eps,\delta, \gamma \in (0,\nfrac{1}{2})$ with $\eps\geq \gamma > {\nfrac{\opt}{\alpha}}$.  
        Consider the \pERM{} instance defined with respect to $n=\wt{O}\inparen{\inparen{\nfrac{1}{\eps}}\cdot \inparen{\vc{}(\hyH)+\log{\nfrac{1}{\delta}}}}$  positive and unlabeled samples obtained from the two streams in the PIU learning model (\cref{prob:PUlearning}).
        Then, with probability $1-\delta$, any hypothesis $H\in \hyH$ that satisfies:
        \begin{enumerate}
            \item $H$ is feasible for \pERM{} with tolerance $\rho= \gamma $; and 
            \item $H$ is optimal (though not necessarily feasible) for \pERM{} with tolerance $\rho= {5\eps} $,
        \end{enumerate}
        also satisfies
        \[
            \cDtrue\inparen{H\triangle \optset} \leq 
                \frac{144}{\sigma^2}\cdot \, \eps^{1/q^2}
                \,.
        \]
    \end{theorem}
        

    
    \begin{figure}[h]
        \centering
        \vspace{4mm}
        \begin{tikzpicture}[node distance=2.5cm, auto]
        \node[myNodeNarrow, line width=1pt, text width=7.25cm] (thmAgnostic) at (-4.5,0) 
            {{Sample Efficient Algorithm (Agnostic)\\[1mm] \cref{new:thm:perm:robust}}};

        \node[myNodeNarrow, line width=1pt, text width=7.25cm] (thmRealizable) at (4.5,0) 
            {{Sample Efficient Algorithm (Realizable)\\[1mm] \cref{thm:sampleComplexity}}};

        \draw[-stealth, line width=0.5mm] (thmAgnostic) -- (thmRealizable); 

        \node[myNode, line width=1pt, text width=7.25cm] (lemEasy) at (-4.5,-2.2) 
            {{Upper bound on False Negative Error\\ $\cD(\optset\setminus H)~{\leq}~\poly_q(\eps)$\\[1mm]\cref{lem:perm:easy}}};

        \node[myNode, line width=1pt, text width=7.25cm] (lemHard) at (4.5,-2.2) 
            {{Upper bound on False Positive Error\\ $\cD(H\setminus\optset)~{\leq}~\poly_q(\eps)$\\[1mm]\cref{lem:perm:hard}}};

        \draw[-stealth, line width=0.5mm] (lemEasy) -- (thmAgnostic); 
        \draw[-stealth, line width=0.5mm] (lemHard) -- (thmAgnostic); 

        \node[myNode, line width=1pt, text width=4.5cm] (lemTerms23) at (0.15,-4.65) 
            {$\cD(H'){\,-\,}\cD(\optset){\,\leq\,}\poly_q(\eps)$\\$\cD(\optset\setminus H){\,\leq\,}\poly_q(\eps)$\\[1mm] \cref{lem:sampleComplexity:Hard:UBterms23}};
        \node[myNode, line width=1pt, text width=4.5cm, minimum height=1.75cm] (lemTerm1) at (5.88,-4.65) 
            {$\cD(H)-\cD(H')~{\leq}~\poly_q(\eps)$\\[3mm]\cref{lem:sampleComplexity:Hard:UBterm1}};

        \draw[-stealth, line width=0.5mm] (lemTerms23) -- (lemHard); 
        \draw[-stealth, line width=0.5mm] (lemTerm1.north) -- ($(lemTerm1.north |- lemHard.south)$);


        \node[myNode, line width=1pt, text width=4.5cm] (lemFeasibility) at (5.88,-6.85) 
            {Feasibility of $H'$\\[1mm]\cref{lem:sampleComplexity:Hard:feasibility}};

        \draw[-stealth, line width=0.5mm] (lemFeasibility) -- (lemTerm1); 

        \node[myNode, dashed, line width=1pt, text width=7.25cm] (unifConv) at (-4.5,-6.85) 
            {Second-Order Uniform Convergence\\[1mm]\cref{thm:unifConvergence} from, \eg{},  \cite{Boucheron_Bousquet_Lugosi_2005}};

        \draw[-stealth, line width=0.5mm] (unifConv) -- (lemFeasibility); 
        \draw[-stealth, line width=0.5mm] (unifConv) -- (lemTerms23); 
        \draw[-stealth, line width=0.5mm] (unifConv) -- (lemEasy); 
        
        \end{tikzpicture}
        \caption{Outline of the proofs of \cref{thm:sampleComplexity:bothSided,new:thm:perm:robust}.}
        \label{fig:outline:sampleComplexity}
    \end{figure}
    
    \noindent The proof of \cref{new:thm:perm:robust} is divided into two parts: \cref{lem:perm:easy} bounds $\cDtrue(H \backslash H^\star)$ and \cref{lem:perm:hard} bounds $\cDtrue(H^\star \backslash H)$.
        Taking a union bound over the events in \cref{lem:perm:easy,lem:perm:hard} and recalling that $\optset\triangle H=\inparen{H\setminus \optset}\cup \inparen{\optset\setminus H}$ completes the proof of \cref{new:thm:perm:robust}.
        \begin{lemma}[Upper Bound on False Negative Error]\label{lem:perm:easy} 
            Fix any $\eta\in (0,1).$
            Under the setting of \cref{new:thm:perm:robust}, consider any hypothesis $T\in \hyH$ feasible for the \pERM{} instance in \cref{new:thm:perm:robust} with tolerance $\rho=\eta$.
            Then, with probability $1-\inparen{\nfrac{\delta}{2}}$,
            \[
                \cDtrue\inparen{\optset\setminus T} \leq \eta+\sqrt{\eta\eps}+\eps\,.
            \]
            In particular, for $T=H$ and $\gamma\leq \eps$, with probability $1-\inparen{\nfrac{\delta}{2}}$, $\cDtrue\inparen{\optset \setminus H}\leq 3\eps$.
        \end{lemma}
        \begin{lemma}[Upper Bound on False Positive Error]\label{lem:perm:hard} 
            Under the setting of \cref{new:thm:perm:robust}, with probability $1-\inparen{\nfrac{\delta}{2}}$,  
            \[
                \cDtrue\inparen{H\setminus\optset} \leq 
                \frac{13}{\sigma^2} \, (\nfrac{9\eps}{\alpha})^{1/q^2}\,.
            \]
        \end{lemma}
        The proofs of both \cref{lem:perm:easy,lem:perm:hard} rely on the following second-order uniform convergence bound, which gives sharper uniform convergence rates for hypotheses with small errors.

        \begin{remark}
            For simplicity, in \cref{new:thm:perm:robust}, we state the result in the regime $\opt\leq\eps$. %
            In the regime $\opt > \eps$,  one can deduce the following error bounds of the form $O(\nfrac{\opt\cdot \eps}{\alpha})^{1/O(q^2)} + O(\eps)^{1/O(q^2)}$ from the same proof with additional bookkeeping.
        \end{remark}
        
        \subsection{Proof of \cref{lem:perm:easy} (Upper Bound on False Negative Error)}
        \begin{proof}[Proof of \cref{lem:perm:easy}]
            Let $T\in\hyH$ be any solution of \pERM{} with $\rho=\eta$. By feasibility, it satisfies $\nfrac{\abs{ T \cap P}}{\abs{P}}\geq 1- \eta $ and, hence, $\nfrac{\abs{ P\setminus T}}{\abs{P}}\leq \eta$.
            Due to the choice of the sample size $n$ and \cref{thm:unifConvergence}, with probability $1-\inparen{\nfrac{\delta}{2}}$, for each $R\in \hyH\cup \inbrace{G^c\colon G\in \hyH}$,\footnote{Here, we use the facts that $\vc{(\inbrace{G^c\colon G\in \hyH})}=\vc{(\hyH)}$ and $\vc{(\hyH_1\cup\hyH_2)}\leq {\vc{(\hyH_1)}+\vc{(\hyH_2)}+1}$ \cite{mohri2018foundations}.}
            \[
                \cPtrue(R) 
                \leq 
                \frac{\abs{R\cap P}}{\abs{P}}
                +
                \sqrt{\frac{\abs{R\cap P}}{\abs{P}}\cdot \eps}
                + \eps\,.
            \]
            In particular, applying this to $R=T^c$ (note that $T^c\in\{G^c: G\in \hyH\}$), and using $\nfrac{\abs{ P\cap T^c}}{\abs{P}}=\nfrac{\abs{ P\setminus T}}{\abs{P}}\leq \eta$, implies that 
            \[
                \cPtrue(T^c) \leq \eta+\sqrt{\eta\eps}+\eps\,.
            \]
            Finally, since $\cDtrue( R )\leq \cPtrue( R )$ for any $ R \subseteq \optset$, it follows that 
            \[
                \cDtrue(\optset\setminus  T )
                \leq \cPtrue(\optset\setminus  T )
                = \cPtrue(T^c)
                \leq  \eta +\sqrt{\eta\eps}+\eps\,.
            \]  
        \end{proof}


        
        \subsection{Proof of \cref{lem:perm:hard} {(Upper Bound on False Positive Error)}}
            Our goal is to upper bound the mass of $H\setminus \optset$ under $\cDtrue$. 
            One natural idea is to use the fact the $\eps$-optimality of $H$ for \pERM{}, \ie{},  $\cDgiven(H)$ is smaller than $\cD(H')$ for any other $H'$ feasible for \pERM{}, and then convert from $\cDgiven$ to $
            \cDtrue$ via \cref{asmp:smoothness}.
            This approach does not go through as $\optset$ is not guaranteed to be feasible for \pERM{} (recall that in \cref{new:thm:perm:robust} we do not assume that $\optset\in \hyH$). %
            
            Instead, we construct a hypothesis $H'$ that is both ``close'' to $\optset$ and is also feasible for \pERM{}.
            Recalling that $\opt\coloneqq \inf_{H'\in \hyH}\cDtrue(H'\triangle \optset)$, one can see that for any $\zeta>0$, there is an hypothesis $H'=H'_\zeta$ in $\hyH$ satisfying
            \[
                \cDtrue\sinparen{H'\triangle \optset} 
                \leq \opt + \zeta\,.
                \yesnum\label{eq:perm:closeToInf}
            \]
            Consider any $0<\zeta\leq \alpha\eps$ satisfying $\gamma \geq \frac{\opt+\zeta}{\alpha}$, which exists since $\gamma > \nfrac{\opt}{\alpha}$.
            Fix $H'=H'_\zeta$ for the remainder of the proof.
            We now show that, with high probability, $H'$ is feasible for \pERM{} with tolerance $\rho=O(\eps)$.
            \begin{lemma}[Feasibility of $H'$]
                \label{lem:sampleComplexity:Hard:feasibility}
                Let $\evE$ be the event that the uniform convergence bounds in \cref{thm:unifConvergence} hold for the positive samples $P$ and the hypothesis class $\hyH\cup\inbrace{G^c\colon G\in \hyH}$.
                Then, $\Pr[\evE]\geq 1-(\nfrac{\delta}{4})$.
                Moreover, conditioned on $\evE$, $H'$ is feasible for \pERM{} defined by samples $(P,U)$ with $\rho=5\eps$.
            \end{lemma}
            \begin{proof}[Proof of \cref{lem:sampleComplexity:Hard:feasibility}]
                The lower bound on the probability of event $\evE$ follows directly from \cref{thm:unifConvergence}.
                We now prove the feasibility of $H'$.
                Since $\cPtrue$ is a truncation of $\cDtrue$ to $\optset$, $ \cPtrue\sinparen{H'\triangle \optset} \leq \nfrac{\cDtrue\sinparen{H'\triangle \optset}}{\cDtrue(\optset)}\leq \nfrac{\inparen{\opt+ \zeta }}{{\cDtrue(\optset)}}$ and, further, since $\cDtrue(\optset)\geq  \alpha $, 
                \[
                    \cPtrue\inparen{H'\triangle \optset} \leq \frac{\opt +  \zeta }{\alpha}\,.
                \]
                Condition on event $\evE$.
                The uniform convergence bound from \cref{thm:unifConvergence} implies that 
                \[
                    \frac{\abs{P\setminus H'}}{\abs{P}}
                    \leq
                        \cPtrue\inparen{{H'}^c}
                        + \sqrt{ \cPtrue\inparen{{H'}^c} \cdot \eps}
                        + \eps\,.
                \]
                Combining the last two inequalities implies
                \[
                    \frac{\abs{P\setminus H'}}{\abs{P}}
                    \leq
                        \frac{\opt +  \zeta }{\alpha}
                        + \sqrt{ \frac{\opt +  \zeta }{\alpha} \cdot \eps}
                        + \eps\,.
                \]
                Further as $\zeta,\opt\leq \alpha\eps$ and $0<\alpha\leq 1$, we have $\nfrac{\abs{P\setminus H'}}{\abs{P}}
                    \leq 5\eps$ which ensures the feasibility of $H'$.
            \end{proof}
            Now, we are ready to upper bound the mass of $H\setminus \optset$ under $\cDgiven$.
            For this, we decompose $H\setminus \optset$
            \[
                H\setminus \optset 
                    = H - \optset + \inparen{\optset \setminus H}
                    = \inparen{H - H'} 
                        + \inparen{H' - \optset} 
                        + \inparen{\optset \setminus H}\,,
            \]
            so that 
            \[
                \cDgiven\inparen{H\setminus \optset}
                    = \inparen{
                        \cDgiven\inparen{H} -    \cDgiven\inparen{H'}
                      } + 
                      \inparen{
                        \cDgiven\inparen{H'} - \cDgiven\inparen{\optset}
                      } +
                      \cDgiven\inparen{{\optset \setminus H}}\,.
                    \yesnum\label{eq:sampleComplexity:decomposition}
            \]
            The bulk of the remaining proof will focus on bounding the first term; before that, we turn to upper bounding the second and third terms. 
            \begin{lemma}
                \label{lem:sampleComplexity:Hard:UBterms23}
                Conditioned on event $\evE$ from \cref{lem:sampleComplexity:Hard:feasibility}, the following bounds hold
                \[
                    \cD(H') - \cD(\optset) \leq \frac{1}{\sigma}\cdot \inparen{2\eps}^{\sfrac{1}{q}} 
                    \qquadand
                    \cD(\optset\setminus H)
                    \leq \frac{1}{\sigma}\cdot\inparen{3\eps}^{1/q}\,.
                \]
            \end{lemma}
            \begin{proof}[Proof of \cref{lem:sampleComplexity:Hard:UBterms23}]
                As $\cDgiven\sinparen{H'} - \cDgiven\inparen{\optset}\leq \cDgiven\sinparen{H'\triangle \optset}$, \cref{asmp:smoothness} implies that 
                \[
                    \cDgiven\sinparen{H'} - \cDgiven\inparen{\optset}
                    \leq \cDgiven\sinparen{H'\triangle \optset}
                    \leq \frac{1}{\sigma}\cdot \cDtrue\sinparen{H'\triangle \optset}^{\sfrac{1}{q}}
                    \,.
                \]
                By the definition of $H'$ (\cref{eq:perm:closeToInf}), 
                \[
                    \cDgiven\inparen{H'} - \cDgiven\inparen{\optset}
                    \leq \frac{1}{\sigma} \inparen{\opt+{ \zeta }}^{\sfrac{1}{q}}
                    \quad\Stackrel{\opt,~\zeta\leq \eps}{\leq}\quad \frac{1}{\sigma}\inparen{2\eps}^{1/q}\,.
                \]
                Next, since $H$ is feasible for \pERM{} with tolerance $\rho=\gamma$, \cref{lem:perm:easy} implies that, conditioned on $\evE$, $\cDtrue(\optset\setminus H)\leq \gamma +\sqrt{\gamma\eps}+\eps$.
                Since $\gamma\leq \eps$, we deduce $\cDtrue(\optset\setminus H)\leq 3\eps$.
                Applying \cref{asmp:smoothness} implies the result.
            \end{proof} 
            Next, we upper bound the first term in \cref{eq:sampleComplexity:decomposition}.
            \begin{lemma}
                \label{lem:sampleComplexity:Hard:UBterm1}
                Let $\evF$ be the event that the uniform convergence bounds in \cref{thm:unifConvergence} hold for both the positive samples $P$ and the unlabeled samples $U$, and the hypothesis classes $\hyH\cup\inbrace{G^c\colon G\in \hyH}\cup \inbrace{G\setminus (\optset \cap H\cap H')\colon G\in \hyH}$.
                Then, $\Pr[\evF]\geq 1-(\nfrac{\delta}{4})$.
                Moreover, conditioned on $\evF$, 
                \[
                    \cD(H) - \cD(H')
                    \leq \frac{11}{\sigma} \, (9\eps )^{1/q}
                    \,.
                \]
            \end{lemma}
            Note that \cref{lem:sampleComplexity:Hard:UBterms23,lem:sampleComplexity:Hard:UBterm1} imply that 
            \[
                \cDgiven\inparen{H\setminus \optset}
                    \leq 
                    \frac{1}{\sigma}\,(2\eps)^{1/q}
                    + \frac{1}{\sigma}\,(3\eps)^{1/q}
                    + \frac{11}{\sigma} \, (9\eps )^{1/q}
                    \leq \frac{13}{\sigma} \, (9\eps )^{1/q}\,.
            \]
            Now, \cref{lem:perm:hard} follows by using \cref{asmp:smoothness}.

            It remains to prove \cref{lem:sampleComplexity:Hard:UBterm1}.
            The key idea is that $H$ is $\eps$-optimal for \pERM{} with $\rho=5\eps$ while $H'$ is feasible (\cref{lem:sampleComplexity:Hard:feasibility}), and, hence,
            \[
                \frac{\abs{H\cap U}}{\abs{U}}
                \leq 
                \frac{\abs{H'\cap U}}{\abs{U}}\,.
                \yesnum\label{eq:sampleComplexity:optimalityOfH}
            \]
            Now, one can control $\cD(H)-\cD(H')$ by combining uniform convergence and the above result.
            However, this requires $O(\nfrac{1}{\eps^2})$ samples since the error of $H'$, $\cD(H')$, may not be small. 
            To get the optimal sample complexity of $O(\nfrac{1}{\eps})$, we use the fact that both $H$ and $H'$ are close to $\optset$, \ie{}, their errors are small for the labeling $\mathds{1}\sinbrace{x\in \optset},$ and then use the second-order uniform convergence bounds (\cref{thm:unifConvergence}).
            \begin{proof}[Proof of \cref{lem:sampleComplexity:Hard:UBterm1}]
                The lower bound on $\Pr[\evF]$ follows from \cref{thm:unifConvergence}, the choice of the sample size $n$, and the fact that $\vc{(\inbrace{G^c\colon G\in \hyH})}\leq \vc{(\hyH)}$ and $\vc{(\inbrace{G\setminus (\optset \cap H\cap H')\colon G\in \hyH})}\leq \vc{(\hyH)}$\footnote{To see this, fix any sets $A,B,C.$
                    Let \(\hyF=\{G\setminus (A\cap B\cap C):G\in \hyH\}\). For any finite set of samples \(X\), define the labelings \(L_{\hyH}=\{X\cap G: G\in \hyH\}\) and \(L_{\hyF}=\{X\cap (G\setminus (A\cap B\cap C)):G\in \hyH\}\). 
                    Note that $\{X\cap (G\setminus (A\cap B\cap C)):G\in \hyH\}=\{(X\cap G)\setminus (A\cap B\cap C)):G\in{\hyH}\}=\{L\setminus (A\cap B\cap C):L\in L_{\hyH}\}$
                    The map \(g\colon L_{\hyH}\to L_{\hyF}\) given by \(g(L)=L\setminus (A\cap B\cap C)\) is injective, so \(|L_{\hyF}|\le |L_{\hyH}|\).} and, hence,
                \[
                    \vc{\bbinparen{\hyH\cup\inbrace{G^c\colon G\in \hyH}\cup \inbrace{G\setminus (\optset \cap H\cap H')\colon G\in \hyH}}}
                    \leq O(\vc{(\hyH)})
                    \,.
                \] 
                Define 
                \[
                    T = \optset \cap H\cap H'\,.
                \]
                Since
                \[
                    \cD(H) = \cD(T) + \cD(H\setminus T)
                    \qquadand
                    \cD(H') = \cD(T) + \cD(
                    H'\setminus T)\,,
                \]
                we have
                \[
                    \cD(H)-\cD(H')
                    = 
                    \cD(H\setminus T)-\cD(H'\setminus T)
                    \leq \cD(H\setminus T)\,.
                    \yesnum\label{eq:sampleComplexity:UBterm1:ub}
                \]
                We now bound $\cD(H\setminus T)$ in two steps.

                \paragraph{Step 1 (Upper bound on $\nfrac{\abs{U\setminus T}}{\abs{U}}$).}
                    Notice that 
                    \begin{align*}
                        \cD(H'\setminus T)
                        = \cD(H'\setminus \optset)
                        + \cD\inparen{\sinparen{H'\cap \optset}\setminus T}\,.
                    \end{align*}
                    \cref{eq:perm:closeToInf} and $\opt,~\zeta\leq \eps$ imply that $\cD(H'\setminus \optset)\leq 2\eps$.
                    Therefore, 
                    \[
                        \cD(H'\setminus T)
                        \leq 2\eps
                        + \cD\inparen{\sinparen{H'\cap \optset}\setminus T}\,.
                    \]
                    Further, since $\optset\setminus T \subseteq \inparen{\optset \setminus H'} \cup \inparen{\optset\setminus H}$,
                    \[
                        \cD\inparen{\sinparen{H'\cap \optset}\setminus T}\leq \cD(\optset \setminus T)\leq \cD(\optset\setminus H')+ \cD(\optset\setminus H)\,.
                    \]
                    Therefore, 
                    \begin{align*}
                        \cD(H'\setminus T)
                        \leq 2\eps 
                        + \cD(\optset\setminus H')
                        + \cD(\optset\setminus H)\,.
                        \yesnum\label{eq:sampleComplexity:UBterm3:step1}
                    \end{align*}
                    We have the following upper bounds on the last two terms:
                    \begin{enumerate}
                        \item Since $H$ is feasible for \pERM{} with $\rho=\gamma\leq \eps$, conditioned on event $\evF$, \cref{lem:sampleComplexity:Hard:UBterms23}, implies that $\cD(\optset\setminus H)\leq (\nfrac{1}{\sigma})\,\sinparen{3\eps}^{1/q}$.
                        (Here, we use the fact that $\evF$ implies $\evE$ by construction.)
                        \item Conditioned on $\evF$ (which implies $\evE$), $H'$ is feasible for \pERM{} with $\rho=5\eps$.
                        Hence, \cref{lem:perm:easy} implies that $\cDtrue(\optset\setminus H')\leq 5\eps + \sqrt{5}\cdot \eps+\eps\leq 9\eps $.
                        Finally, \cref{asmp:smoothness} implies that $\cD(\optset\setminus H')\leq (\nfrac{1}{\sigma})\,(9\eps )^{1/q}$.
                    \end{enumerate}
                    Condition only event $\evF$.
                    Substituting the above two bounds into \cref{eq:sampleComplexity:UBterm3:step1}, implies that 
                    \[
                        \cD(H'\setminus T) \leq 2\eps + \frac{1}{\sigma}\,(3\eps)^{1/q} + \frac{1}{\sigma}\,(9\eps )^{1/q}
                        ~~~\Stackrel{}{\leq}~~~ 2\eps + \frac{2}{\sigma}\,(9\eps )^{1/q}
                        ~~~\Stackrel{q\geq 1}{\leq}~~~
                        \frac{3}{\sigma}\,(9\eps )^{1/q}\,.
                    \]
                    Now, uniform convergence over the unlabeled samples and hypotheses class $\inbrace{G\setminus T\colon G\in \hyH}$ implies that,
                    \[
                        \frac{\abs{(H'\setminus T)\cap U}}{\abs{U}}
                        \leq 
                        \frac{3}{\sigma}\,(9\eps )^{1/q}
                        + 
                        \sqrt{\frac{3}{\sigma}\,(9\eps )^{1/q}\cdot \eps}
                        +
                        \eps 
                        =
                        \frac{3}{\sigma}\, (9\eps )^{1/q}
                        + 
                        \sqrt{\frac{3}{\sigma}}\,(9\eps )^{\frac{1+q}{2q}}
                        +\eps\,.
                    \]
                    Since $\nfrac{1}{q}\leq \nfrac{(1+q)}{2q}$ (as $q\geq 1$), $9\eps \leq 1$ (which can be ensured by rescaling $\eps$ by a constant factor), and $\sigma\leq \sqrt{\sigma}$, it follows that 
                    \[
                        \frac{\abs{(H'\setminus T)\cap U}}{\abs{U}}
                        \leq 
                        \frac{3+\sqrt{3}+1}{\sigma}\, (9\eps )^{1/q}
                        \leq  \frac{6}{\sigma}\, (9\eps )^{1/q}\,.
                        \yesnum\label{eq:sampleComplexity:boundSamples}
                    \]
                \paragraph{Step 2 (Upper bound on $\cD(H\setminus T)$).}
                    \cref{eq:sampleComplexity:optimalityOfH} and the fact that for any set $S$, $\nfrac{\abs{(S\setminus T)\cap U}}{\abs{U}}=\nfrac{\abs{S\cap (U\setminus T)}}{\abs{U}}$ implies that 
                    \[
                        \frac{\abs{(H\setminus T)\cap U}}{\abs{U}}
                        \leq 
                        \frac{\abs{(H'\setminus T)\cap U}}{\abs{U}}
                        +\eps\,.
                    \]
                    Hence, \cref{eq:sampleComplexity:boundSamples} implies that,
                    \[
                        \frac{\abs{(H\setminus T)\cap U}}{\abs{U}}
                        \leq  \frac{7}{\sigma}(9\eps )^{1/q}\,.
                    \]
                    Now, uniform convergence over $U$ and the hypothesis $\inbrace{G\setminus T\colon G\in \hyH}$ implies that,
                    \begin{align*}
                        \cD(H\setminus T)
                        \leq 
                            \frac{7}{\sigma}\,(9\eps )^{1/q}
                            + \sqrt{\frac{7}{\sigma}\,(9\eps )^{1/q} \cdot \eps}
                            + \eps
                            = 
                                \frac{7}{\sigma}\,(9\eps )^{1/q}
                                +
                                \sqrt{\frac{7}{\sigma}}\,(9\eps )^{\frac{1+q}{2q}}
                                +\eps \,.
                    \end{align*}
                    Noting again that $\sigma\leq \sqrt{\sigma}\leq 1$ and that the square root term can be absorbed as $\nfrac{1}{q}\leq \nfrac{(1+q)}{2q}$, we obtain
                    \[
                            \cD(H\setminus T)
                            \leq \frac{7+\sqrt{7}+1}{\sigma}\, (9\eps )^{1/q}
                            \leq \frac{11}{\sigma} \, (9\eps )^{1/q}
                            \,.
                    \]
                    Substituting this into \cref{eq:sampleComplexity:UBterm1:ub} completes the proof.
            \end{proof}


\section{Proof of \cref{thm:impossibilityPositiveSamples} (Impossibility of Learning from Positive Samples)}\label{sec:asmp:examples}  
    \label{sec:lowerBound:positiveSamples}
    In this section, we prove \cref{thm:impossibilityPositiveSamples}, which shows that most of the interesting hypothesis classes are not learnable from positive samples when the distribution of covariates is allowed to be arbitrary.
    This starkly contrasts our positive result in \cref{sec:intro:application:smooth}, which shows that if the distribution of covariates is smooth, then all finite VC dimension classes are learnable from positive examples.

    \impossiblityPositiveSampleLearning*
    \noindent 
    To see that most common hypothesis classes satisfy the above requirement and, hence, are not learnable from positive samples alone, note that intersections of $k$ one-dimensional intervals or intersections of $k$ two-dimensional halfspaces have a VC dimension $\Omega(k)$, which tends to $\infty$ with $k$.
    Moreover, this also implies that any hypothesis class containing one-dimensional intervals or two-dimensional halfspaces (\eg{}, decision trees) are also not learnable from positive examples.
    
    An analogous result is proved by \citet{natarajan1987learning} in the special care where the learner is required to be proper.
    To the best of our knowledge, the extension to improper learners is novel.

     
    \noindent 
    \begin{proof}[Proof of \cref{thm:impossibilityPositiveSamples}]
        Consider any $k\geq 1$.
        We reduce the problem of realizable $(\eps,\delta)$-PAC learning {$\hyH_{\cap k}$} from (labeled) samples $(x,y)$ where $x\sim \cDtrue$ to $(\eps,\delta)$-PAC learning {$\hyH$} from \textit{only} positive examples under the assumptions in the theorem.
        
        This reduction will work when the optimal hypothesis $\optset$ has sufficient mass in given PAC learning instance, in particular, $\cDtrue(\optset)\geq \nfrac{1}{2}$.
        The desired result can be deduced from this reduction due to the following sample complexity lower bound for learning intersections of $k$ halfspaces by taking the limit $k\to\infty$.
        \begin{theorem}[\protect{\cite[Theorem 5.3]{neuralNetwork2009bartlett}}]
                \label{thm:PAClearningLowerbound}
                The sample complexity of $(\eps,\delta)$-PAC learning {$\hyH_{\cap k}$} under the assumption that $\cD(\optset)\geq \nfrac{1}{2}$ is $S_{\eps,\delta}(k)\coloneqq \wt{\Omega}\inparen{(\nfrac{1}{\eps})(\vc{}(\hyH_{\cap k})+\log{\nfrac{1}{\delta}})}$ for $\eps<\nfrac{1}{8}$ and $\delta < \nfrac{1}{100}$.
                {In particular, if $\lim_{k\to \infty} \vc{}(\hyH_{\cap k})=\infty$, then $\lim_{k\to \infty} S_{\eps,\delta}(k)=\infty$.}
        \end{theorem}
        \begin{proof}[Proof sketch of \cref{thm:PAClearningLowerbound}]
                The proof follows by verifying that the hard instance created in Theorem 5.3 of \cite{neuralNetwork2009bartlett} satisfies $\cD(\optset)\geq \nfrac{1}{2}$ and that the VC-dimension of {$\hyH_{\cap k}$ goes to $\infty$ with $k$ by definition}. 
            \end{proof}
        

        \noindent In the remainder of the proof, we describe the reduction and prove its correctness.

        \paragraph{Reduction.}
        Fix an instance of PAC-learning {$\hyH_{\cap k}$}, covariate distribution $\cDtrue$, and the optimal set $\optset$ satisfying $\cDtrue(\optset)\geq \nfrac{1}{2}$.
        The reduction is as follows.
        \begin{mdframed}
            \begin{enumerate}[itemsep=-1pt]
                \item[]  \hspace{-6mm} \textbf{Input:} 
                    An accuracy parameter $\eps>0$, a confidence parameter $\delta>0$, and
                    $n$ labeled samples $\inbrace{\inparen{x_i,y_i}\colon 1\leq i\leq n}$ where $x_1,x_2,\dots,x_n$ are \iid{} from $\cD$ and $y_i=\mathds{1}\sinbrace{x_i\in \optset}$
                \item[]  \hspace{-6mm}  \textbf{Oracle Access:}
                    Query access to a PAC learner $\cL$ that learns from only positive samples 
                \vspace{2mm}
                \item Let the set of positive examples be $P=\inbrace{x_i \colon i\in \inbrace{1, 2, \dots,n},~~~ y_i=1}$.
                \item Return the hypothesis $\wh{H}$ obtained by querying the learner $\cL$ with $\inparen{\eps,\delta,P}$.
            \end{enumerate}
        \end{mdframed} 
        Since $\cDtrue(\optset)\geq \nfrac{1}{2}$, we get that with high probability $\abs{P}\geq \Omega(n)$ and, hence, we only lose a constant factor of samples in the reduction.
        
        
        \paragraph{Soundness and Correctness.}
        To prove the correctness of the reduction, suppose $\wh{H}$ satisfies the requirement of PAC learning from positive examples:
            with probability $1-\delta$, 
            for each distribution $\cD'$, 
                satisfying $\cD'(\optset)\geq \nfrac{1}{2}$ and
                $\cD'_{\optset}=\cP$, it holds that 
            \[
                \cD'(\wh{H}\triangle \optset)\leq\eps\,.
                \yesnum\label{eq:lowerbound:PIUGuarantee}
            \]
        We will show that it implies $\cDtrue\sinparen{\wh{H}\triangle \optset}\leq 2\eps$, proving the result.
        Let the optimal hypothesis be $\optset=H^\star_1\cap H^\star_2\cap \dots \cap H^\star_k$ where $H^\star_1,H^\star_2,\dots,H^\star_k \in \hyH$.
        Without loss of generality, suppose that there is no set $S\subsetneq[k]$ with $\optset\neq \bigcap_{i\in S} H^\star_i$.
        Otherwise, we can select the smallest such $S$, and continue with the remainder of the proof with $k=\abs{S}$ and the halfspaces $\inbrace{H^\star_i \colon i\in S}$.
        Define the following sets, which are useful in the analysis:
        \[
            T_1 = \R^2 \setminus H^\star_1\,,\quad 
            T_2 = H^\star_1 \setminus H^\star_2\,,\quad 
            T_3 = \sinparen{H^\star_1\cap H^\star_2} \setminus H^\star_3\,,\quad \dots\quad 
            T_k = {\bigcap_{i=1}^{k-1}H^\star_i} \setminus H^\star_k
            \,.
        \]
        We will only need the following properties of $T_1,T_2,\dots,T_k$: for distinct $1\leq i,j\leq k$
        \[
            T_i \neq \emptyset\,,\quad 
            T_i \cap T_j = \emptyset\,,\quad 
            T_i \cap \optset = \emptyset\,,\quadand
            T_i \subseteq \R^2\setminus H^\star_i\,.
            \yesnum\label{eq:onesided:propertiesofT}
        \]
        The first condition follows from the assumption that there is no set $S\subsetneq [k]$ with $\optset=\bigcap_{i\in S}H^\star_i$.
        The remaining conditions follow from construction.
        Consider the following sequence of distributions: for each $1\leq i\leq k,$
        \[
            \cD'_i = \frac{1}{2}
                        \cDtrue_{\optset} + \frac{1}{2}\cDtrue_{T_i}\,,
        \]
        where, for each set $S$ and distribution $\cQ$, $\cQ_S$ is the truncation of $\cQ$ to $S$.
        First, we verify that $\cD'_i$ satisfies the guarantees specified above.
        \begin{fact}
            For each $1\leq i\leq k$, 
            $\cD'(\optset)\geq \nfrac{1}{2}$, and 
            $\cD'_{\optset}=\cP$. %
        \end{fact}
        \begin{proof}
            Fix any $1\leq i\leq k$.
            It holds that, $\cD'(\optset)=\sfrac{\cDtrue_{\optset}(\optset)}{2}=\nfrac{1}{2}$.
            Further, by construction, $\cD'_{\optset}=\cP$.






            
            
            
            
        \end{proof}
        Now, using \cref{eq:lowerbound:PIUGuarantee} for each of $\cD'_1,\dots,\cD'_k$ implies that: for each $1\leq i\leq k$,
        \[
            \cDtrue_{\optset}(\wh{H}\triangle \optset)
            + 
            \cDtrue_{T_i}(\wh{H}\triangle \optset)
            \leq 2\eps\,.
        \]
        Hence, 
        \[
            \cDtrue_{\optset}(\wh{H}\triangle \optset)\leq 2\eps
            \qquadand
            \max_{1\leq i\leq k}\cDtrue_{T_i}(\wh{H}\triangle \optset)
            \leq 2\eps\,.
            \yesnum\label{eq:onesided:errorGuarantee}
        \]
        Now, we are ready to bound $\cDtrue(\wh{H}\triangle \optset)\leq O(\eps)$:
        \begin{align*}
            \cDtrue(\wh{H}\triangle \optset)
            &= \cDtrue(\optset) \cDtrue_{\optset}(\wh{H}\triangle \optset)
            + \sum_i \cDtrue(T_i) \cDtrue_{T_i}(\wh{H}\triangle \optset)
                \tag{since $T_1,T_2,\dots,T_k,\optset$ are disjoint; \eqref{eq:onesided:propertiesofT}}\\
            &\leq 2\cDtrue(\optset)\eps
            + 2\eps \sum_i \cDtrue(T_i)
                \tag{by \eqref{eq:onesided:errorGuarantee}}\\
            &= 2\eps\,.
        \end{align*}
        
    \end{proof}



\section{Proofs Deferred from the Main Body} %
    \subsection{Proof of \cref{lem:constReg:unifConvergenceHoldsWHP} (Uniform Convergence in Each Repetition)}
                \label{sec:proofof:lem:constReg:unifConvergenceHoldsWHP}
                Recall that \cref{lem:constReg:unifConvergenceHoldsWHP} requires us to show that the following events hold with probability $1-\delta$:
                \begin{itemize}
                    \item For every repetition $1 \leq i \leq T$ of Subroutine B of \cref{alg:constReg}, the samples in each repetition satisfy $\zeta$-uniform convergence with respect to degree-$k$ PTFs, and
                \end{itemize}
                The proof follows from standard uniform convergence arguments and a union bound.
                \begin{proof}
                    Recall that the VC-dimensions of degree-$k$ polynomials is $\wt{O}(d^{k})$.
                    Therefore, for any $\eta\in (0,1)$, 
                    \[
                        \wt{O}\inparen{
                            \frac{1}{\zeta^2}
                            \cdot \inparen{d^k + \log{\frac{1}{\eta}}}
                        }\,.
                    \]
                    samples are sufficient to achieve $\zeta$-uniform convergence with respect to degree-$k$ polynomials with probability $1-\eta$ for any single repetition of Subroutine B.
                    We set $\eta=\nfrac{\delta}{T}$ and take a union bound over all $T$ events, which are all the $T$ independent repetitions of Subroutine B, to get that if $n$ is at least 
                    \[
                        \wt{O}\inparen{
                        \frac{1}{\zeta^2}
                        \cdot \inparen{d^k + \log{\frac{T}{\delta}}}
                        }\,,
                    \]
                    then the result follows.
                    Now, substituting the value of $T$, namely, $T=O\inparen{\inparen{\nfrac{1}{\sqrt{\zeta}}}\cdot \log{\nfrac{1}{\delta}}}$, and simplifying implies that the required number of samples is 
                    \[
                         \wt{O}\inparen{
                         \frac{1}{\zeta^2}
                            \cdot \inparen{d^k + \log{\frac{1}{\delta}}}
                        }\,.
                    \]
                    The choice of $n$ in \cref{alg:constReg} satisfies this.
                     
                \end{proof}
    
            
    \subsection{Proof of \cref{lem:constReg:boost} (Boosting via Repetition)}\label{sec:proofof:lem:constReg:boost}
                \begin{proof}[Proof of \cref{lem:constReg:boost}]
                Fix $\eta$ corresponding to the guarantee in \cref{lem:constReg:expectedValue}
                \[
                    \eta= O\inparen{\sqrt{\frac{1}{\alpha\sigma}}} \cdot \zeta^{1/(4q)}
                        + O\sinparen{\delta}\,.
                \]
                Observe that for each $i$, $\nfrac{\abs{H_i\cap U}}{\abs{U}}\geq 0$ and, hence, Markov's inequality implies that 
                \[
                    \Pr\insquare{\frac{\abs{H_i\cap U_i}}{\abs{U_i}} 
                    \geq \cD(H_{\opt})+2\eta} 
                    \leq {\frac{\cD(H_{\opt})+\eta}{\cD(H_{\opt})+2\eta}}\,.
                \]
                Hence, since each repetition of Subroutines A and B in \cref{alg:constReg} is independent, it follows that 
                \[
                    \Pr\insquare{
                        \min_{1\leq i\leq T}
                    \frac{\abs{H_i\cap U_i}}{\abs{U_i}} \geq \cD(H_{\opt})+2\eta 
                    } 
                    \leq \inparen{{\frac{\cD(H_{\opt})+\eta}{\cD(H_{\opt})+2\eta}}}^T
                    ~~~~~~\Stackrel{\substack{\cD(H_{\opt})\leq 1,\\0\leq\eta\leq \nfrac{1}{2}}}{\leq}~~~~~~ 
                        \inparen{1-\frac{\eta}{2}}^T
                    \leq e^{-\frac{1}{2}\eta T}
                    \leq \delta
                    \,.
                \]
                Where in the last inequality, we used $T=O\inparen{\inparen{\nfrac{1}{\sqrt{\zeta}}}\log\inparen{\nfrac{1}{\delta}}}$ and, hence, $\eta T\geq 2\log{\nfrac{1}{\delta}}$.
                Since with probability $1-\delta$, $\zeta$-uniform convergence holds with respect to each $U_1,U_2,\dots,U_T$ and hypothesis $\hyH$ (see \cref{lem:constReg:unifConvergenceHoldsWHP}), it follows that with probability $1-2\delta$,
                \[
                    \Pr\insquare{
                    \min_{1\leq i\leq T} \cD(H_i) \geq \cD(H_{\opt}) + 2\eta + \zeta
                    }
                    \leq \delta\,.
                \]
                Define $\wh{H}=H_{i^\star}$ for $i^\star\in \argmin_{1\leq i\leq T} \abs{H_i\cap U'}$.
                Now, uniform convergence, with respect to $\inbrace{H_1,\dots,H_N}$ and $U'$ implies that with probability $1-\delta$, $\abs{\cDgiven(H_i) - \inparen{\nfrac{\abs{H_i\cap U_i'}}{\abs{U_i'}}}}\leq \zeta$ and, hence, with probability $1-3\delta$,
                \[
                    \cD(\wh{H})
                    \leq \frac{\sabs{\wh{H}\cap U'}}{\abs{U'}} + \zeta
                    \leq \cD(H_{\opt}) + 2\eta + 3\zeta 
                    ~~\Stackrel{\zeta\leq \eta}{\leq}~~ \cDgiven(H_{\opt}) + 5\eta\,. 
                \] 
            \end{proof} 
     
        
            
        
    \subsection{Framing \cref{eq:constReg:prog} in \cref{alg:constReg} as a Linear Program}\label{sec:module:linearProgram}
            In this section, we show how to formulate the optimization program in \cref{alg:constReg} as a linear program. 
            Recall that \cref{eq:constReg:prog} is as follows.
            \[
                \min_{{\rm deg}(p)\leq k}~\frac{\sum_{x\in U} \abs{{p(x)}}}{\abs{U}}
                    \,,\quad 
                    \text{such that}\,,\quad 
                    {\frac{\sum_{x\in P} \min\inbrace{p(x), 1}}{\abs{P}} \geq {1-\rho}} \,.
            \]
            Let $\hypo{R}$ be the set of all multisets $R\subseteq [d]$ of size at most $k$.
            For each $z\in \R^d$ and $R\in \hypo{R}$, define the monomial $z(R)\coloneqq \Pi_{i\in R} z_i$.
            For each monomial $z(R)$, let $c(R)$ be the corresponding coefficient in a polynomial $p$.
            The above program is equivalent to the following linear program.
            \begin{align*} 
                \min_{c\in \R^{\hypo{R}},\ z\in \R^{\abs{U}},\ w\in \R^{\abs{P}}}~
                    & {\sum_{x\in U} z(x)} \,,\\
                \forall_{x\in U}\,,\quad& \sum_{R\in \hypo{R}} c(R) x(R) ~\leq~ z(x) ~~\leq~ -\sum_{R\in \hypo{R}} c(R) x(R)\,,\\
                \forall_{x\in P}\,,\quad& \sum_{R\in \hypo{R}} c(R) x(R)~\leq~ w(x) ~\leq~ ~ 1\,,\\
                & {\frac{\sum_{x\in P} w(x)}{\abs{P}} \geq {1-\rho}}\,.
            \end{align*} 
    
    
    \subsection{\cref{asmp:smoothness} Is Weaker than Bounding $\chi^2$--Square-Divergence}\label{sec:distCloseness}
        In this section, we show that if $\chidiv{\cDtrue}{\cD}<\infty$, then \cref{asmp:smoothness} holds. %
        \begin{lemma}
            \label{lem:closenessFromChiSquare}
            For $B\geq 1$, if $\chidiv{\cDtrue}{\cDgiven}\leq B$, then \cref{asmp:smoothness} holds with $\sigma=\frac{1}{\sqrt{1+B}}$ and $q=2$.
        \end{lemma}
        \begin{proof}
            Fix measurable set $S$.
            Change of measure and the Cauchy--Schwarz inequality imply that
            \[
                \cDtrue(S) 
                = \int \mathds{1}_S(x) \frac{\cDtrue(x)}{\cD(x)}\d\cD(x)
                \leq 
                    \inparen{
                        \int \mathds{1}_S(x)^{2}~ \d\cD(x)
                    }^{1/2}
                    \inparen{
                        \int \inparen{\frac{\cDtrue(x)}{\cD(x)}}^2 \d \cD(x)
                    }^{1/2}\,.
            \]
        Since $\mathds{1}_S(x)^{2}=\mathds{1}_S(x)$ and $
                \chidiv{\cDtrue}{\cD}= 
            {\int \inparen{\nfrac{\d \cDtrue(x)}{\d \cD(x)}}^2 \d \cD(x)} - 1$,
        the above is the same as 
        \[
            \cDtrue(S) \leq \cD(S)^{1/2}\cdot 
                \inparen{\chidiv{\cDtrue}{\cD}+1}^{1/2}
                \leq 
                    \sqrt{1+B}\cdot 
                    \cD(S)^{1/2}
                \,.
        \]
        The result follows since this holds for all measurable sets $S$.
        \end{proof}
        More generally, generalizing the above proof (by using Hölder's inequality instead of the Cauchy--Schwarz inequality) shows that \cref{asmp:smoothness} also follows from bounds on Rényi divergences:
        \begin{lemma}
            \label{lem:closenessFromRenyi}
            For $B,r\geq 1$, if $\renyi{r}{\cDtrue}{\cDgiven}\leq B$, then \cref{asmp:smoothness} holds with $\sigma=e^{-B/q}$ and $q=\frac{r}{r-1}$.
        \end{lemma}
        Note that the first result follows from the second one when $r=2$.
        To see this, recall that $\renyi{2}{\cDtrue}{\cD} = \log\sinparen{1+\chidiv{\cDtrue}{\cD}}$, so $\chidiv{\cDtrue}{\cD}\le B$ implies $\renyi{2}{\cDtrue}{\cD}\le\log(1+B)$. 
        Applying \cref{lem:closenessFromRenyi} with $r=2$ (thus, $q=2$) gives $\cDtrue(S) \le \cD(S)^{1/2}\,e^{\frac{1}{2}\renyi{2}{\cDtrue}{\cD}} \le \cD(S)^{1/2}\,\sqrt{1+B}$, which is precisely the conclusion of \cref{lem:closenessFromChiSquare}.
        
    
            
        
    

    \subsection{Satisfying \cref{eq:mostGeneral:smoothness} via Bounds on KL-Divergence}
        \label{sec:klBoundsSatisfyGeneralSmoothness}
        Next, we show that if $\kl{\cDtrue}{\cD}<\infty$, then \cref{eq:mostGeneral:smoothness}, a weakening of \cref{asmp:smoothness}, holds.

        We restate \cref{eq:mostGeneral:smoothness} here  for convenience:
            there is a known rate function $r\colon \R_{\geq 0}\to \R_{\geq 0}$ satisfying $\lim_{m\to 0^+} r(m)=0$ such that the pair of distributions $\inparen{\cD, \cDtrue}$ satisfy
            \[
                \text{for any measurable set $S$}\,,\quad 
                r(\cDtrue(S))\leq \cD(S)\,.
                \tag{4}
            \] 
        \vspace{-5mm}
        \begin{theorem}\label{thm:capturingKL}
            
            If $\kl{\cDtrue}{\cD} = C$, then \eqref{eq:mostGeneral:smoothness} holds with $r(m)=e^{\sfrac{-3\cdot \max\inbrace{1, C}}{m}}$; satisfying \mbox{$\lim_{m\to 0^+}r(m)=0$.}
        \end{theorem}
        In contrast, $\kl{\cDtrue}{\cD}<\infty$ is insufficient to imply \cref{asmp:smoothness} for any finite $\sigma$ and $q$.
        Hence, the weakening of \cref{asmp:smoothness} in \cref{eq:mostGeneral:smoothness} is necessary to capture $\kl{\cDtrue}{\cD}<\infty$.

        To prove \cref{thm:capturingKL}, we will use the following inequality, which we believe is well-known but since we could not find an explicit reference, we prove it at the end of this section.
        \begin{lemma}\label{lem:kl_transfer}
            Let $\cP$ and $\cQ$ be probability distributions on $\R^d$ and assume $\kl{\cQ}{\cP} < \infty$. 
            Then, for any measurable set $S$,
            \[
                \cQ(S) \;\le\; 
                \inf_{t>1} \Bigl\{t\,P(S) \;+\; \frac{\kl{\cQ}{\cP}}{\ln t}\Bigr\}\,.
            \]
        \end{lemma}
        Now we are ready to prove \cref{thm:capturingKL}.
        \begin{proof}[Proof of \cref{thm:capturingKL}]
            First, \cref{lem:kl_transfer} implies that
            \[
            \cDtrue(S)\le\inf_{t>1}\Bigl\{t\,\cD(S)+\frac{C}{\ln t}\Bigr\}\,.
            \]
            Choosing $t=\nfrac{1}{\sqrt{\cD(S)}}$ and using that $\sqrt{z}\leq \frac{1}{\ln(1/z)}$ for all $z>0$, gives
            \[
            \cDtrue(S)\le\sqrt{\cD(S)}+\frac{2C}{\ln(1/\cD(S))}\le\frac{3\max\inbrace{1, C}}{\ln(1/\cD(S))}\,.
            \]
            which is equivalent to 
            \[
                \cD(S)\geq \exp\Bigl(-\frac{3 \max\inbrace{1, C}}{\cDtrue(S)}\Bigr)\,,
            \]
            that satisfies \eqref{eq:mostGeneral:smoothness} with $r(m)=\exp(\nfrac{-3\max\inbrace{1, C}}{m})$ and $\lim_{m\to0^+}r(m)=0$.
        \end{proof}
        \begin{proof}[Proof of \cref{lem:kl_transfer}]
        Fix any $t>1$, and define the set
        \[
        T_t \;=\; \Bigl\{\, x \in \R^d : \frac{\cQ(x)}{\cP(x)} \;>\; t \Bigr\}.
        \]
        On $T_t$, we have $\ln\inparen{\sfrac{\cQ(x)}{\cP(x)}} > \ln t$. Hence,
        \[
        \kl{\cQ}{\cP} \;=\; \int \cQ(x)\,\ln{\frac{\cQ(x)}{\cP(x)}}\,\d x
        \;\ge\; \int_{T_t} \cQ(x)\,\ln{\frac{\cQ(x)}{\cP(x)}}\,\d x
        \;>\; (\ln t)\,Q(T_t)\,,
        \]
        which gives $Q(T_t) \le \kl{\cQ}{\cP}/\ln t$. Next, split
        \[
        Q(S) \;=\; Q(S \cap T_t) \;+\; Q(S \cap T_t^c)\,.
        \]
        On $S \cap T_t^c$, we have $\nfrac{\cQ(x)}{\cP(x)} \le t$, so
        \[
        Q(S \cap T_t^c)
        \;=\; \int_{S \cap T_t^c} \cQ(x)\,\d x
        \;\le\; \int_{S \cap T_t^c} t\,\cP(x)\,\d x
        \;=\; t\,P(S \cap T_t^c)
        \;\le\; t\,P(S)\,.
        \]
        Therefore,
        \[
        Q(S) \;\le\; Q(S \cap T_t) + Q(S \cap T_t^c)
        \;\le\; Q(T_t) + t\,P(S)
        \;\le\;
        \frac{\kl{\cQ}{\cP}}{\ln t} + t\,P(S)\,.
        \]
        Since $t>1$ was arbitrary, taking the infimum over $t>1$ completes the proof.
        \end{proof}


\subsection{Generating List of Distributions From Adversarially Corrupted Sources}\label{sec:advCorruption}
            In this section, we show how one can satisfy the List Decoding model's requirements (\cref{def:list}) in the following types of data sources -- where all but an $\gamma$-fraction of the samples can be \textit{adversarially} corrupted.

            \begin{mdframed}
                \textbf{Example (Adversarially Corrupted Gaussian Source):}\quad
                Suppose that $\cDtrue$ is a Gaussian distribution $\cN(\muStar,\SigmaStar)$.
                Consider a data source that is $(1-\gamma)$-corrupted (for some $\gamma\in (0,1]$):
                    with probability $\gamma$, the source outputs a sample from $\cDtrue$, and, otherwise, outputs an arbitrary sample that may depend on the past samples and the knowledge of the learning algorithm.
            \end{mdframed}
            We show that one can efficiently convert the above $(1-\gamma)$-corrupted source for any $\gamma>0$, into one satisfying the requirements of the List Decoding PIU Model with $\ell= \poly(\nfrac{1}{\gamma})$.

            In the remainder of this section, we will prove a generalization of this result where $\cDtrue$ is any exponential family distribution satisfying mild assumptions.
            The result for Gaussian $\cDtrue$ follows as a special case of this.
            
            

            \paragraph{Notation.}
            We begin by recalling the definition of an exponential family. %
            \begin{definition}[\textbf{Exponential Family}]
                Given $m\geq 1$, a function $h\colon \R^d\to \R_{\geq 0}$ called the \emph{carrier measure}, a function $t\colon \R^d\to \R^m$ called the \emph{sufficient statistic}, the exponential family $\cE_{m,h,t}$ is a set of distributions parameterized by a vector $\theta\in \R^m$ where, for each $\theta$, $\cE_{m,h,t}(\theta)\propto h(x) \cdot \exp\inparen{\theta^\top t(x)}$.
            \end{definition}
            Henceforth, $m$, $h$, and $t$ will be clear from convex and, hence, we drop them from the subscript. %
            The set of $\theta$ for which the density of $\cE(\theta)$ is well defined is called the \textit{natural parameter space} $\overline{\Theta}$.
            For each $\theta\in \overline{\Theta}$, define the \textit{log-partition function} $A(\theta)\coloneqq \log~{\int h(x) \cdot \exp\inparen{\theta^\top t(x)}\d x}$.
            For any canonical exponential family, $\overline{\Theta}$ is a convex set and $A$ is a convex function \cite{dasgupta2008asymptotic}.


            \paragraph{Assumptions.}
            Suppose $\cDtrue = \cE(\theta^\star)$ for some unknown parameter $\theta^\star$.
            We will be interested in a subset $\Theta\subseteq \overline{\Theta}$ where the Fisher Information Matrix $\cov_{\cE(\theta)}[t(x)]$ has bounded eigenvalues.
            Concretely, we make the following assumption on the family $\cE$. 
            \begin{assumption}\label{asmp:exponentialFamily:appendix}
                There is a subset $\Theta$ of the natural parameter space and $\Lambda\geq \lambda > 0$ and $\eta>0$ such that:
                \begin{enumerate}
                    \item \textbf{(Interiority)}           %
                        $\theta^\star$ is in the $\eta$-relative-interior $\Theta(\eta)$ of $\Theta$.
                    \item \textbf{(Projection Oracle)} There is an efficient projection oracle to the $\eta$-relative-interior of $\Theta$.
                    \item \textbf{(Bounded Fisher-Information)} For any $\theta\in \Theta$,
                    $\lambda I \preceq \cov_{x\sim \cE(\theta)}[t(x)]\preceq \Lambda I$.
                \end{enumerate}
            \end{assumption}
            \begin{remark}[\cref{asmp:exponentialFamily:appendix} Holds Unconditionally for Gaussians]
                Before proceeding, we note when $\cDtrue$ is a Gaussian distribution, then \cref{asmp:exponentialFamily:appendix} always holds.
                This is due to pre-processing routines by \citet{lee2024unknown} which show how to convert an arbitrary Gaussian distribution with parameter $\theta^\star=(\muStar,\SigmaStar)$ into one satisfying \cref{asmp:exponentialFamily:appendix}.
            \end{remark}
 
            \noindent In this section, we prove the following result.
            \begin{theorem}\label{thm:listDecoding:expFamily}
                Suppose \cref{asmp:exponentialFamily} holds with constants $\eta,\lambda,\Lambda>0$.
                There is an algorithm that, given $\gamma>0$ and $n={\inparen{\nfrac{m}{\gamma}}\log\inparen{\nfrac{1}{\delta}}}$ samples from $\cE(\theta^\star)$ out of which $(1-\gamma)$-fraction are adversarially corrupted, %
                outputs a list of $\inparen{\nfrac{2}{\gamma}}$-parameters $\theta_1,\theta_2,\dots,\theta_k$ such that, with probability at least $1-\delta$,
                for at least one $\theta_i$ ($1\leq i\leq k$), 
                \[
                    \renyi{w}{\cE(\theta^\star)}{\cE(\theta_i)}
                    \leq 
                    \frac{
                        \Lambda^4
                    }{
                        \gamma^3\lambda^3 
                    }
                    \quadwhere
                    w = 1 + \Omega\inparen{\frac{\gamma\lambda}{\Lambda}}
                    \,.
                \]
            \end{theorem}
            Now \cref{lem:closenessFromRenyi} implies that the above closeness in Rényi divergence implies \cref{asmp:smoothness} with $q=\nfrac{w}{(w-1)}=O\inparen{\nfrac{\Lambda}{(\gamma \lambda)}}$ and $\sigma=\exp(-\nfrac{\Lambda^3}{(\gamma^2\lambda^2)})$.
            %
            \begin{proof}[Proof of \cref{thm:listDecoding:expFamily}]
                The proof of this result uses techniques from \citet{charikar2017learning}.
            
                
                Fix $\beta= \nfrac{1}{2}$. %
                Let the (corrupted) samples be $X=\inbrace{x_1,x_2,\dots,x_n}$ and let $x_1,\dots,x_{r}$ be the set of \textit{uncorrupted} samples from $X$.
                For each $i$, define $f_i\colon \R^m\to \R^m$ to be the negative log-likelihood function, \ie{}, 
                $f_i(\theta) \coloneqq -\log\inparen{\cE(x_i; \theta)}$.
                Let $\overline{f}\colon \R^m \to \R^m$ be the population log-likelihood, \ie{}, $\overline{f}(\theta)=\Ex_{z\sim \cE(\theta^\star)}\insquare{-\log\inparen{z;\theta}}$.
                Standard properties of the exponential density (see, \eg{}, \citet{dasgupta2008asymptotic}) imply that:
                for all $i$ and $\theta$
                \begin{align*}
                    \nabla f_i(\theta) &= \nabla A(\theta) - t(x_i)
                    ~~\quad\quadand~~
                    \nabla^2 f_i(\theta) = \nabla^2 A(\theta)\,,\\
                    \nabla \overline{f}(\theta) &= \nabla A(\theta) - \nabla A(\theta^\star)
                    \quadand~~
                    \nabla^2 \overline{f}(\theta) = \nabla^2 A(\theta)\,.
                \end{align*}
                Where $A\colon \R^m \to \R_{>0}$ is the log-partition function of the exponential family and $t\colon \R^d\to \R^m$ is the sufficient statistic of the exponential family.
                
                \cref{asmp:exponentialFamily} implies that $f_i(\theta)$ and $\overline{f}(\theta)$ are $\lambda$-strongly convex and, hence, we can use \cite[Theorem 6.1]{charikar2017learning}.
                Before using \cite[Theorem 6.1]{charikar2017learning} though, we need to establish a spectral bound required by the theorem.
                Toward this, observe that for $i$ drawn uniformly from $[r]$ the expectation and variance of $\nabla f_i(\theta) - \nabla \overline{f}(\theta)$ is bounded for any $\theta\in \Theta$ as
                \begin{align*}
                    \Ex_{i\sim \unif{}([r])}\insquare{\nabla f_i(\theta) - \nabla \overline{f}(\theta)} 
                        &= \nabla A(\theta^\star) - \Ex_{z\sim \cE(\theta^\star)}[t(z)]
                        \tag{as $x_1,\dots,x_r$ are \iid{} from $\cE(\theta^\star)$}\\
                        &= 0 \tag{as $\Ex_{z\sim \cE(\theta)}[t(z)]=\nabla A(\theta)$ \cite{dasgupta2008asymptotic}}\,,\\
                    \cov_{i\sim \unif{}([r])}\insquare{\nabla f_i(\theta) - \nabla \overline{f}(\theta)} 
                    &= \cov_{i\sim \unif{}([r])}\insquare{
                        A(\theta) - t(x_i)
                        - \nabla A(\theta) + \nabla A(\theta^\star)
                    }\\
                    &= \cov_{i\sim \unif{}([r])}\insquare{
                        - t(x_i)
                    }\\
                    &=\cov_{z\sim \cE(\theta^\star)}\insquare{t(z)}\tag{as $\Ex_{z\sim \cE(\theta)}[t(z)]=\nabla A(\theta)$ \cite{dasgupta2008asymptotic}}\\
                    &\preceq \Lambda I\,.
                \end{align*}
                Now, we can establish a bound on $\min_{1\leq i\leq k}\norm{\theta_i-\theta^\star}_2$ using a straightforward adaptation of the proof of \cite[Corollary 9.1]{charikar2017learning} and then, we show that this implies a bound on $\cR_{w}$-divergence due to \cref{asmp:exponentialFamily}. 
                For completeness, we also include the proof of the bound on $\min_{1\leq i\leq k}\norm{\theta_i-\theta^\star}_2$ below. %
                
                \itparagraph{Bound on the $\cR_q$ divergence.}
                    Using \cite[Proposition B.1]{charikar2017learning} and above bounds on on the covariance and expectation of $\nabla f_i(\theta) - \nabla \overline{f}(\theta)$, implies that with probability at least $1-e^{\Omega(\beta^2\gamma n)}$, there exists a subset $G\subseteq [n]$ of $(1-(\nfrac{\beta}{2}))\gamma n$  ``good'' points that have a spectral bounds of $S=O(\nfrac{\Lambda}{\sqrt{\beta}})$ and (by \cite[Proposition 5.6]{charikar2017learning}) $S_{\sfrac{\beta}{2}}=O(\nfrac{\Lambda}{\beta})$.
                    Now we are ready to use \cite[Theorem 6.1]{charikar2017learning}, which gives parameters $\theta_1,\theta_2,\dots,\theta_k$ for $k\leq \inparen{(1-\beta)\gamma}^{-1}$ such that 
                    \[
                        \min_{1\leq i\leq k}\norm{\theta_i-\theta^\star}_2^2
                        \leq 
                        O\inparen{
                            \inparen{\frac{\Lambda}{\beta\lambda}}^2
                            \cdot
                            \frac{\log\inparen{\nfrac{2}{\gamma}}}{\gamma}
                        }
                        \yesnum\label{eq:application:listDecoding:paramDistance}
                        \,.
                    \] 

                \medskip 
                
                \noindent Next, we project $\theta_1,\dots,\theta_k$ to $\Theta(\eta)$.
                With some abuse of notation, let the projections be $\theta_1,\dots,\theta_k$.
                Since $\theta^\star\in \Theta(\eta)$, this only reduces the distance between $\theta_i$ and $\theta^\star$ for all $1\leq i\leq k$.
                Define 
                \[
                    R = O\inparen{
                            {\frac{\Lambda}{\lambda}}
                            \cdot
                            \sqrt{
                                \frac{\log\inparen{\nfrac{2}{\gamma}}}{\gamma}
                            }
                        }
                    \,.
                \]
                We claim that for any $i$, the $\cR_{w}$ divergence between $\cE(\theta_i)$ and $\cE(\theta^\star)$ is bounded as desired (in both directions).
                To show this, we prove the following fact. 
                    \begin{fact}\label{fact:app}
                        For any $\theta_1,\theta_2\in \Theta(\eta)$ that satisfy that $\norm{\theta_1-\theta_2}_2^2\leq R$, it holds that 
                        \[
                            \Ex_{x\sim \cE(\theta_2)}\insquare{
                            \inparen{
                                \frac{
                                    \cE(x; \theta_1)
                                }{
                                    \cE(x; \theta_2)
                                }
                            }^{w}
                            }
                            \leq \exp\inparen{
                                \Lambda ({w}R)^2 
                            }\,.
                        \]
                        Where the ${w}$ is as in the statement of \cref{thm:listDecoding:expFamily}.
                    \end{fact}
                    \begin{proof}[Proof of \cref{fact:app}]
                        Observe that 
                        \begin{align*}
                            \Ex_{x\sim \cE(\theta_2)}\insquare{
                                \inparen{
                                    \frac{
                                        \cE(x; \theta_1)
                                    }{
                                        \cE(x; \theta_2)
                                    }
                                }^{w}
                            }
                            &=
                            \int  e^{{w}(\theta_1-\theta_2)^\top t(x) + ({w}-1)A(\theta_2)-{w}A(\theta_1)} \cE(x; \theta_2)\d x\\
                            &=
                            \int  e^{
                                    ({w}-1)A(\theta_2)+ A\inparen{\theta_1+({w}-1)(\theta_1 - \theta_2)} - {w}A(\theta_1) 
                                }
                             \,.
                        \end{align*}
                        Since $\theta_1$ is in the $\eta$-relative interior of $\Theta$ and $\theta_1-\theta_2$ lies in the affine subspace of $\Theta$ with $\norm{\theta_1-\theta_2}\leq R$, for the chosen ${w}$, $\theta_1+({w}-1)(\theta_1 - \theta_2)\in \Theta$.
                        Since $A$ is a $\Lambda$-smooth function over $\Theta$, 
                        \[
                            ({w}-1)A(\theta_2)+ A\inparen{\theta_1+({w}-1)(\theta_1 - \theta_2)} - {w}A(\theta_1) \leq \Lambda \norm{\theta_1-\theta_2}_2^2
                            + \Lambda({w}-1)^2\norm{\theta_1-\theta_2}_2^2\,.
                        \]
                        Using that $\norm{\theta_1-\theta_2}_2\leq R$, it follows that
                        \[
                            \Ex_{x\sim \cE(\theta_2)}\insquare{
                            \inparen{
                                \frac{
                                    \cE(x; \theta_1)
                                }{
                                    \cE(x; \theta_2)
                                }
                                }^{w}
                            }
                            \leq 
                            \exp\inparen{
                                \Lambda ({w}R)^2
                            }
                            \,.
                        \]
                    \end{proof}
                Therefore, as for all $i$, $\theta_i\in \Theta(\eta)$ and $\theta^\star\in \Theta(\eta)$, we can use the above fact with the bounds on the parameter distance (\cref{eq:application:listDecoding:paramDistance}), 
                to conclude that there exists an $1\leq i\leq k$ such that 
                \[
                    \renyi{{w}}{\cE(\theta^\star)}{\cE(\theta_i)}
                    ~~\leq~~
                    \frac{\Lambda ({w}R)^2}{{w}-1}
                    ~~\leq~~ 
                    \frac{
                        \Lambda^4
                    }{
                        \gamma^3\lambda^3 
                    }\,.
                \]
                This implies the desired result by observing that $e^{-\Omega(\gamma n)}\leq \delta$.
            \end{proof}


\end{document}
